\paragraph{谱临界、输出慢模、端点有限部与熵缺损常数的统一闭合}
在碰撞位势的解析半径、输出位势的四阶累积结构、地基层端点能垒、两步 primitive 原子、极端偏置有限部以及 bin-fold 的信息位率阈值都已建立之后，还可把这些分散结论进一步压缩为下列七条彼此闭合的终端后果。它们分别刻画碰撞 RH 临界与解析分歧的严格分离、输出模扭曲慢模的四阶刚性、端点代价与地基熵的精确配平、输出大偏差相图的双坐标代数化、七条分支二回路在 primitive 层的单原子塌缩、极端偏置有限部对单一四次代数参数的依赖，以及熵缺损常数在三种信息尺度上的同一化。

\begin{theorem}[碰撞 RH 临界先于解析分歧出现]\label{thm:xi-time-part68-collision-rh-critical-before-analytic-branch}
设
\[
P_2(\theta):=\log \rho(e^\theta),
\]
并记 \(R_\theta\) 为其在 \(\theta=0\) 处 Taylor 级数的收敛半径。又设 \(u_R>1\) 为核版 RH/Ramanujan 条件
\[
\Lambda(u)\le \lambda_1(u)^{1/2}
\]
失效的唯一正临界点，\(\theta_R:=\log u_R\)。则
\[
0<\theta_R<R_\theta.
\]
更具体地，
\[
u_R\approx 3.59262181344,\qquad
\theta_R\approx 1.27888,\qquad
R_\theta\approx 2.76230289346087.
\]
因此，碰撞核的 RH/Ramanujan 失效并非由解析延拓的首个分歧点触发，而是在解析域内部由实谱不等式
\[
\Lambda(u)\le \lambda_1(u)^{1/2}
\]
的翻转所触发；该临界属于纯谱型临界，而非分支型临界。
\end{theorem}
\begin{proof}
定理 \ref{thm:xi-rh-critical-before-nearest-complex-branch} 已证明：\(P_2\) 的全部代数分歧点恰由判别式方程
\[
4u^3+25u^2+88u+8=0
\]
的三根经
\[
\theta=\log u_j+2\pi i k
\qquad (k\in\ZZ)
\]
产生，且最近分歧点的模长正是
\[
R_\theta\approx 2.76230289346087.
\]
同一定理又给出：\(u_R\) 是五次方程
\[
u^5+4u^4+3u^3-96u^2-42u-14=0
\]
的唯一正根，\(\theta_R=\log u_R\approx 1.27888\)，并且
\[
\textsf{RH}(u)\Longleftrightarrow 0<u\le u_R.
\]
比较数值便得
\[
0<\theta_R<R_\theta.
\]
于是对任意 \(\theta\in(\theta_R,R_\theta)\)，一方面 \(|\theta|<R_\theta\)，故 \(P_2(\theta)\) 仍位于 \(\theta=0\) 的解析收敛盘内；另一方面 \(u=e^\theta>u_R\)，故 RH/Ramanujan 条件已经失效。再由同一定理中的
\[
\Lambda(u_R)=\lambda_1(u_R)^{1/2},
\]
可知临界边界确由谱不等式的等号翻转刻画，而非由解析分支奇性触发。证毕。
\end{proof}

\begin{theorem}[输出模扭曲慢模的四阶累积刚性]\label{thm:xi-time-part68-output-twist-slowmode-fourth-cumulant-rigidity}
记
\[
\lambda(u):=\rho(M(u)),
\qquad
\lambda:=\lambda(1)=\varphi^2,
\qquad
P(\theta):=\log \lambda(e^\theta),
\qquad
\kappa_2:=P''(0),
\qquad
\kappa_4:=P^{(4)}(0).
\]
则
\[
\kappa_2,\kappa_4\in\QQ(\sqrt5),
\qquad
\kappa_2>0.
\]
对主单位根扭曲通道设
\[
t_m:=\frac{2\pi}{m},
\qquad
\rho_{m,1}:=\rho\!\bigl(M(e^{it_m})\bigr).
\]
则有严格渐近展开
\[
\log \frac{\rho_{m,1}}{\lambda}
=
-\frac{\kappa_2}{2}\Bigl(\frac{2\pi}{m}\Bigr)^2
+\frac{\kappa_4}{24}\Bigl(\frac{2\pi}{m}\Bigr)^4
+O(m^{-6}).
\]
并且在已审计模数 \(m=5,10,20\) 上，四阶预测的绝对误差分别约为
\[
1.08\times 10^{-6},\qquad
2.78\times 10^{-8},\qquad
6.21\times 10^{-10},
\]
而仅保留二阶项时的相应误差分别约为
\[
5.67\times 10^{-4},\qquad
3.66\times 10^{-5},\qquad
2.30\times 10^{-6}.
\]
因此，输出模扭曲的近平凡慢模在四阶以内已被前四阶累积量刚性锁定；四阶项并非可忽略修饰，而是恢复慢模几何所必需的主导校正。
\end{theorem}
\begin{proof}
命题 \ref{prop:real-input-40-output-cumulants-closed} 已给出
\[
P^{(k)}(0)\in\QQ(\sqrt5),
\]
特别是 \(\kappa_2,\kappa_4\in\QQ(\sqrt5)\) 且 \(\kappa_2>0\)。再由定理 \ref{thm:xi-output-potential-dirichlet-fixed-j-fourth-order} 在 \(j=1\) 时的结论，
\[
\log \frac{\rho_{m,1}}{\lambda}
=
-\frac{\kappa_2}{2}\Bigl(\frac{2\pi}{m}\Bigr)^2
+\frac{\kappa_4}{24}\Bigl(\frac{2\pi}{m}\Bigr)^4
+O(m^{-6}),
\]
这就给出所示四阶展开。

另一方面，附录注 \ref{rem:real-input-40-dirichlet-cumulant-error-order} 已逐项记录 \(m=5,10,20\) 上的四阶预测误差与纯二阶高斯近似误差，恰为题述数值。再由定理 \ref{thm:xi-output-potential-dirichlet-fourth-order-gain}，
\[
\frac{\widetilde\varepsilon_m^{(2)}}{\widetilde\varepsilon_m^{(4)}}=\Omega(m^2),
\]
故四阶修正不仅改善常数，而且在阶数量级上严格压缩误差。于是四阶项确为慢模重建所必需的主导校正。证毕。
\end{proof}

\begin{proposition}[端点代价缺口恰等于地基熵]\label{prop:xi-time-part68-endpoint-gap-ground-entropy-identity}
\leanverified{paper\_xi\_time\_part68\_endpoint\_gap\_ground\_entropy\_identity}
设 \(I(\alpha)\) 为输出位势
\[
P(\theta)=\log\lambda(e^\theta)
\]
的归一化大偏差速率函数，并记
\[
\alpha_{\max}:=\lim_{\theta\to+\infty}P'(\theta)=\frac12,
\qquad
h_{\mathrm{ground}}:=\log c.
\]
则有精确恒等式
\[
I(0)-I(\alpha_{\max})=h_{\mathrm{ground}}.
\]
更具体地，
\[
I(0)=\log(\varphi^2),
\qquad
I\!\left(\frac12\right)=\log\!\left(\frac{\varphi^2}{c}\right),
\qquad
h_{\mathrm{ground}}=\log c.
\]
因此，右端点的大偏差代价相对于左端点的全部减免，恰由地基层的组合熵指数精确补偿。
\end{proposition}
\begin{proof}
定理 \ref{thm:xi-output-potential-endpoint-gap-ground-state} 已直接给出
\[
I(0)-I\!\left(\frac12\right)=h_{\mathrm{ground}}.
\]
而推论 \ref{cor:real-input-40-alpha-max} 说明
\[
\alpha_{\max}=\frac12,
\]
推论 \ref{cor:real-input-40-ground-entropy} 则给出
\[
h_{\mathrm{ground}}=\log c.
\]
再结合推论 \ref{cor:real-input-40-rate-endpoints} 中的端点闭式
\[
I(0)=\log(\varphi^2),
\qquad
I(1/2)=\log(\varphi^2/c),
\]
即可得到题述恒等式与具体数值表达。证毕。
\end{proof}

\begin{theorem}[输出大偏差相图的代数一致化]\label{thm:xi-time-part68-output-ldp-phase-diagram-algebraization}
设 \(u=e^\theta\)，\(\lambda(u)\) 为方程
\[
G(\lambda,u)=0
\]
的 Perron 正根，并定义平衡态下输出 \(1\) 的典型密度
\[
\alpha(u):=\frac{\mathrm d}{\mathrm d\theta}\log\lambda(e^\theta).
\]
则存在非零多项式
\[
F(\alpha,u)\in\QQ[\alpha,u],
\qquad
H(\alpha,\lambda)\in\QQ[\alpha,\lambda]
\]
使得沿 Perron 分支恒有
\[
F(\alpha(u),u)=0,
\qquad
H(\alpha(u),\lambda(u))=0.
\]
并且归一化大偏差速率函数满足参数公式
\[
I(\alpha(u))
=
\alpha(u)\log u-\log\frac{\lambda(u)}{\lambda(1)}.
\]
因此，输出大偏差相图在 \((\alpha,u)\) 与 \((\alpha,\lambda)\) 两组坐标中都可代数化；真正的超越性只保留在 Legendre 对偶中的单一对数项。
\end{theorem}
\begin{proof}
命题 \ref{prop:real-input-40-pressure-freezing} 已给出输出位势的特征多项式
\[
G(\lambda,u)\in\QQ[\lambda,u],
\]
并说明 \(\lambda(u)\) 是其 Perron 正根。另一方面，附录 \ref{app:real-input-40-zeta-u} 中的输出位势大偏差代数参数化已把 \(\alpha(u)\) 写成
\[
\alpha(u)=\left.\frac{N(\lambda,u)}{D(\lambda,u)}\right|_{\lambda=\lambda(u)},
\qquad
N,D\in\QQ[\lambda,u].
\]
因此沿 Perron 分支有
\[
\alpha\,D(\lambda,u)-N(\lambda,u)=0,
\qquad
G(\lambda,u)=0.
\]
对 \(\lambda\) 取结果式并约去共同因子，便得到某个非零多项式
\[
F(\alpha,u)\in\QQ[\alpha,u]
\]
满足 \(F(\alpha(u),u)=0\)。同理，对 \(u\) 取结果式并约去共同因子，可得某个非零多项式
\[
H(\alpha,\lambda)\in\QQ[\alpha,\lambda]
\]
满足 \(H(\alpha(u),\lambda(u))=0\)。

最后，由同一附录中的参数公式，或者直接由严格凸压力的驻点 Legendre 公式，
\[
I(\alpha(u))
=
\alpha(u)\log u-\log\frac{\lambda(u)}{\lambda(1)}.
\]
于是 \((\alpha,u)\) 与 \((\alpha,\lambda)\) 两种坐标下的相图都由代数关系控制，而超越部分只出现在该单一对数项中。证毕。
\end{proof}

\begin{theorem}[七条分支二回路在 primitive 层塌缩为唯一原子]\label{thm:xi-time-part68-seven-branch-two-cycles-collapse-single-atom}
设输出位势的 primitive 生成函数为
\[
P(z,u)=\sum_{n\ge 1}p_n(u)z^n.
\]
则在真实输入 \(40\) 状态核上有精确分解
\[
P(z,u)=uz^2+P_{\mathrm{core}}(z,u).
\]
另一方面，branch 层审计表明，满足
\[
(|\tau|,E)=(2,1)
\]
的两步回路恰有 \(7\) 条。因而，这 \(7\) 条几何上可区分的 branch 二回路在 primitive--Witt 化之后并不留下 \(7\) 个独立原子，而是完全塌缩为唯一长度 \(2\) 原子
\[
uz^2.
\]
\end{theorem}
\begin{proof}
定理 \ref{thm:xi-time-part62d-atomic-two-step-witt-rankone-collapse} 已直接给出
\[
P(z,u)=uz^2+P_{\mathrm{core}}(z,u),
\]
并同时指出 branch 层中满足
\[
(|\tau|,E)=(2,1)
\]
的两步回路总数恰为 \(7\)。逐项比较系数可知，full/core 的 primitive 差分只支撑在长度 \(2\) 的单一坐标上；等价地，
\[
\dim_{\QQ(u)}
\operatorname{span}\{\,p_n(u)-p_n^{\mathrm{core}}(u):n\ge 1\,\}=1.
\]
因此，这七条 branch 两步回路在 primitive--Witt 层只能留下唯一原子 \(uz^2\)。证毕。
\end{proof}

\begin{theorem}[极端偏置有限部由单一四次代数参数完全控制]\label{thm:xi-time-part68-extreme-bias-finitepart-single-quartic-control}
记
\[
b:=b_\infty:=\rho(B_\infty)^{-1}\in(0,1),
\]
等价地，\(b\) 是四次方程
\[
1-6b+9b^2-b^3-b^4=0
\]
在区间 \((0,1)\) 内的最小正根。则极端偏置 \(u\to+\infty\) 下，主极点常数 \(C_\infty\) 与 Abel 有限部常数 \(\mathfrak M_\infty\) 满足
\[
C_\infty=
\frac{1}{2b(1-b)(6-18b+3b^2+4b^3)},
\]
\[
\log\mathfrak M_\infty
=
\log C_\infty
-\sum_{m\ge 2}\frac{\mu(m)}{m}
\log\Bigl[(1-b^m)\bigl(1-6b^m+9b^{2m}-b^{3m}-b^{4m}\bigr)\Bigr].
\]
数值上
\[
b\approx 0.27807480214113,\qquad
C_\infty\approx 1.89745149363,\qquad
\log\mathfrak M_\infty\approx 0.282290255262.
\]
因此，极端偏置下的全部有限部数据都由单一四次代数参数 \(b\) 完整统辖。
\end{theorem}
\begin{proof}
命题 \ref{prop:real-input-40-logM-infty} 已直接证明：令
\[
b=\lim_{u\to\infty}u\,z_\star(u)^2=\frac{1}{c^2},
\]
则 \(b\) 满足题述四次方程，并且 \(C_\infty\) 与 \(\log\mathfrak M_\infty\) 的闭式恰为上式。另一方面，推论 \ref{cor:xi-zero-temp-leading-pole-ground-dominance} 说明
\[
b=b_\infty=\rho(B_\infty)^{-1}
\]
正是零温缩放极限中控制主导正实零点的位置参数。故无论从主极点常数还是从 Abel 有限部常数看，极端偏置下的主导有限部数据都确实由同一个四次代数参数 \(b\) 完整决定。证毕。
\end{proof}

\begin{theorem}[熵缺损常数的三重同一化]\label{thm:xi-time-part68-entropy-deficit-triple-identification}
记
\[
\mu_m(w):=\frac{d_m^{\mathrm{bin}}(w)}{2^m},
\qquad
W_m\sim\mu_m,
\qquad
\gamma_m:=\frac{1}{m}\log d_m^{\mathrm{bin}}(W_m).
\]
则对任意 \(q>0\)（其中 \(q=1\) 按 Shannon 熵解释）都有
\[
\lim_{m\to\infty}\frac{1}{m}\bigl(m\log 2-H_q(\mu_m)\bigr)
=
\log\frac{2}{\varphi},
\qquad
\gamma_m\longrightarrow \log\frac{2}{\varphi}.
\]
若一列二进制预算 \(B_m\in\NN\) 满足
\[
\mathrm{Succ}_m^{\mathrm{bin}}(B_m)\ge 1-\varepsilon
\qquad
(0<\varepsilon<1\ \text{固定}),
\]
则必有
\[
\liminf_{m\to\infty}\frac{B_m}{m}
\ge
\log_2\frac{2}{\varphi}.
\]
因此，同一个常数 \(\log(2/\varphi)\) 同时控制全部正阶 Rényi 熵的广延缺损、典型微态的纤维多重度指数以及保持固定成功率所需的最小外置比特率。
\end{theorem}
\begin{proof}
由定理 \ref{thm:fold-bin-renyi-rate-collapse}，
\[
\frac{1}{m}H_q(\mu_m)\longrightarrow \log\varphi
\qquad (q>0),
\]
并且
\[
\gamma_m\longrightarrow \log\frac{2}{\varphi}.
\]
于是
\[
\frac{1}{m}\bigl(m\log 2-H_q(\mu_m)\bigr)
=
\log 2-\frac{1}{m}H_q(\mu_m)
\longrightarrow
\log 2-\log\varphi
=
\log\frac{2}{\varphi},
\]
这就得到前两条极限。

另一方面，若用 \(B_m\) 个二进制外置信息位实现恢复，则其状态数为 \(M_m=2^{B_m}\)。推论 \ref{cor:xi-reconstruction-cardinality-strong-converse} 对 \(\delta=1-\varepsilon\) 给出
\[
B_m
\ge
m-\log_2F_{m+2}+\log_2(1-\varepsilon).
\]
两边同除以 \(m\)，再用 Binet 渐近
\[
\log_2F_{m+2}=m\log_2\varphi+O(1),
\]
便得
\[
\liminf_{m\to\infty}\frac{B_m}{m}\ge \log_2\frac{2}{\varphi}.
\]
故题述三种量确由同一个常数统一控制。证毕。
\end{proof}

\noindent
上述七条结果表明，碰撞核的 RH 破裂、输出单位根慢模、地基层端点能垒、两步 primitive 原子、零温有限部常数与 bin-fold 的熵缺损阈值并不是彼此分离的局部现象。它们共同揭示：一方面，碰撞侧的临界首先发生在解析盘内部，因而属于纯谱型翻转；另一方面，输出侧的慢模、端点代价与大偏差相图都受同一条代数压力曲线及其四阶累积结构支配；最终，这种谱--代数刚性又在 primitive 原子剥离、零温有限部与信息预算下界中外显为同一组可审计的不变量。因而，从碰撞 RH 窗口到输出大偏差，再到部分可逆预算与极端偏置有限部，全文中的主导常数与临界门槛实际上共享同一条谱--组合--信息几何闭合链。

\paragraph{低预算线性相、二点实验塌缩、局部化周期谱与形式线性化的闭合后果}
在审计偶窗的最小退化 Fibonacci 律、末位两层塌缩、局部化 solenoid 的自覆盖分类、素数寄存器的幂等估值算术化以及 fold-gauge 缺陷的纯奇次重整化已经建立之后，还可把这些链条继续压缩为下列七条结果。它们分别刻画审计偶窗中的低预算严格线性段及其首个偏离阈值、bin-fold 输出实验的 Le Cam 二点塌缩、一般局部化有理乘法自同态的固定点与 Artin--Mazur \(\zeta\) 函数、相应周期点几何的指数率、素数寄存器幂等元的精确闭式计数，以及 fold-gauge 缺陷的形式 Koenigs 线性化。

\begin{theorem}[bin-fold 审计偶数窗上的低预算严格线性律]\label{thm:xi-time-part69-audited-even-pure-pigeonhole-linear-phase}
对审计偶数窗
\[
m\in\{6,8,10,12\},
\]
定义
\[
d_{\min}(m):=\min_{w\in X_m}d_m^{\mathrm{bin}}(w),
\qquad
C_m^{\mathrm{bin}}(B):=\sum_{w\in X_m}\min\bigl(d_m^{\mathrm{bin}}(w),2^B\bigr),
\]
\[
\mathrm{Succ}_m^{\mathrm{bin}}(B):=\frac{C_m^{\mathrm{bin}}(B)}{2^m}.
\]
若整数预算 \(B\ge 0\) 满足
\[
2^B\le d_{\min}(m)=F_{m/2},
\]
则有精确公式
\[
C_m^{\mathrm{bin}}(B)=2^B\,F_{m+2},
\qquad
\mathrm{Succ}_m^{\mathrm{bin}}(B)=2^{B-m}F_{m+2}.
\]
因此，在最小退化阈值之前，bin-fold 的每增加一比特辅助信息都会使可逆微态总数严格翻倍，且不出现任何次级修正。
\end{theorem}
\begin{proof}
定理 \ref{thm:fold-bin-min-degeneracy-fib-audited} 已证明：当
\[
m\in\{6,8,10,12\}
\]
时，
\[
d_{\min}(m)=F_{m/2}.
\]
若 \(2^B\le F_{m/2}\)，则对所有 \(w\in X_m\) 都有
\[
2^B\le d_{\min}(m)\le d_m^{\mathrm{bin}}(w),
\]
从而
\[
\min\bigl(d_m^{\mathrm{bin}}(w),2^B\bigr)=2^B.
\]
因此
\[
C_m^{\mathrm{bin}}(B)
=
\sum_{w\in X_m}2^B
=
2^B|X_m|.
\]
另一方面，由稳定词计数律
\[
|X_m|=F_{m+2},
\]
即得
\[
C_m^{\mathrm{bin}}(B)=2^BF_{m+2}.
\]
再除以 \(2^m\)，得到
\[
\mathrm{Succ}_m^{\mathrm{bin}}(B)=2^{B-m}F_{m+2}.
\]
证毕。
\end{proof}

\begin{corollary}[bin-fold 容量曲线的首个折点恰位于 \texorpdfstring{$F_{m/2}$}{F(m/2)}]\label{cor:xi-time-part69-binfold-first-kink}
在定理 \ref{thm:xi-time-part69-audited-even-pure-pigeonhole-linear-phase} 的设定下，容量曲线
\[
B\longmapsto C_m^{\mathrm{bin}}(B)
\]
的首个偏离线性区间发生在最小整数
\[
B_{\mathrm{kink}}(m):=\min\{B\in\NN:\ 2^B>F_{m/2}\}
=
\bigl\lfloor \log_2 F_{m/2}\bigr\rfloor+1.
\]
更准确地，
\[
C_m^{\mathrm{bin}}(B)=2^BF_{m+2}
\quad\Longleftrightarrow\quad
2^B\le F_{m/2},
\]
而一旦 \(2^B>F_{m/2}\)，便有严格不等式
\[
C_m^{\mathrm{bin}}(B)<2^BF_{m+2}.
\]
\end{corollary}
\begin{proof}
“若”方向即定理 \ref{thm:xi-time-part69-audited-even-pure-pigeonhole-linear-phase}。反之，由 \(X_m\) 的有限性与
\[
d_{\min}(m)=F_{m/2}
\]
可取 \(w_0\in X_m\) 使
\[
d_m^{\mathrm{bin}}(w_0)=F_{m/2}.
\]
若 \(2^B>F_{m/2}\)，则
\[
\min\bigl(d_m^{\mathrm{bin}}(w_0),2^B\bigr)=F_{m/2}<2^B,
\]
而其余各项至多为 \(2^B\)。于是
\[
C_m^{\mathrm{bin}}(B)<2^B|X_m|=2^BF_{m+2}.
\]
故首个折点恰由阈值 \(F_{m/2}\) 决定。证毕。
\end{proof}

\begin{theorem}[bin-fold 统计实验的 Le Cam 二点塌缩]\label{thm:xi-time-part69-binfold-lecam-two-point-collapse}
\leanverified{paper\_xi\_time\_part69\_binfold\_lecam\_two\_point\_collapse}
设
\[
\mu_m(w):=\frac{d_m^{\mathrm{bin}}(w)}{2^m},
\qquad
u_m(w):=\frac{1}{|X_m|},
\qquad
\ell_m(w):=w_m\in\{0,1\},
\]
并定义两块
\[
A_0:=\{w\in X_m:\ w_m=0\},
\qquad
A_1:=\{w\in X_m:\ w_m=1\}.
\]
令
\[
T_m:=\ell_m:X_m\to\{0,1\},
\qquad
K_m(i,\cdot):=\mathrm{Unif}(A_i)\qquad(i=0,1).
\]
再定义极限二点实验
\[
Q_\infty:=\left(\frac1\varphi,\frac1{\varphi^2}\right),
\qquad
P_\infty:=\left(\frac{\varphi}{\sqrt5},\frac{1}{\varphi\sqrt5}\right),
\]
其中坐标顺序对应 \((0,1)\)。则原始两点统计实验
\[
\mathcal E_m:=\{u_m,\mu_m\}\quad\text{定义在 }X_m
\]
与二元实验
\[
\mathcal E_\infty:=\{Q_\infty,P_\infty\}\quad\text{定义在 }\{0,1\}
\]
在 Le Cam 意义下渐近等价。更具体地，
\[
D_{\mathrm{TV}}\!\bigl(T_{m\#}u_m,Q_\infty\bigr)=O(\varphi^{-2m}),
\qquad
D_{\mathrm{TV}}\!\bigl(T_{m\#}\mu_m,P_\infty\bigr)=O\!\Bigl(\Bigl(\frac{\varphi}{2}\Bigr)^m\Bigr),
\]
并且
\[
D_{\mathrm{TV}}\!\bigl(u_m,K_{m\#}Q_\infty\bigr)=O(\varphi^{-2m}),
\qquad
D_{\mathrm{TV}}\!\bigl(\mu_m,K_{m\#}P_\infty\bigr)=O\!\Bigl(\Bigl(\frac{\varphi}{2}\Bigr)^m\Bigr).
\]
因此
\[
\Delta(\mathcal E_m,\mathcal E_\infty)\xrightarrow[m\to\infty]{}0.
\]
\end{theorem}
\begin{proof}
先看均匀基线 \(u_m\)。由
\[
|A_0|=F_{m+1},
\qquad
|A_1|=F_m,
\qquad
|X_m|=F_{m+2},
\]
可得
\[
T_{m\#}u_m=
\left(\frac{F_{m+1}}{F_{m+2}},\frac{F_m}{F_{m+2}}\right)
\]
且由 Binet 公式比较相邻 Fibonacci 比值可知，其两坐标分别与 \(Q_\infty\) 的对应坐标相差 \(O(\varphi^{-2m})\)。特别地，
\[
D_{\mathrm{TV}}\!\bigl(T_{m\#}u_m,Q_\infty\bigr)=O(\varphi^{-2m}).
\]
另一方面，\(u_m\) 在每个块 \(A_i\) 上本来就是严格均匀分布，因此
\[
u_m=K_{m\#}(T_{m\#}u_m).
\]
由 Markov 核的收缩性，
\[
D_{\mathrm{TV}}\!\bigl(u_m,K_{m\#}Q_\infty\bigr)
\le
D_{\mathrm{TV}}\!\bigl(T_{m\#}u_m,Q_\infty\bigr)
=
O(\varphi^{-2m}).
\]

再看 \(\mu_m\)。定理 \ref{thm:xi-fold-lastbit-asymptotic-sufficiency} 已给出
\[
D_{\mathrm{TV}}\!\bigl(\mu_m,K_{m\#}(T_{m\#}\mu_m)\bigr)
=
O\!\Bigl(\Bigl(\frac{\varphi}{2}\Bigr)^m\Bigr).
\]
另一方面，定理 \ref{thm:xi-fold-lastbit-two-layer-collapse} 证明 \(\mu_m\) 与末位两层模型 \(\nu_m(w)=\varphi^{-m-w_m}\) 在全变差意义下相差
\[
O\!\Bigl(\Bigl(\frac{\varphi}{2}\Bigr)^m\Bigr).
\]
而由
\[
\nu_m(A_0)=\varphi^{-m}F_{m+1}=\frac{\varphi}{\sqrt5}+O(\varphi^{-2m}),
\qquad
\nu_m(A_1)=\varphi^{-m-1}F_m=\frac{1}{\varphi\sqrt5}+O(\varphi^{-2m}),
\]
再与
\[
D_{\mathrm{TV}}(\mu_m,\nu_m)=O\!\Bigl(\Bigl(\frac{\varphi}{2}\Bigr)^m\Bigr)
\]
合并，可知
\[
D_{\mathrm{TV}}\!\bigl(T_{m\#}\mu_m,P_\infty\bigr)
=
O\!\Bigl(\Bigl(\frac{\varphi}{2}\Bigr)^m\Bigr).
\]
于是再用 Markov 核收缩性，
\[
\begin{aligned}
D_{\mathrm{TV}}\!\bigl(\mu_m,K_{m\#}P_\infty\bigr)
&\le
D_{\mathrm{TV}}\!\bigl(\mu_m,K_{m\#}(T_{m\#}\mu_m)\bigr)
+
D_{\mathrm{TV}}\!\bigl(K_{m\#}(T_{m\#}\mu_m),K_{m\#}P_\infty\bigr)\\
&\le
O\!\Bigl(\Bigl(\frac{\varphi}{2}\Bigr)^m\Bigr)
+
D_{\mathrm{TV}}\!\bigl(T_{m\#}\mu_m,P_\infty\bigr)\\
&=
O\!\Bigl(\Bigl(\frac{\varphi}{2}\Bigr)^m\Bigr).
\end{aligned}
\]

最后，对参数集 \(\{0,1\}\) 上的两点实验，前向缺陷由核 \(T_m\) 控制，后向缺陷由核 \(K_m\) 控制，因此 Le Cam 距离满足
\[
\Delta(\mathcal E_m,\mathcal E_\infty)
\le
\max\Bigl\{
D_{\mathrm{TV}}(T_{m\#}u_m,Q_\infty),
D_{\mathrm{TV}}(T_{m\#}\mu_m,P_\infty),
D_{\mathrm{TV}}(u_m,K_{m\#}Q_\infty),
D_{\mathrm{TV}}(\mu_m,K_{m\#}P_\infty)
\Bigr\},
\]
右端趋于 \(0\)，故结论成立。证毕。
\end{proof}

\begin{theorem}[局部化 solenoid 上有理乘法自同态的固定点与 Artin--Mazur \texorpdfstring{$\zeta$}{zeta} 函数]\label{thm:xi-time-part69-localized-rational-endomorphism-fixedpoints-zeta}
\leanverified{paper\_xi\_time\_part69\_localized\_rational\_endomorphism\_fixedpoints\_zeta}
设 \(S\subset\mathbb P\) 为有限素数集，并记
\[
G_S:=\ZZ[S^{-1}],
\qquad
\Sigma_S:=\widehat{G_S}.
\]
取有理数
\[
r=\frac{a}{b}\in G_S,
\qquad
\gcd(a,b)=1,
\]
其中 \(b\) 的全部素因子都属于 \(S\)，并设 \(f_r:\Sigma_S\to\Sigma_S\) 为乘法映射
\[
x\longmapsto rx\qquad (x\in G_S)
\]
的 Pontryagin 对偶连续自同态。若 \(|a|\neq |b|\)，则对每个 \(n\ge 1\)，都有
\[
\#\mathrm{Fix}(f_r^n)=|a^n-b^n|_{S^\perp},
\]
其中对任意非零整数 \(N\)，记
\[
|N|_{S^\perp}:=(|N|)_{S^\perp}
\]
为其 \(S\)-自由部分。并且 Artin--Mazur 动力学 \(\zeta\) 函数满足
\[
\zeta_{f_r}(z)
:=
\exp\!\left(
\sum_{n\ge 1}\frac{\#\mathrm{Fix}(f_r^n)}{n}z^n
\right)
=
\exp\!\left(
\sum_{n\ge 1}\frac{|a^n-b^n|_{S^\perp}}{n}z^n
\right).
\]
换言之，局部化 solenoid 上的周期点几何只感知 \(a^n-b^n\) 的 \(S\)-外部。
\end{theorem}
\begin{proof}
由定理 \ref{thm:xi-localized-solenoid-self-cover-degree-classification}，任意连续自同态都唯一对应于某个 \(r\in G_S\)，其对偶离散同态为
\[
x\longmapsto rx
\qquad (x\in G_S).
\]
因此对每个 \(n\ge 1\)，
\[
f_r^n
\]
在对偶离散侧对应于乘法
\[
x\longmapsto r^n x.
\]
按定理 \ref{thm:xi-localized-solenoid-fixedpoints-artin-mazur-series} 的同样论证，
\[
\mathrm{Fix}(f_r^n)\cong \widehat{G_S/(r^n-1)G_S}.
\]

现把
\[
r^n-1=\frac{a^n-b^n}{b^n}
\]
代入。由于 \(b\) 的全部素因子都属于 \(S\)，故 \(b^n\in G_S^\times\) 是单位，从而
\[
(r^n-1)G_S=(a^n-b^n)G_S.
\]
又因 \(|a|\neq |b|\)，故 \(a^n-b^n\neq 0\)。于是由定理 \ref{thm:xi-localized-quotient-exact-sfree-externalization}，
\[
G_S/(a^n-b^n)G_S
\cong
\ZZ/|a^n-b^n|_{S^\perp}\ZZ.
\]
取 Pontryagin 对偶便得到
\[
\#\mathrm{Fix}(f_r^n)=|a^n-b^n|_{S^\perp}.
\]
再将该固定点计数代入 Artin--Mazur \(\zeta\) 函数定义，即得所示级数公式。证毕。
\end{proof}

\begin{theorem}[局部化有理自覆盖的周期点指数率只由分子分母的最大模控制]\label{thm:xi-time-part69-localized-rational-endomorphism-growth-rate}
沿用定理 \ref{thm:xi-time-part69-localized-rational-endomorphism-fixedpoints-zeta} 的记号，并进一步假设
\[
|a|\neq |b|.
\]
则有极限
\[
\lim_{n\to\infty}\frac{1}{n}\log \#\mathrm{Fix}(f_r^n)
=
\log\max\{|a|,|b|\}.
\]
换言之，局部化 solenoid 上有理自覆盖的周期点指数率不受 \(S\)-内素因子的影响，而仅由分子与分母模长中的较大者决定。
\end{theorem}
\begin{proof}
由定理 \ref{thm:xi-time-part69-localized-rational-endomorphism-fixedpoints-zeta} 证明中的同一对偶化论证，对每个 \(n\ge 1\) 都有
\[
\mathrm{Fix}(f_r^n)\cong \widehat{G_S/(r^n-1)G_S}.
\]
又因 \(|a|\neq |b|\)，故 \(a^n-b^n\neq 0\)。并且
\[
r^n-1=\frac{a^n-b^n}{b^n},
\]
其中 \(b^n\in G_S^\times\) 为单位，于是
\[
(r^n-1)G_S=(a^n-b^n)G_S.
\]
再由定理 \ref{thm:xi-localized-quotient-exact-sfree-externalization}，
\[
\#\mathrm{Fix}(f_r^n)=|a^n-b^n|_{S^\perp}.
\]
因此只需比较 \(a^n-b^n\) 的 \(S\)-内与 \(S\)-外部分的增长。

先固定 \(p\in S\)。若 \(p\mid ab\)，则由 \(\gcd(a,b)=1\) 可知 \(p\) 不能同时整除 \(a\) 与 \(b\)，从而对任意 \(n\ge 1\) 都有
\[
p\nmid (a^n-b^n).
\]
若 \(p\nmid ab\)，则 \(a/b\) 在 \((\ZZ/p\ZZ)^\times\) 中有定义。记其模 \(p\) 的乘法阶为 \(t_p\)。只有在 \(t_p\mid n\) 时，\(p\) 才可能整除 \(a^n-b^n\)；并且由标准 lifting-the-exponent 估计，存在仅依赖于 \(p,a,b\) 的常数 \(C_p\)，使得
\[
v_p(a^n-b^n)\le C_p+v_p(n)
\qquad(n\ge 1).
\]
由于 \(S\) 有限，遂得
\[
\sum_{p\in S}v_p(a^n-b^n)\log p=O(\log n),
\]
因而
\[
\frac{1}{n}\sum_{p\in S}v_p(a^n-b^n)\log p\longrightarrow 0.
\]

另一方面，若 \(|a|>|b|\)，则
\[
|a^n-b^n|
=
|a|^n\Bigl|1-(b/a)^n\Bigr|
=
|a|^n(1+o(1)),
\]
故
\[
\frac{1}{n}\log |a^n-b^n|\longrightarrow \log|a|.
\]
若 \(|b|>|a|\)，同理得到极限为 \(\log|b|\)。

最后，由
\[
\log |a^n-b^n|_{S^\perp}
=
\log |a^n-b^n|
-\sum_{p\in S}v_p(a^n-b^n)\log p
\]
即知
\[
\frac{1}{n}\log |a^n-b^n|_{S^\perp}
\longrightarrow
\log\max\{|a|,|b|\}.
\]
再代回
\(\#\mathrm{Fix}(f_r^n)=|a^n-b^n|_{S^\perp}\)，即得所示极限。证毕。
\end{proof}

\begin{theorem}[素数寄存器半群幂等元的精确计数律]\label{thm:xi-time-part69-prime-register-idempotent-count-law}
沿用定理 \ref{thm:killo-finite-semigroup-prime-valuation-arithmetization} 的记号，设
\[
S_n:=\left\{\prod_{i=0}^{n-1}q_i^{a_i}: a_i\in\{0,\dots,n-1\}\right\},
\]
并记
\[
E_n:=\{N\in S_n:\ N\star N=N\},
\qquad
E_{n,k}:=\{N\in E_n:\ |\operatorname{Im}(f_N)|=k\}
\qquad (1\le k\le n).
\]
则有精确计数公式
\[
|E_{n,k}|=\binom{n}{k}k^{n-k}
\qquad (1\le k\le n),
\]
从而
\[
|E_n|=\sum_{k=1}^{n}\binom{n}{k}k^{n-k}.
\]
换言之，素数寄存器中的幂等投影按像集大小完全分层，而其总数与 \([n]\) 上幂等自映射的经典计数严格一致。
\end{theorem}
\begin{proof}
由定理 \ref{thm:killo-finite-semigroup-prime-valuation-arithmetization} 与推论 \ref{cor:killo-projection-idempotent-prime-register}，\(N\in S_n\) 为幂等元当且仅当对应映射
\[
f_N:[n]\to[n]
\]
为幂等变换。

固定 \(k\in\{1,\dots,n\}\)，并设 \(|\operatorname{Im}(f_N)|=k\)。记
\[
A:=\operatorname{Im}(f_N)\subseteq [n],
\qquad |A|=k.
\]
幂等条件
\[
f_N\circ f_N=f_N
\]
等价于：每个 \(a\in A\) 都必须是定点，即
\[
f_N(a)=a
\qquad (a\in A),
\]
而对每个 \(x\in[n]\setminus A\)，其像只需任意落在 \(A\) 中。因而，一旦像集 \(A\) 固定，幂等变换数恰为
\[
k^{n-k}.
\]
再对全部 \(\binom{n}{k}\) 个 \(k\)-元像集求和，即得
\[
|E_{n,k}|=\binom{n}{k}k^{n-k}.
\]
最后对 \(k=1,\dots,n\) 求和，得到
\[
|E_n|=\sum_{k=1}^{n}\binom{n}{k}k^{n-k}.
\]
证毕。
\end{proof}

\begin{theorem}[fold-gauge 缺陷的形式 Koenigs 线性化]\label{thm:xi-time-part69-fold-gauge-formal-koenigs-linearization}
\leanverified{paper\_xi\_time\_part69\_fold\_gauge\_formal\_koenigs\_linearization}
记
\[
E_m:=\log C_m-\log(2\pi),
\qquad
x:=\frac{\varphi}{2}\in(0,1).
\]
由定理 \ref{thm:xi-time-part60ab-gauge-constant-recursive-bernoulli-stripping}，存在奇形式级数
\[
F(y):=\sum_{r\ge 1}a_r y^{2r-1}\in y\RR[[y^2]],
\qquad
a_1=\frac{1}{3\varphi}\neq 0,
\]
使
\[
E_m=F(x^m).
\]
则存在唯一奇形式幂级数
\[
\mathcal K(z)=3\varphi\,z+\kappa_3 z^3+\kappa_5 z^5+\cdots \in z\RR[[z^2]]
\]
使得对任意 \(N\ge 1\)，若记 \(\mathcal K_N\) 为其截断到 \(z^{2N+1}\) 的多项式，则
\[
\mathcal K_N(E_m)=x^m+O\!\bigl(x^{(2N+1)m}\bigr)
\qquad (m\to\infty).
\]
等价地，若 \(\mathcal R(z)\in z\RR[[z^2]]\) 为定理 \ref{thm:xi-time-part64a-fold-gauge-purely-odd-renormalization} 中的唯一形式重整化映射，则
\[
\mathcal K(\mathcal R(z))=x\,\mathcal K(z)
\]
在形式幂级数意义下成立。其首个非平凡系数为
\[
\kappa_3=\frac{135+63\sqrt5}{40}.
\]
\end{theorem}
\begin{proof}
由于 \(a_1\neq 0\)，级数 \(F\) 在形式幂级数环中存在唯一复合逆
\[
F^{-1}(z)\in z\RR[[z]].
\]
又因为 \(F\) 为奇级数，其逆也必为奇级数，于是
\[
\mathcal K(z):=F^{-1}(z)\in z\RR[[z^2]].
\]
由构造立得
\[
\mathcal K(E_m)=\mathcal K(F(x^m))=x^m
\]
作为形式恒等式成立。若 \(\mathcal K_N\) 为截断到 \(z^{2N+1}\) 的多项式，则
\[
\mathcal K(z)-\mathcal K_N(z)=O(z^{2N+3}),
\]
而 \(E_m\sim a_1x^m\)，故
\[
\mathcal K_N(E_m)=x^m+O(x^{(2N+3)m}),
\]
特别地也有较弱的
\[
\mathcal K_N(E_m)=x^m+O(x^{(2N+1)m}).
\]

另一方面，定理 \ref{thm:xi-time-part64a-fold-gauge-purely-odd-renormalization} 已给出唯一形式重整化映射
\[
\mathcal R(z)=F\!\bigl(xF^{-1}(z)\bigr).
\]
于是
\[
\mathcal K(\mathcal R(z))
=
F^{-1}\!\bigl(F(xF^{-1}(z))\bigr)
=
xF^{-1}(z)
=
x\,\mathcal K(z),
\]
这正是 Koenigs 线性化方程。

最后计算首项系数。由 \(r=1,2\) 的系数公式，
\[
a_1=\frac{1}{3\varphi},
\qquad
a_2=\frac{B_4}{2\cdot 3}\,(\varphi^{-1}+\varphi)
=
-\frac{\sqrt5}{180},
\]
其中用到 \(B_4=-1/30\) 与 \(\varphi+\varphi^{-1}=\sqrt5\)。设
\[
\mathcal K(z)=k_1 z+k_3 z^3+O(z^5).
\]
把它代入
\[
\mathcal K(F(y))=y
\]
并比较 \(y\) 与 \(y^3\) 的系数，得到
\[
k_1a_1=1,
\qquad
k_1a_2+k_3a_1^3=0.
\]
第一式给出
\[
k_1=\frac{1}{a_1}=3\varphi.
\]
第二式给出
\[
k_3=-\frac{a_2}{a_1^4}
=
\frac{135+63\sqrt5}{40}.
\]
故 \(\kappa_3=k_3\)。证毕。
\end{proof}

\noindent
上述七条结果表明，审计偶窗的低预算容量、其首个非线性折点、bin-fold 的完整统计实验、局部化 solenoid 的周期点几何、素数寄存器的有限幂等结构以及 fold-gauge 缺陷的形式重整化之间存在同一类低维闭合机制：在有限分辨率层面，它表现为严格线性段、精确折点与精确的幂等计数；在统计层面，它表现为末位二点实验对全模型的 Le Cam 塌缩；在局部化动力学层面，它既表现为 \(S\)-外部因子对周期点计数的完全控制，也表现为指数率仅由 \(\max\{|a|,|b|\}\) 锁定；而在形式动力学层面，它最终收束为 Koenigs 坐标下的严格乘法线性化。于是，有限预算盲区、容量折点、二点统计极限、局部化周期谱与重整化正规形不再是彼此分离的技术现象，而被统一压缩为同一条可审计的降维主线。

\paragraph{单位自同构、周期指数与 prime-register 盲性}
在一般局部化有理自同态的固定点与指数率已经获得闭式之后，单位参数区间还可进一步收束为一条更强的动力学链：连续自同构群完全由单位格给出，周期点公式退化为 \(S\)-自由差分的精确读数，而遍历、mixing 与同构分类之间则出现了严格分离。

\begin{theorem}[局部化 solenoid 的连续自同构群完全分类]\label{thm:xi-time-part69-localized-unit-automorphism-group-classification}
\leanverified{paper\_xi\_time\_part69\_localized\_unit\_automorphism\_group\_classification}
设 \(S\subset\mathbb P\) 为有限素数集，并记
\[
G_S:=\ZZ[S^{-1}],
\qquad
\Sigma_S:=\widehat{G_S}.
\]
则连续群自同构群满足显式同构
\[
\boxed{
\operatorname{Aut}_{\mathrm{cts}}(\Sigma_S)
\cong
G_S^\times
=
\left\{\pm\prod_{p\in S}p^{n_p}:n_p\in\ZZ\right\}
\cong
C_2\times \ZZ^{|S|}.
}
\]
\end{theorem}
\begin{proof}
定理 \ref{thm:xi-localized-solenoid-self-cover-degree-classification} 已给出
\[
\operatorname{Aut}_{\mathrm{cts}}(\Sigma_S)\cong G_S^\times
=
\left\{\pm\prod_{p\in S}p^{n_p}:n_p\in\ZZ\right\}.
\]
由于 \(S\) 有限，映射
\[
\varepsilon\prod_{p\in S}p^{n_p}
\longmapsto
\bigl(\varepsilon,(n_p)_{p\in S}\bigr),
\qquad
\varepsilon\in\{\pm1\},
\]
给出该单位群与 \(C_2\times \ZZ^{|S|}\) 的群同构，故结论成立。证毕。
\end{proof}

\begin{theorem}[单位比值自同构的周期点公式]\label{thm:xi-time-part69-unit-automorphism-periodic-point-formula}
\leanverified{paper\_xi\_time\_part69\_unit\_automorphism\_periodic\_point\_formula}
沿用定理 \ref{thm:xi-time-part69-localized-unit-automorphism-group-classification} 的记号。取
\[
q\in G_S^\times\setminus\{\pm1\},
\]
并将其写成既约形式
\[
q=\frac{a}{b},
\qquad
a\in\ZZ\setminus\{0\},
\qquad
b\in\NN,
\qquad
\gcd(a,b)=1,
\qquad
\operatorname{Supp}(|ab|)\subseteq S.
\]
令
\[
\alpha_q:\Sigma_S\to\Sigma_S
\]
为与离散端乘法 \(x\mapsto qx\) 对偶的连续群自同构。则对任意 \(m\ge 1\)，有
\[
\boxed{
\operatorname{Fix}(\alpha_q^m)
\cong
\ZZ/|a^m-b^m|_{S^\perp}\ZZ,
\qquad
\#\operatorname{Fix}(\alpha_q^m)=|a^m-b^m|_{S^\perp}.
}
\]
\end{theorem}
\begin{proof}
由定理 \ref{thm:xi-time-part69-localized-unit-automorphism-group-classification}，\(\alpha_q\) 确为 \(\Sigma_S\) 的连续群自同构。再把定理 \ref{thm:xi-time-part69-localized-rational-endomorphism-fixedpoints-zeta} 应用于 \(r=q=a/b\)，便得到
\[
\#\operatorname{Fix}(\alpha_q^m)=|a^m-b^m|_{S^\perp},
\]
以及对应的有限循环群结构。证毕。
\end{proof}

\begin{theorem}[单位自同构的周期指数与 prime-register 的多项式级盲性]\label{thm:xi-time-part69-unit-automorphism-polynomial-primeblindness}
\leanpartial{paper\_xi\_time\_part69\_unit\_automorphism\_polynomial\_primeblindness}{当前 Lean 目标按给定签名为假：取 $a=b=1$ 时假设成立而结论要求正的下界，但左侧恒为 $0$。}
沿用定理 \ref{thm:xi-time-part69-unit-automorphism-periodic-point-formula} 的记号，并设
\[
q\neq \pm1.
\]
则有极限
\[
\boxed{
\lim_{m\to\infty}\frac1m\log \#\operatorname{Fix}(\alpha_q^m)
=
\log\max\{|a|,b\}.
}
\]
更精确地，存在常数 \(C=C(q,S)>0\)，使得对充分大的 \(m\) 都有
\[
\boxed{
C^{-1}\,\frac{\max\{|a|,b\}^{\,m}}{m^{|S|}}
\le
\#\operatorname{Fix}(\alpha_q^m)
\le
C\,\max\{|a|,b\}^{\,m}.
}
\]
因此，prime-register \(S\) 只能改变周期点计数的多项式级修正，而不能改变其指数主率。
\end{theorem}
\begin{proof}
由定理 \ref{thm:xi-time-part69-unit-automorphism-periodic-point-formula}，
\[
\#\operatorname{Fix}(\alpha_q^m)
=
|a^m-b^m|_{S^\perp}
=
\frac{|a^m-b^m|}{\prod_{p\in S}p^{v_p(a^m-b^m)}}.
\]
记
\[
M:=\max\{|a|,b\},
\qquad
\lambda:=\frac{\min\{|a|,b\}}{M}\in(0,1).
\]
由于 \(q\neq \pm1\)，故 \(|a|\neq b\)。若 \(M=|a|\)，则
\[
|a^m-b^m|
=
M^m\left|\operatorname{sgn}(a)^m-\lambda^m\right|;
\]
若 \(M=b\)，则
\[
|a^m-b^m|
=
M^m\left|(\operatorname{sgn}(a)\lambda)^m-1\right|.
\]
两种情形都表明：存在常数 \(c_0,C_0>0\)，使得对充分大的 \(m\)，有
\[
c_0 M^m\le |a^m-b^m|\le C_0 M^m.
\]

现估计 \(S\)-内素因子。对固定 \(p\in S\)：

若 \(p\mid ab\)，则由 \(\gcd(a,b)=1\) 可知 \(p\) 只能整除 \(a,b\) 之一，从而对任意 \(m\ge 1\) 都有
\[
v_p(a^m-b^m)=0.
\]
在这种情形下取 \(C_p:=0\)。
若 \(p\nmid ab\)，则 \(a/b\) 在 \((\ZZ/p\ZZ)^\times\) 中有定义。记其模 \(p\) 的乘法阶为 \(d_p\)。当 \(d_p\nmid m\) 时，显然
\[
v_p(a^m-b^m)=0.
\]
当 \(d_p\mid m\) 时，由标准的 lifting-the-exponent 估计（在 \(p=2\) 时使用其 dyadic 变体）可知，存在仅依赖于 \(p,a,b\) 的常数 \(C_p\)，使得
\[
v_p(a^m-b^m)\le C_p+v_p(m).
\]
因此对所有 \(p\in S\) 与所有 \(m\ge1\)，统一有
\[
v_p(a^m-b^m)\le C_p+v_p(m).
\]
于是
\[
\prod_{p\in S}p^{v_p(a^m-b^m)}
\le
\left(\prod_{p\in S}p^{C_p}\right)\prod_{p\in S}p^{v_p(m)}
\le
C_1\,m^{|S|},
\]
其中 \(C_1:=\prod_{p\in S}p^{C_p}\)。

把上述估计与 \(c_0 M^m\le |a^m-b^m|\le C_0 M^m\) 合并，即得
\[
\#\operatorname{Fix}(\alpha_q^m)
=
\frac{|a^m-b^m|}{\prod_{p\in S}p^{v_p(a^m-b^m)}}
\ge
\frac{c_0}{C_1}\,\frac{M^m}{m^{|S|}}
\]
以及
\[
\#\operatorname{Fix}(\alpha_q^m)\le C_0 M^m.
\]
重新吸收常数便得到所示双边界。两边取对数、除以 \(m\) 并令 \(m\to\infty\)，立得指数极限。证毕。
\end{proof}

\begin{theorem}[rank-\texorpdfstring{$1$}{1} 局部化世界中遍历与 mixing 完全塌缩为 \texorpdfstring{$q\neq\pm1$}{q!=±1}]\label{thm:xi-time-part69-unit-automorphism-ergodic-mixing-collapse}
\leanverified{paper\_xi\_time\_part69\_unit\_automorphism\_ergodic\_mixing\_collapse}
沿用定理 \ref{thm:xi-time-part69-localized-unit-automorphism-group-classification} 的记号。对任意 \(q\in G_S^\times\)，记
\[
\alpha_q:\Sigma_S\to\Sigma_S
\]
为对偶于乘法 \(x\mapsto qx\) 的连续群自同构。则在 Haar 概率测度意义下，
\[
\boxed{
\alpha_q\ \text{遍历}
\iff
\alpha_q\ \text{mixing}
\iff
q\neq \pm1.
}
\]
\end{theorem}
\begin{proof}
记 \(\mu_S\) 为 \(\Sigma_S\) 上的 Haar 概率测度。对每个 \(x\in G_S\)，记对应字符为
\[
\chi_x(\omega):=\omega(x)
\qquad(\omega\in \Sigma_S).
\]
则 Pontryagin 对偶关系给出
\[
\chi_x\circ \alpha_q^n=\chi_{q^n x}
\qquad(n\ge 0).
\]
因而对任意 \(x,y\in G_S\)，利用字符正交性有
\[
\int_{\Sigma_S}(\chi_x\circ \alpha_q^n)\,\overline{\chi_y}\,d\mu_S
=
\int_{\Sigma_S}\chi_{q^n x-y}\,d\mu_S
=
\begin{cases}
1,& q^n x=y,\\
0,& q^n x\neq y.
\end{cases}
\]

现设 \(q\neq \pm1\)。若存在 \(x,y\neq 0\) 与两个不同整数 \(m>n\) 使
\[
q^m x=y=q^n x,
\]
则
\[
(q^{m-n}-1)x=0.
\]
由于 \(G_S=\ZZ[S^{-1}]\subset\QQ\) 无挠，必有 \(q^{m-n}=1\)，与 \(q\neq \pm1\) 矛盾。故对任意非零 \(x,y\in G_S\)，方程 \(q^n x=y\) 至多成立一次；于是上式对所有充分大的 \(n\) 恒为 \(0\)。

由三角多项式在 \(L^2(\Sigma_S,\mu_S)\) 中稠密，便得到：对任意零均值函数 \(f,g\in L^2(\Sigma_S,\mu_S)\)，有
\[
\langle f\circ \alpha_q^n,g\rangle\longrightarrow 0.
\]
故 \(\alpha_q\) mixing，从而也遍历。

反之，若 \(q=1\)，则 \(\alpha_q=\mathrm{id}\)，显然不遍历。若 \(q=-1\)，取任意非零 \(x\in G_S\)。则
\[
\chi_x\circ \alpha_{-1}^{2n}=\chi_x
\qquad(n\ge 0),
\]
故 \(\alpha_{-1}\) 不是 mixing。并且
\[
\Re\chi_x\circ \alpha_{-1}
=
\Re\chi_{-x}
=
\Re\chi_x,
\]
而 \(\Re\chi_x\) 非常值，故 \(\alpha_{-1}\) 亦不遍历。综上得到所示等价关系。证毕。
\end{proof}

\begin{corollary}[同一指数主率下存在无穷多 pairwise non-isomorphic 的 mixing 局部化系统]\label{cor:xi-time-part69-same-exponential-rate-infinite-mixing-family}
\leanverified{paper\_xi\_time\_part69\_same\_exponential\_rate\_infinite\_mixing\_family}
固定任意既约有理数
\[
q=\frac{a}{b}>1,
\qquad
a,b\in\NN,
\qquad
\gcd(a,b)=1.
\]
对任意有限素数集
\[
S\supseteq \operatorname{Supp}(ab),
\]
记
\[
\Sigma_S:=\widehat{\ZZ[S^{-1}]},
\qquad
\alpha_q:\Sigma_S\to\Sigma_S
\]
为与离散端乘法 \(x\mapsto qx\) 对偶的连续群自同构。则：
\begin{enumerate}
  \item 对所有这类 \(S\)，都有
  \[
  \lim_{m\to\infty}\frac1m\log \#\operatorname{Fix}(\alpha_q^m)=\log a;
  \]
  \item 若 \(S\neq T\)，则
  \[
  \Sigma_S\not\cong \Sigma_T;
  \]
  \item 因而对每个固定的 \(q>1\)，存在无穷多个 pairwise non-isomorphic 的 compact connected mixing 系统
  \[
  (\Sigma_S,\alpha_q),
  \]
  它们具有完全相同的周期指数主率 \(\log a\)。
\end{enumerate}
\end{corollary}
\begin{proof}
第 \(1\) 条是定理 \ref{thm:xi-time-part69-unit-automorphism-polynomial-primeblindness} 的直接特例，因为此时
\[
\max\{|a|,b\}=a.
\]
第 \(2\) 条由定理 \ref{thm:xi-time-part56-localized-pontryagin-isomorphism-rigidity} 立即得到。

对第 \(3\) 条，只需注意：包含 \(\operatorname{Supp}(ab)\) 的有限素数集 \(S\) 有无穷多个；由定理 \ref{thm:xi-time-part69-unit-automorphism-ergodic-mixing-collapse}，对每个这样的 \(S\)，因 \(q>1\) 故 \(\alpha_q\) mixing；再结合前两条，便得到一列两两不同构而指数主率完全相同的局部化动力系统。证毕。
\end{proof}

\endinput

\paragraph{有限角色审计的单圆塌缩与 prime-register 尾部失明}
在局部化 solenoid 的连续自同构、周期点指数与 prime-register 的多项式级盲性已经闭合之后，有限角色族的审计能力还可进一步压缩为一条双重塌缩链：在连通方向上，任意有限角色族都退化为单一 primitive 圆相位；在全不连通方向上，它们对 prime-register 核只能读取有限深度截断，因此有限 trigonometric 审计代数同时失去高维相位自由度与无限 \(p\)-进尾部。

\begin{theorem}[局部化 solenoid 上任意有限角色族都塌缩为单一 primitive 相位]\label{thm:xi-time-part69-finite-character-family-single-primitive-phase}
\leanverified{paper\_xi\_time\_part69\_finite\_character\_family\_single\_primitive\_phase}
设 \(S\subset\mathbb P\) 为有限素数集，并记
\[
G_S:=\ZZ[S^{-1}],
\qquad
\Sigma_S:=\widehat{G_S}.
\]
任取有限族连续角色
\[
\chi_1,\dots,\chi_r\in \widehat{\Sigma_S}\cong G_S.
\]
则存在一个角色 \(\chi_\ast\in\widehat{\Sigma_S}\) 与整数
\[
n_1,\dots,n_r\in\ZZ
\]
使得
\[
\chi_j=\chi_\ast^{\,n_j}
\qquad(1\le j\le r).
\]
若 \(\chi_1,\dots,\chi_r\) 不全为平凡角色，则 \(\chi_\ast\) 可取为 primitive 角色，并且在此规范下唯一到反演；相应地，整数族 \((n_1,\dots,n_r)\) 自动满足
\[
\gcd(n_1,\dots,n_r)=1.
\]
\end{theorem}
\begin{proof}
在 Pontryagin 对偶同构
\(
\widehat{\Sigma_S}\cong G_S
\)
下，令 \(\chi_j\) 对应于 \(q_j\in G_S\)。记
\[
H:=\langle q_1,\dots,q_r\rangle\subseteq G_S\subseteq \QQ.
\]
取正整数 \(D\) 使得
\[
Dq_j\in\ZZ
\qquad(1\le j\le r).
\]
则 \(DH\) 是 \(\ZZ\) 的有限生成子群，故必为循环群。于是存在唯一整数 \(d\ge 0\) 使
\[
DH=d\ZZ.
\]

若 \(d=0\)，则 \(H=\{0\}\)，亦即全部 \(q_j=0\)，从而所有 \(\chi_j\) 都平凡。此时取 \(\chi_\ast=1\) 与 \(n_j=0\) 即可。

若 \(d>0\)，则
\[
H=\frac{d}{D}\ZZ.
\]
记
\[
q_\ast:=\frac{d}{D}\in G_S.
\]
则 \(H=\ZZ q_\ast\)，故存在整数 \(n_j\) 使
\[
q_j=n_j q_\ast
\qquad(1\le j\le r).
\]
令 \(\chi_\ast\) 为与 \(q_\ast\) 对应的角色，则
\[
\chi_j=\chi_{q_j}=\chi_{q_\ast}^{\,n_j}=\chi_\ast^{\,n_j}.
\]
若 \(q_\ast'\) 也是 \(H\) 的生成元，则 \(q_\ast'=\pm q_\ast\)，故对应角色只差反演 \(\chi_\ast^{-1}\)。最后，由
\[
H=\langle q_1,\dots,q_r\rangle=\langle n_1 q_\ast,\dots,n_r q_\ast\rangle
\]
可知 \(\langle n_1,\dots,n_r\rangle=\ZZ\)，从而 \(\gcd(n_1,\dots,n_r)=1\)。证毕。
\end{proof}

\begin{corollary}[有限相位审计的全部 Borel 结构都等价于单圆相位审计]\label{cor:xi-time-part69-finite-phase-audit-borel-collapse-single-circle}
沿用定理 \ref{thm:xi-time-part69-finite-character-family-single-primitive-phase} 的记号，并定义联合相位映射
\[
\boldsymbol{\chi}:\Sigma_S\to \TT^r,
\qquad
x\longmapsto \bigl(\chi_1(x),\dots,\chi_r(x)\bigr).
\]
则存在一个角色 \(\chi_\ast:\Sigma_S\to\TT\) 与整数指数单项式映射
\[
M:\TT\to \TT^r,
\qquad
M(z):=(z^{n_1},\dots,z^{n_r}),
\]
使得
\[
\boldsymbol{\chi}=M\circ \chi_\ast.
\]
从而对任意 Borel 集 \(B\subseteq \TT^r\)，都存在 Borel 集 \(B_\ast\subseteq \TT\) 使
\[
\boldsymbol{\chi}^{-1}(B)=\chi_\ast^{-1}(B_\ast).
\]
并且有限角色族所生成的 \(\sigma\)-代数恰满足
\[
\sigma(\chi_1,\dots,\chi_r)=\sigma(\chi_\ast).
\]
\end{corollary}
\begin{proof}
由定理 \ref{thm:xi-time-part69-finite-character-family-single-primitive-phase}，存在 \(\chi_\ast\) 与整数 \(n_1,\dots,n_r\) 使
\[
\chi_j=\chi_\ast^{\,n_j}.
\]
故对任意 \(x\in \Sigma_S\)，有
\[
\boldsymbol{\chi}(x)
=
\bigl(\chi_\ast(x)^{n_1},\dots,\chi_\ast(x)^{n_r}\bigr)
=
M(\chi_\ast(x)),
\]
即 \(\boldsymbol{\chi}=M\circ \chi_\ast\)。

任取 Borel 集 \(B\subseteq \TT^r\)，令
\[
B_\ast:=M^{-1}(B)\subseteq \TT.
\]
由于 \(M\) 连续，\(B_\ast\) 仍为 Borel 集，并且
\[
\boldsymbol{\chi}^{-1}(B)
=
\chi_\ast^{-1}(M^{-1}(B))
=
\chi_\ast^{-1}(B_\ast).
\]
这给出
\[
\sigma(\chi_1,\dots,\chi_r)\subseteq \sigma(\chi_\ast).
\]

若全部 \(\chi_j\) 平凡，则两侧 \(\sigma\)-代数都只含空集与全集，结论已得。下设并非全部平凡。由定理 \ref{thm:xi-time-part69-finite-character-family-single-primitive-phase}，
\[
\gcd(n_1,\dots,n_r)=1.
\]
故可取整数 \(a_1,\dots,a_r\) 使
\[
\sum_{j=1}^{r}a_j n_j=1.
\]
定义连续映射
\[
N:\TT^r\to \TT,
\qquad
N(z_1,\dots,z_r):=\prod_{j=1}^{r}z_j^{a_j}.
\]
则对任意 \(x\in \Sigma_S\)，有
\[
N(\boldsymbol{\chi}(x))
=
\prod_{j=1}^{r}\chi_j(x)^{a_j}
=
\chi_\ast(x)^{\sum_j a_j n_j}
=
\chi_\ast(x).
\]
于是
\[
\chi_\ast=N\circ \boldsymbol{\chi}.
\]
对任意 Borel 集 \(C\subseteq \TT\)，有
\[
\chi_\ast^{-1}(C)=\boldsymbol{\chi}^{-1}(N^{-1}(C)),
\]
从而
\[
\sigma(\chi_\ast)\subseteq \sigma(\chi_1,\dots,\chi_r).
\]
综上得到两侧 \(\sigma\)-代数相等。证毕。
\end{proof}

\begin{theorem}[有限角色族对 prime-register 核只能看到有限层截断]\label{thm:xi-time-part69-finite-character-prime-ledger-finite-depth}
\leanverified{paper\_xi\_time\_part69\_finite\_character\_prime\_ledger\_finite\_depth}
沿用定理 \ref{thm:cdim-star-z1s-dual-extension} 的短正合列
\[
0\longrightarrow K_S:=\prod_{p\in S}\ZZ_p
\longrightarrow \Sigma_S
\longrightarrow \TT
\longrightarrow 0.
\]
给定任意有限角色族 \(\chi_1,\dots,\chi_r\in \widehat{\Sigma_S}\cong G_S\)，令 \(q_j\in G_S\) 为其对应元素，并对每个 \(p\in S\) 定义
\[
m_{j,p}:=\max\{0,-v_p(q_j)\}.
\]
当 \(q_j\neq 0\) 时，等价地，\(q_j\) 可唯一写成
\[
q_j=a_j\prod_{p\in S}p^{-m_{j,p}},
\qquad
a_j\in\ZZ,
\qquad
v_p(a_j)=0\ \ (p\in S).
\]
若 \(q_j=0\)，则取全部 \(m_{j,p}=0\)。
再令
\[
M_p:=\max_{1\le j\le r}m_{j,p},
\qquad
U(\chi_1,\dots,\chi_r):=\prod_{p\in S}p^{M_p}\ZZ_p\subseteq K_S.
\]
则：
\begin{enumerate}
  \item 每个限制角色 \(\chi_j|_{K_S}\) 都在 \(U(\chi_1,\dots,\chi_r)\) 的陪集上恒定；
  \item 若记
  \[
  T_{\boldsymbol{\chi}}:=\{p\in S:\ M_p\ge 1\},
  \]
  则联合映射
  \[
  K_S\longrightarrow \TT^r,
  \qquad
  x\longmapsto \bigl(\chi_1(x),\dots,\chi_r(x)\bigr)
  \]
  必经由有限商
  \[
  K_S/U(\chi_1,\dots,\chi_r)
  \cong
  \prod_{p\in T_{\boldsymbol{\chi}}}\ZZ/p^{M_p}\ZZ
  \]
  因子化；
  \item 因而任何有限角色审计都严格失明于 \(K_S\) 的无限 \(p\)-进尾部。
\end{enumerate}
\end{theorem}
\begin{proof}
由定理 \ref{thm:cdim-star-z1s-dual-extension}，有互为 Pontryagin 对偶的短正合列
\[
0\longrightarrow \ZZ
\longrightarrow G_S
\longrightarrow \bigoplus_{p\in S}C_{p^\infty}
\longrightarrow 0
\]
与
\[
0\longrightarrow K_S
\longrightarrow \Sigma_S
\longrightarrow \TT
\longrightarrow 0.
\]
因此，\(\chi_j\) 在 \(K_S\) 上的限制正对应于 \(q_j\) 在商群
\[
G_S/\ZZ\cong \bigoplus_{p\in S}C_{p^\infty}
\]
中的像。

固定 \(j\) 与 \(p\in S\)。由 \(m_{j,p}=\max\{0,-v_p(q_j)\}\) 的定义，\(q_j\) 在 \(G_S/\ZZ\) 中的 \(p\)-主分量落在 \(C_{p^\infty}\) 的唯一 \(p^{m_{j,p}}\)-阶循环子群内。对偶化以后，这意味着 \(\chi_j|_{K_S}\) 的第 \(p\)-坐标分量只能通过有限商
\[
\ZZ_p\twoheadrightarrow \ZZ_p/p^{m_{j,p}}\ZZ_p\cong \ZZ/p^{m_{j,p}}\ZZ
\]
被读取；等价地，
\[
p^{m_{j,p}}\ZZ_p
\subseteq
\ker\bigl(\chi_j|_{\ZZ_p}\bigr).
\]
于是定义
\[
U_j:=\prod_{p\in S}p^{m_{j,p}}\ZZ_p\subseteq K_S,
\]
便有
\[
U_j\subseteq \ker(\chi_j|_{K_S}).
\]

对全部 \(j\) 取交，得到
\[
\bigcap_{j=1}^{r}U_j
=
\prod_{p\in S}p^{M_p}\ZZ_p
=
U(\chi_1,\dots,\chi_r).
\]
故每个 \(\chi_j|_{K_S}\) 都在 \(U(\chi_1,\dots,\chi_r)\) 的陪集上恒定，第 \(1\) 条成立。

由于 \(U(\chi_1,\dots,\chi_r)\) 是 \(K_S\) 的开子群，联合映射必经由商群 \(K_S/U(\chi_1,\dots,\chi_r)\) 因子化。又由各坐标逐素数分解，
\[
K_S/U(\chi_1,\dots,\chi_r)
\cong
\prod_{p\in T_{\boldsymbol{\chi}}}\ZZ/p^{M_p}\ZZ.
\]
这给出第 \(2\) 条。最后，所有更深层的 \(p\)-进尾部都落在 \(U(\chi_1,\dots,\chi_r)\) 内，因而不会改变任何有限角色的取值，第 \(3\) 条随即成立。证毕。
\end{proof}

\begin{theorem}[有限角色生成的 trigonometric 审计代数在一致闭包后恰为单圆代数]\label{thm:xi-time-part69-finite-trigonometric-audit-algebra-single-circle-finite-tail}
设
\[
\mathcal A(\chi_1,\dots,\chi_r)\subset C(\Sigma_S)
\]
为由有限角色族 \(\chi_1,\dots,\chi_r\) 及其复共轭生成的含幺 \(^\ast\)-代数。令 \(\chi_\ast\) 为定理 \ref{thm:xi-time-part69-finite-character-family-single-primitive-phase} 给出的生成角色，并令
\[
U:=U(\chi_1,\dots,\chi_r)\subseteq K_S
\]
为定理 \ref{thm:xi-time-part69-finite-character-prime-ledger-finite-depth} 的开子群。则成立：
\begin{enumerate}
  \item 代数恒等式
  \[
  \mathcal A(\chi_1,\dots,\chi_r)
  =
  \chi_\ast^{*}\bigl(\CC[z,z^{-1}]\bigr);
  \]
  \item 若 \(\chi_1,\dots,\chi_r\) 不全平凡，则
  \[
  \overline{\mathcal A(\chi_1,\dots,\chi_r)}^{\|\cdot\|_\infty}
  =
  \chi_\ast^{*}C(\TT)
  \cong
  C(\TT);
  \]
  若全部平凡，则
  \[
  \overline{\mathcal A(\chi_1,\dots,\chi_r)}^{\|\cdot\|_\infty}
  =
  \CC\mathbf 1;
  \]
  \item 由 \(\mathcal A(\chi_1,\dots,\chi_r)\) 生成的任意有限 Borel 分割的每个非空原子都包含某个 \(U\)-陪集；特别地，它不可能把 \(K_S\) 的点分离到有限精度以下。
\end{enumerate}
\end{theorem}
\begin{proof}
由定理 \ref{thm:xi-time-part69-finite-character-family-single-primitive-phase}，存在整数 \(n_1,\dots,n_r\) 使
\[
\chi_j=\chi_\ast^{\,n_j}.
\]
因此 \(\mathcal A(\chi_1,\dots,\chi_r)\) 的每个元素都是 \(\chi_\ast\) 的 Laurent 多项式，故有包含
\[
\mathcal A(\chi_1,\dots,\chi_r)
\subseteq
\chi_\ast^{*}\bigl(\CC[z,z^{-1}]\bigr).
\]

若全部 \(\chi_j\) 平凡，则两边都等于常数代数，结论立即成立。下设并非全部平凡。由定理 \ref{thm:xi-time-part69-finite-character-family-single-primitive-phase}，可取 \(\chi_\ast\) primitive，且
\[
\gcd(n_1,\dots,n_r)=1.
\]
故可取整数 \(a_1,\dots,a_r\) 使
\[
\sum_{j=1}^{r}a_j n_j=1.
\]
于是
\[
\chi_\ast=\prod_{j=1}^{r}\chi_j^{a_j}.
\]
由于负指数可由复共轭实现，故 \(\chi_\ast\in \mathcal A(\chi_1,\dots,\chi_r)\)。从而
\[
\chi_\ast^{*}\bigl(\CC[z,z^{-1}]\bigr)
\subseteq
\mathcal A(\chi_1,\dots,\chi_r),
\]
第 \(1\) 条得证。

对第 \(2\) 条，只需证明非平凡情形下 \(\chi_\ast\) 满射到 \(\TT\)。由于 \(G_S\) 无挠，\(\Sigma_S=\widehat{G_S}\) 是紧连通交换群；故 \(\chi_\ast(\Sigma_S)\) 是 \(\TT\) 的闭连通子群。又 \(\chi_\ast\) 非平凡，故其像不能为 \(\{1\}\)，从而
\[
\chi_\ast(\Sigma_S)=\TT.
\]
于是 \(\chi_\ast^{*}:C(\TT)\to C(\Sigma_S)\) 是等距单射。由 Stone--Weierstrass 定理，
\[
\overline{\CC[z,z^{-1}]}^{\|\cdot\|_\infty}=C(\TT),
\]
拉回后即得
\[
\overline{\mathcal A(\chi_1,\dots,\chi_r)}^{\|\cdot\|_\infty}
=
\chi_\ast^{*}C(\TT)
\cong
C(\TT).
\]

最后，由定理 \ref{thm:xi-time-part69-finite-character-prime-ledger-finite-depth}，每个生成元 \(\chi_j\) 都在 \(U\)-陪集上恒定。故 \(\mathcal A(\chi_1,\dots,\chi_r)\) 的每个元素，以及其一致闭包中的每个元素，也都在 \(U\)-陪集上恒定。于是由这些函数所生成的全部 Borel 集都是 \(U\)-饱和集，即都是 \(U\)-陪集的并。因而任意由该代数决定的有限 Borel 分割，其每个非空原子也必为 \(U\)-陪集的并；特别地，取其中任一点 \(x\)，整个陪集 \(xU\) 都包含在该原子中。故该分割不可能把 \(K_S\) 的点分离到比 \(U\) 更细的精度。证毕。
\end{proof}

\endinput

\paragraph{有限连续可见性、有限指数/有限商谱与连续自同态环的闭合压缩}
在局部化数系的对偶短正合列、有限商分类、连续自同态与有限角色审计已经分别获得闭式之后，还可把连续有限可见性、有限指数格、有限商轮廓与连续动力学几条线路进一步压缩为下列若干条直接后果。

\begin{theorem}[局部化数系对偶在有限连续读出下的完全塌缩]\label{thm:xi-time-part69c-localized-dual-finite-visible-collapse}
设 \(S\subset\mathbb P\) 为有限素数集，并记
\[
G_S:=\ZZ[S^{-1}],
\qquad
\Sigma_S:=\widehat{G_S}.
\]
则 \(\Sigma_S\) 是连通紧交换群。因而对任意有限离散集 \(A\) 与任意连续映射
\[
F:\Sigma_S\longrightarrow A,
\]
映射 \(F\) 必为常值。
\end{theorem}
\begin{proof}
由定理 \ref{thm:cdim-star-z1s-dual-extension}，\(\Sigma_S\) 为紧交换群，并落在短正合列
\[
0\longrightarrow \prod_{p\in S}\ZZ_p
\longrightarrow \Sigma_S
\longrightarrow \TT
\longrightarrow 0.
\]
另一方面，
\[
G_S=\ZZ[S^{-1}]\subset \QQ
\]
作为离散加法群无挠。按 Pontryagin 对偶的标准连通性判据，\(\Sigma_S\) 必连通。

现令 \(A\) 为有限离散集，\(F:\Sigma_S\to A\) 连续。则 \(F(\Sigma_S)\) 是 \(A\) 的连通子空间；而有限离散空间的连通子空间只能是单点，故 \(F(\Sigma_S)\) 仅含一个元素，从而 \(F\) 为常值。证毕。
\end{proof}

\begin{theorem}[局部化数系有限商阶谱的完全闭式]\label{thm:xi-time-part69c-localized-finite-quotient-order-spectrum}
\leanverified{paper\_xi\_time\_part69c\_localized\_finite\_quotient\_order\_spectrum}
沿用定理 \ref{thm:xi-time-part69c-localized-dual-finite-visible-collapse} 的记号，并定义
\[
\mathcal Q(G_S)
:=
\left\{
n\ge 1:\ \exists\ \text{满同态 }G_S\twoheadrightarrow \ZZ/n\ZZ
\right\}.
\]
则有严格等式
\[
\mathcal Q(G_S)
=
\left\{
n\ge 1:\ \gcd\!\left(n,\prod_{p\in S}p\right)=1
\right\}.
\]
等价地，\(G_S\) 的全部有限商恰为循环群
\[
\ZZ/n\ZZ
\qquad
\bigl(\gcd(n,\prod_{p\in S}p)=1\bigr).
\]
\end{theorem}
\begin{proof}
由定理 \ref{thm:xi-localized-finite-quotient-phase-layer-classification} 第 \(1\) 条，对任意 \(n\ge 1\) 都有
\[
G_S/nG_S\cong \ZZ/n_{S^\perp}\ZZ.
\]
若
\[
\gcd\!\left(n,\prod_{p\in S}p\right)=1,
\]
则 \(n_{S^\perp}=n\)，故自然商映射给出满同态
\[
G_S\twoheadrightarrow G_S/nG_S\cong \ZZ/n\ZZ,
\]
于是 \(n\in \mathcal Q(G_S)\)。

反之，若 \(n\in \mathcal Q(G_S)\)，则存在满同态
\[
G_S\twoheadrightarrow \ZZ/n\ZZ.
\]
这说明 \(\ZZ/n\ZZ\) 是 \(G_S\) 的有限商。再由定理
\ref{thm:xi-localized-finite-quotient-phase-layer-classification}
第 \(2\) 条，\(G_S\) 的每个有限商都必为阶数与 \(S\) 中所有素数互素的循环群。故 \(n\) 的全部素因子都不属于 \(S\)，亦即
\[
\gcd\!\left(n,\prod_{p\in S}p\right)=1.
\]
于是所示集合恒等式成立；最后一句只是同一定理的重述。证毕。
\end{proof}

\begin{theorem}[局部化数系的有限指数子群完全分类]\label{thm:xi-time-part69c-localized-finite-index-subgroup-classification}
沿用定理 \ref{thm:xi-time-part69c-localized-finite-quotient-order-spectrum} 的记号。则 \(G_S\) 的任意有限指数子群 \(H\le G_S\) 都唯一形如
\[
H=nG_S
\]
其中
\[
n\ge 1,
\qquad
\gcd\!\left(n,\prod_{p\in S}p\right)=1.
\]
并且其指数满足精确公式
\[
[G_S:nG_S]=n.
\]
\leanverified{paper\_xi\_time\_part69c\_localized\_finite\_index\_subgroup\_classification}
\end{theorem}
\begin{proof}
先设
\[
\gcd\!\left(n,\prod_{p\in S}p\right)=1.
\]
定义模 \(n\) 约化映射
\[
\rho_n:G_S\to \ZZ/n\ZZ,
\qquad
\rho_n\!\left(\frac{a}{\prod_{p\in S}p^{k_p}}\right)
:=
a\cdot\left(\prod_{p\in S}p^{k_p}\right)^{-1}\pmod n.
\]
由于每个 \(p\in S\) 在 \(\ZZ/n\ZZ\) 中可逆，\(\rho_n\) 良定义。它显然满射，因为整数子群 \(\ZZ\subset G_S\) 已满射到 \(\ZZ/n\ZZ\)。

现证
\[
\ker(\rho_n)=nG_S.
\]
包含 \(nG_S\subseteq \ker(\rho_n)\) 直接成立。反向地，若
\[
x=\frac{a}{b}\in G_S,
\qquad
b=\prod_{p\in S}p^{k_p},
\]
且 \(\rho_n(x)=0\)，则
\[
a\,b^{-1}\equiv 0\pmod n.
\]
由于 \(\gcd(b,n)=1\)，\(b^{-1}\) 在 \(\ZZ/n\ZZ\) 中可逆，故 \(a\equiv 0\pmod n\)。写 \(a=nc\)，便有
\[
x=\frac{nc}{b}=n\cdot \frac{c}{b}\in nG_S.
\]
因此
\[
G_S/nG_S\cong \ZZ/n\ZZ,
\]
从而
\[
[G_S:nG_S]=n.
\]

反之，设 \(H\le G_S\) 为有限指数子群，并记
\[
A:=G_S/H.
\]
则 \(A\) 为有限商。由定理
\ref{thm:xi-time-part69c-localized-finite-quotient-order-spectrum}，存在唯一满足
\[
\gcd\!\left(n,\prod_{p\in S}p\right)=1
\]
的正整数 \(n\)，使得
\[
A\cong \ZZ/n\ZZ.
\]
因而 \([G_S:H]=n\)，并且 \(nA=0\)。这说明对每个 \(x\in G_S\) 都有
\[
nx\in H,
\]
故
\[
nG_S\subseteq H.
\]
而第一部分已证明
\[
[G_S:nG_S]=n=[G_S:H].
\]
既然 \(nG_S\subseteq H\) 且两者指数相同，遂得
\[
H=nG_S.
\]
唯一性则由指数唯一确定。证毕。
\end{proof}

\begin{corollary}[有限指数子群格与 \(S\)-自由整除格的反同构及其子群 \texorpdfstring{$\zeta$}{zeta} 函数]\label{cor:xi-time-part69c-localized-finite-index-lattice-zeta}
沿用定理 \ref{thm:xi-time-part69c-localized-finite-index-subgroup-classification} 的记号。则映射
\[
n\longmapsto nG_S
\qquad
\left(
\gcd\!\left(n,\prod_{p\in S}p\right)=1
\right)
\]
把 \(S\)-自由正整数的整除偏序反同构到 \(G_S\) 的有限指数子群格。更具体地，对任意与 \(S\) 互素的正整数 \(m,n\)，有
\[
nG_S\subseteq mG_S
\quad\Longleftrightarrow\quad
m\mid n,
\]
并且
\[
nG_S\cap mG_S=\operatorname{lcm}(n,m)\,G_S,
\]
\[
nG_S+mG_S=\gcd(n,m)\,G_S.
\]
因此 \(G_S\) 的子群 \(\zeta\) 函数满足完全闭式
\[
\zeta_{G_S}^{\le}(s)
:=
\sum_{H\le_f G_S}[G_S:H]^{-s}
=
\sum_{\gcd(n,\prod_{p\in S}p)=1}n^{-s}
=
\zeta(s)\prod_{p\in S}(1-p^{-s}).
\]
\end{corollary}
\begin{proof}
先证包含判据。若 \(m\mid n\)，写 \(n=mk\)，则
\[
nG_S=m(kG_S)\subseteq mG_S,
\]
故 \(nG_S\subseteq mG_S\)。

反之，若 \(nG_S\subseteq mG_S\)，则特别有 \(n\in mG_S\)，故存在 \(q\in G_S\) 使
\[
n=mq.
\]
于是
\[
q=\frac{n}{m}\in G_S.
\]
但 \(n,m\) 都与 \(S\) 互素，故既约分数 \(n/m\) 的分母不含任何 \(S\)-中素数。要使 \(n/m\in G_S=\ZZ[S^{-1}]\)，其分母只能为 \(1\)。因此 \(n/m\in \ZZ\)，即 \(m\mid n\)。

现设
\[
K:=nG_S\cap mG_S.
\]
由定理 \ref{thm:xi-time-part69c-localized-finite-index-subgroup-classification}，存在唯一 \(S\)-自由正整数 \(k\) 使
\[
K=kG_S.
\]
由 \(K\subseteq nG_S\) 与 \(K\subseteq mG_S\)，按上面的包含判据得
\[
n\mid k,
\qquad
m\mid k,
\]
故 \(\operatorname{lcm}(n,m)\mid k\)。另一方面，
\[
\operatorname{lcm}(n,m)\,G_S\subseteq nG_S\cap mG_S=K,
\]
再次应用包含判据得
\[
k\mid \operatorname{lcm}(n,m).
\]
故 \(k=\operatorname{lcm}(n,m)\)，从而
\[
nG_S\cap mG_S=\operatorname{lcm}(n,m)\,G_S.
\]

再设
\[
L:=nG_S+mG_S.
\]
同理，\(L=\ell G_S\) 对某个唯一 \(S\)-自由正整数 \(\ell\) 成立。由
\[
nG_S\subseteq L,
\qquad
mG_S\subseteq L,
\]
可知
\[
\ell\mid n,
\qquad
\ell\mid m,
\]
故 \(\ell\mid \gcd(n,m)\)。另一方面，取整数 \(u,v\) 满足
\[
un+vm=\gcd(n,m).
\]
则对任意 \(x\in G_S\)，有
\[
\gcd(n,m)\,x=u(nx)+v(mx)\in nG_S+mG_S=L.
\]
因此
\[
\gcd(n,m)\,G_S\subseteq L=\ell G_S,
\]
故由包含判据得
\[
\ell\mid \gcd(n,m).
\]
另一方面，由 \(n,m\) 都是 \(\gcd(n,m)\) 的整数倍，显然
\[
nG_S\subseteq \gcd(n,m)\,G_S,
\qquad
mG_S\subseteq \gcd(n,m)\,G_S.
\]
从而
\[
L=nG_S+mG_S\subseteq \gcd(n,m)\,G_S,
\]
再次应用包含判据得到
\[
\gcd(n,m)\mid \ell.
\]
于是 \(\ell=\gcd(n,m)\)，从而
\[
nG_S+mG_S=\gcd(n,m)\,G_S.
\]

最后，由定理 \ref{thm:xi-time-part69c-localized-finite-index-subgroup-classification}，全部有限指数子群恰由
\[
H=nG_S
\qquad
\left(
\gcd\!\left(n,\prod_{p\in S}p\right)=1
\right)
\]
枚举，且其指数正为 \(n\)。故
\[
\zeta_{G_S}^{\le}(s)
=
\sum_{\gcd(n,\prod_{p\in S}p)=1}n^{-s}.
\]
按 Euler 乘积展开即得
\[
\sum_{\gcd(n,\prod_{p\in S}p)=1}n^{-s}
=
\prod_{p\notin S}(1-p^{-s})^{-1}
=
\zeta(s)\prod_{p\in S}(1-p^{-s}).
\]
证毕。
\end{proof}

\begin{corollary}[有限商轮廓对素数账本的完全恢复]\label{cor:xi-time-part69c-localized-finite-quotient-profile-recovers-primes}
对任意有限素数集 \(S,T\subset\mathbb P\)，有
\[
\mathcal Q(G_S)=\mathcal Q(G_T)
\quad\Longleftrightarrow\quad
S=T.
\]
等价地，对任意素数 \(p\) 都有
\[
p\in S
\quad\Longleftrightarrow\quad
\forall n\in \mathcal Q(G_S),\ p\nmid n.
\]
\end{corollary}
\begin{proof}
由定理 \ref{thm:xi-time-part69c-localized-finite-quotient-order-spectrum}，
\[
\mathcal Q(G_S)
=
\left\{
n\ge 1:\ \forall p\in S,\ p\nmid n
\right\}.
\]
因此 \(\mathcal Q(G_S)\) 恰记录了哪些素数被完全排除在全部有限商阶谱之外。若
\[
\mathcal Q(G_S)=\mathcal Q(G_T),
\]
则被排除的素数集合一致，故 \(S=T\)。反向蕴含显然。

第二条只是同一刻画按单个素数逐点展开：若 \(p\in S\)，则由上式 \(p\) 不整除任意 \(n\in\mathcal Q(G_S)\)；反之若 \(p\notin S\)，则 \(n=p\) 已属于 \(\mathcal Q(G_S)\)。证毕。
\end{proof}

\begin{theorem}[局部化数系的连续自同态环完全由局部化环穷尽]\label{thm:xi-time-part69c-localized-dual-endomorphism-ring-exhaustion}
沿用定理 \ref{thm:xi-time-part69c-localized-dual-finite-visible-collapse} 的记号。则加法群 \(G_S\) 的端同态环满足
\[
\operatorname{End}(G_S)\cong \ZZ[S^{-1}],
\]
其中同构由
\[
q\longmapsto m_q,
\qquad
m_q(x):=qx
\]
给出。因而其对偶紧群的连续端同态环满足
\[
\operatorname{End}_{\mathrm{cts}}(\Sigma_S)\cong \ZZ[S^{-1}].
\]
进一步，连续自同构群满足
\[
\operatorname{Aut}_{\mathrm{cts}}(\Sigma_S)
\cong
\operatorname{Aut}(G_S)
\cong
\left\{
\pm\prod_{p\in S}p^{n_p}: n_p\in\ZZ
\right\}.
\]
\end{theorem}
\begin{proof}
定理 \ref{thm:xi-localized-endomorphism-ring-solenoid-indecomposable} 的前两条已经给出
\[
\operatorname{End}(G_S)\cong G_S=\ZZ[S^{-1}],
\]
并且该同构正由乘法映射 \(q\mapsto m_q\) 实现。Pontryagin 对偶把紧交换群上的连续端同态环识别为离散端端同态环的反环，故
\[
\operatorname{End}_{\mathrm{cts}}(\Sigma_S)\cong \operatorname{End}(G_S)^{\mathrm{op}}.
\]
由于 \(\ZZ[S^{-1}]\) 为交换环，其反环与自身一致，于是
\[
\operatorname{End}_{\mathrm{cts}}(\Sigma_S)\cong \ZZ[S^{-1}].
\]

最后，由定理 \ref{thm:xi-time-part69-localized-unit-automorphism-group-classification}，
\[
\operatorname{Aut}_{\mathrm{cts}}(\Sigma_S)
\cong
G_S^\times
=
\left\{
\pm\prod_{p\in S}p^{n_p}: n_p\in\ZZ
\right\}.
\]
而离散端自同构正对应于 \(\ZZ[S^{-1}]\) 的可逆乘法元，故同一公式也描述 \(\operatorname{Aut}(G_S)\)。证毕。
\end{proof}

\noindent
由此，局部化数系的对偶对象在连续有限值读出层面完全塌缩为单点；而在有限指数与有限商层面，其全部可见结构又分别被精确压缩为 \(S\)-自由整数的整除半格、对应的子群 \(\zeta\) 函数 Euler 因子以及有限商阶谱。进一步，在连续自同态动力学层面，全部端同态与自同构又完全坍缩为局部化环 \(\ZZ[S^{-1}]\) 及其单位群本身。于是，素数账本、有限指数格、有限商轮廓与连续动力学四者被同一局部化数据闭合地锁定。

\paragraph{局部化态射矩阵、同构刚性与 Gödel--Tate 仿射动力学的终端分类}
在局部化数系之间的 \(\operatorname{Hom}\)-分类、有限局部化直和的稀疏矩阵态射结构，以及 Gödel--Tate 仿射作用的唯一不动点、闭不变核与轨道闭包都已建立之后，还可把有限局部化账本与 \(2\)-进仿射表示两条主线进一步压缩为下列一组分类性结论。

\begin{theorem}[局部化数系 \texorpdfstring{$G_S=\ZZ[S^{-1}]$}{G\_S=Z[S^{-1}]} 之间的 \texorpdfstring{$\operatorname{Hom}$}{Hom}/\texorpdfstring{$\operatorname{Aut}$}{Aut} 完全分类]\label{thm:xi-time-part69d-localized-hom-aut-complete-classification}
\leanverified{paper\_xi\_time\_part69d\_localized\_hom\_aut\_complete\_classification}
设 \(S,T\subset\mathbb P\) 为有限素数集，并记
\[
G_S:=\ZZ[S^{-1}],
\qquad
G_T:=\ZZ[T^{-1}]
\]
为加法群。则任意群同态 \(\phi:G_S\to G_T\) 都唯一写成
\[
\phi(x)=qx
\qquad (x\in G_S)
\]
的形式，其中 \(q\in \QQ\)。进一步，
\[
\operatorname{Hom}(G_S,G_T)\cong
\begin{cases}
G_T,& S\subseteq T,\\[1mm]
0,& S\not\subseteq T,
\end{cases}
\]
其中同构由 \(q\mapsto (x\mapsto qx)\) 给出。特别地，
\[
\operatorname{End}(G_S)\cong G_S
\]
作为环同构成立，而自同构群满足
\[
\operatorname{Aut}(G_S)\cong G_S^\times
=
\left\{
\pm \prod_{p\in S}p^{n_p}:\ n_p\in\ZZ,\ \text{仅有限个 }n_p\neq 0
\right\}.
\]
\end{theorem}
\begin{proof}
定理 \ref{thm:xi-localized-integers-cross-hom-classification} 已给出
\[
\operatorname{Hom}(G_S,G_T)\cong
\begin{cases}
G_T,& S\subseteq T,\\[1mm]
0,& S\not\subseteq T.
\end{cases}
\]
其证明机制同时说明每个同态都由 \(\phi(1)\) 唯一决定。事实上，设 \(\phi:G_S\to G_T\) 为任意群同态，并记 \(q:=\phi(1)\in G_T\subset \QQ\)。对任意
\[
x=\frac{a}{b}\in G_S,
\qquad
a\in\ZZ,
\qquad
b=\prod_{p\in S}p^{m_p},
\]
在加法群 \(G_S\) 中有 \(bx=a\)。施加 \(\phi\) 得
\[
b\,\phi(x)=\phi(a)=aq.
\]
由于 \(G_T\subset \QQ\) 为无挠群，上式在 \(\QQ\) 中唯一解为
\[
\phi(x)=\frac{aq}{b}=qx.
\]
故全部同态均唯一来自有理数乘法。

当 \(S\subseteq T\) 时，\(q\in G_T\) 必满足 \(qG_S\subseteq G_T\)，故得到全部同态；当 \(S\not\subseteq T\) 时，定理 \ref{thm:xi-localized-integers-cross-hom-classification} 说明只能有零同态。

再由定理 \ref{thm:xi-localized-integers-endomorphism-automorphism-explicit}，
\[
\operatorname{End}(G_S)\cong G_S
\]
且
\[
\operatorname{Aut}(G_S)\cong G_S^\times.
\]
而局部化环 \(\ZZ[S^{-1}]\) 的单位恰为
\[
\left\{
\pm \prod_{p\in S}p^{n_p}:\ n_p\in\ZZ,\ \text{仅有限个 }n_p\neq 0
\right\},
\]
于是结论成立。
\end{proof}

\begin{theorem}[局部化数系紧对偶的连续自同态环完全显式]\label{thm:xi-time-part69d-localized-solenoid-endomorphism-ring}
设 \(S\subset\mathbb P\) 为有限素数集，并记
\[
G_S:=\ZZ[S^{-1}],
\qquad
\Sigma_S:=\widehat{G_S}.
\]
则存在自然环同构
\[
\operatorname{End}_{\mathrm{cts}}(\Sigma_S)\cong \ZZ[S^{-1}].
\]
\end{theorem}
\begin{proof}
Pontryagin 对偶把连续自同态反变地送到离散侧自同态，因此给出自然环反同构
\[
\operatorname{End}_{\mathrm{cts}}(\Sigma_S)
\cong
\operatorname{End}_{\ZZ}(G_S)^{\mathrm{op}}.
\]
而定理 \ref{thm:xi-time-part69d-localized-hom-aut-complete-classification} 与定理 \ref{thm:xi-localized-integers-endomorphism-automorphism-explicit} 已给出
\[
\operatorname{End}_{\ZZ}(G_S)\cong \ZZ[S^{-1}]
\]
作为环同构成立。由于 \(\ZZ[S^{-1}]\) 为交换环，其对立环与自身相同，故结论成立。证毕。
\end{proof}
\leanverified{paper\_xi\_time\_part69d\_localized\_solenoid\_endomorphism\_ring}

\begin{corollary}[局部化数系紧对偶的连续自同构群正是 \texorpdfstring{$S$}{S}-unit 群]\label{cor:xi-time-part69d-localized-solenoid-automorphism-sunits}
沿用定理 \ref{thm:xi-time-part69d-localized-solenoid-endomorphism-ring} 的记号，则
\[
\operatorname{Aut}_{\mathrm{cts}}(\Sigma_S)
\cong
\ZZ[S^{-1}]^\times
=
\left\{
\pm\prod_{p\in S}p^{n_p}: n_p\in\ZZ,\ \text{仅有限个 }n_p\neq 0
\right\}.
\]
\end{corollary}
\begin{proof}
Pontryagin 对偶在自同构层面给出自然同构
\[
\operatorname{Aut}_{\mathrm{cts}}(\Sigma_S)\cong \operatorname{Aut}_{\ZZ}(G_S).
\]
再由定理 \ref{thm:xi-time-part69d-localized-hom-aut-complete-classification} 与定理 \ref{thm:xi-localized-integers-endomorphism-automorphism-explicit}，
\[
\operatorname{Aut}_{\ZZ}(G_S)\cong \ZZ[S^{-1}]^\times.
\]
这正是所述带符号的 \(S\)-unit 群。证毕。
\end{proof}

\begin{theorem}[有限局部化直和的矩阵 \texorpdfstring{$\operatorname{Hom}$}{Hom} 分类与同构判别]\label{thm:xi-time-part69d-localized-directsum-matrix-isomorphism-criterion}
\leanverified{paper\_xi\_time\_part69d\_localized\_directsum\_matrix\_isomorphism\_criterion}
设
\[
\mathbf S=(S_1,\dots,S_r),
\qquad
\mathbf T=(T_1,\dots,T_s)
\]
为有限素数集族，并记
\[
G_{\mathbf S}:=\bigoplus_{i=1}^{r}G_{S_i},
\qquad
G_{\mathbf T}:=\bigoplus_{j=1}^{s}G_{T_j}.
\]
则任意群同态 \(\Phi:G_{\mathbf S}\to G_{\mathbf T}\) 唯一对应于一个 \(s\times r\) 矩阵
\[
Q=(q_{ji})
\]
满足
\[
q_{ji}\in
\begin{cases}
G_{T_j},& S_i\subseteq T_j,\\
0,& S_i\not\subseteq T_j,
\end{cases}
\]
且作用为
\[
\Phi(x_1,\dots,x_r)=
\left(
\sum_{i=1}^{r}q_{1i}x_i,\ \dots,\ \sum_{i=1}^{r}q_{si}x_i
\right).
\]
因而
\[
\operatorname{Hom}(G_{\mathbf S},G_{\mathbf T})
\cong
\bigoplus_{j=1}^{s}\bigoplus_{i=1}^{r}\operatorname{Hom}(G_{S_i},G_{T_j}).
\]
进一步，
\[
G_{\mathbf S}\cong G_{\mathbf T}
\iff
r=s\ \text{且在重排后 }S_i=T_i\ \ (1\le i\le r).
\]
\end{theorem}
\begin{proof}
前两条正是定理 \ref{thm:xi-localized-directsum-hom-sparse-matrix} 的内容。下面证明同构判别。

若在重排后 \(S_i=T_i\)，则逐分量恒等映射的直和显然给出
\(
G_{\mathbf S}\cong G_{\mathbf T}
\)。

反过来，设
\[
\Phi:G_{\mathbf S}\xrightarrow{\sim} G_{\mathbf T}
\]
为一同构，\(\Psi=\Phi^{-1}\)。对任意有限素数集 \(U\)，记
\[
A_U:=\bigoplus_{\,i:\,S_i=U}G_{S_i}\subset G_{\mathbf S},
\qquad
B_U:=\bigoplus_{\,j:\,T_j=U}G_{T_j}\subset G_{\mathbf T}.
\]
在有限集合族
\(
\{S_1,\dots,S_r,T_1,\dots,T_s\}
\)
中取一个关于包含关系极大的素数集 \(U\)。由单对象 \(\operatorname{Hom}\)-分类，若 \(S_i=U\)，则 \(G_{S_i}\) 只能非零映入那些满足 \(U\subseteq T_j\) 的目标分量；而 \(U\) 极大迫使这等价于 \(T_j=U\)。因此
\[
\Phi(A_U)\subseteq B_U.
\]
同理，
\[
\Psi(B_U)\subseteq A_U.
\]
于是 \(\Phi|_{A_U}:A_U\to B_U\) 为同构。另一方面，\(A_U\) 与 \(B_U\) 都是若干个 \(G_U\) 的有限直和，而每个 \(G_U\) 的 \(\QQ\)-秩均为 \(1\)。故
\[
\operatorname{rank}_{\QQ}(A_U)=\#\{i:\ S_i=U\},
\qquad
\operatorname{rank}_{\QQ}(B_U)=\#\{j:\ T_j=U\}.
\]
同构保持 \(\QQ\)-秩，从而
\[
\#\{i:\ S_i=U\}=\#\{j:\ T_j=U\}.
\]
删去两侧全部等于 \(U\) 的分量后，对剩余有限集合族重复同样论证。由于待比较的素数集总数有限，归纳终止并得到两族完全相同的多重集。故
\[
r=s
\]
且经重排后
\[
S_i=T_i\qquad(1\le i\le r).
\]
结论成立。
\end{proof}

\begin{theorem}[Gödel--Tate 仿射表示的中心化子完全描述]\label{thm:xi-time-part69d-godel-tate-affine-centralizer-complete}
沿用定理 \ref{thm:killo-godel-semidir-affine-2adic} 与定理 \ref{thm:xi-godel-tate-affine-closed-subgroup-classification} 的记号。记
\[
T:=T_2(E_{\min})\cong \ZZ_2^2,
\qquad
H:=\mathcal R\rtimes_\Sigma \NN,
\]
\[
\Psi(r,k)(x):=B^k x+\operatorname{ev}_B(r),
\qquad
\operatorname{ev}_B(r)=\sum_{t\ge 1}r(p_t)B^{t-1}v\in L_v:=\ZZ_2v.
\]
则 \(\Psi(H)\) 在连续仿射自映射半群
\(
\mathrm{Aff}_{\mathrm{cts}}(T)
\)
中的中心化子恰为
\[
\operatorname{Cent}_{\mathrm{Aff}_{\mathrm{cts}}(T)}(\Psi(H))
=
\left\{
x\mapsto Ax:\ A\in \operatorname{End}_{\ZZ_2}(T),\ A|_{L_v}=\mathrm{id}_{L_v}
\right\}.
\]
若固定基底 \((v,w)\) 使
\(
T=\ZZ_2v\oplus \ZZ_2w
\)，则上述中心化子恰对应于矩阵集合
\[
\left\{
\begin{pmatrix}
1 & b\\
0 & d
\end{pmatrix}
:\ b,d\in\ZZ_2
\right\}
\subset M_2(\ZZ_2),
\]
而其可逆部分，即
\(
\mathrm{Aff}_{\mathrm{cts}}^\times(T)
\)
中的中心化子，为
\[
\left\{
\begin{pmatrix}
1 & b\\
0 & d
\end{pmatrix}
:\ b\in\ZZ_2,\ d\in\ZZ_2^\times
\right\}.
\]
\end{theorem}
\begin{proof}
任取连续仿射映射
\[
\Phi(x)=Ax+t,
\qquad
A\in \operatorname{End}_{\ZZ_2}(T),
\qquad
t\in T,
\]
并要求它与一切 \(\Psi(r,k)\) 交换。

先令 \(k=0\)。此时
\[
\Psi(r,0)(x)=x+\operatorname{ev}_B(r).
\]
交换关系给出
\[
\Phi\bigl(x+\operatorname{ev}_B(r)\bigr)
=
Ax+A\operatorname{ev}_B(r)+t,
\]
\[
\Psi(r,0)(\Phi(x))
=
Ax+t+\operatorname{ev}_B(r).
\]
故对一切 \(r\in\mathcal R\) 都有
\[
A\operatorname{ev}_B(r)=\operatorname{ev}_B(r).
\]
而定理 \ref{thm:xi-godel-tate-affine-closed-subgroup-classification} 的证明已经给出
\[
\overline{\operatorname{ev}_B(\mathcal R)}=L_v.
\]
由 \(A\) 的连续性，推出
\[
A|_{L_v}=\mathrm{id}_{L_v}.
\]

再令 \(r=0\)。对任意 \(k\ge 1\)，交换关系给出
\[
\Phi(B^k x)=AB^k x+t,
\qquad
\Psi(0,k)(\Phi(x))=B^kAx+B^k t.
\]
由于 \(B^k=2^{Lk}\) 是标量，故 \(AB^k=B^kA\)。于是必须有
\[
t=B^k t
\qquad (\forall k\ge 1).
\]
取 \(k=1\)，即
\[
(1-B)t=0.
\]
而 \(B=2^L\)，故
\[
1-B=1-2^L\in \ZZ_2^\times
\]
为单位，从而 \(t=0\)。

因此任何中心化元都必为线性映射 \(x\mapsto Ax\)，且满足 \(A|_{L_v}=\mathrm{id}_{L_v}\)。

反过来，若 \(A\in \operatorname{End}_{\ZZ_2}(T)\) 且 \(A|_{L_v}=\mathrm{id}_{L_v}\)，则对任意 \((r,k)\in H\) 与 \(x\in T\)，有
\[
A(\Psi(r,k)(x))
=
A(B^k x+\operatorname{ev}_B(r))
=
B^kAx+\operatorname{ev}_B(r)
=
\Psi(r,k)(Ax),
\]
故 \(x\mapsto Ax\) 确实中心化 \(\Psi(H)\)。

最后，在基底 \((v,w)\) 下，条件 \(A(v)=v\) 恰意味着矩阵第一列为 \((1,0)^\top\)，故
\[
A=
\begin{pmatrix}
1 & b\\
0 & d
\end{pmatrix},
\qquad
b,d\in\ZZ_2.
\]
而该矩阵可逆当且仅当 \(d\in\ZZ_2^\times\)。结论成立。
\end{proof}

\begin{theorem}[Gödel--Tate 仿射动力系统中非平凡周期轨道完全不存在]\label{thm:xi-time-part69d-godel-tate-affine-no-nontrivial-periodic-orbits}
沿用定理 \ref{thm:xi-time-part69d-godel-tate-affine-centralizer-complete} 的记号。对任意 \((r,k)\in H\)，考虑仿射映射
\[
\Psi(r,k):T\to T,
\qquad
x\mapsto B^k x+\operatorname{ev}_B(r).
\]
则其周期点结构满足：
\begin{enumerate}
  \item 若 \(k\ge 1\)，则 \(\Psi(r,k)\) 恰有唯一不动点
  \[
  x_{r,k}=(1-B^k)^{-1}\operatorname{ev}_B(r),
  \]
  且它也是 \(\Psi(r,k)\) 的唯一周期点；特别地不存在任何真周期 \(n>1\) 轨道。
  \item 若 \(k=0\) 且 \((r,0)\neq e\)，则 \(\Psi(r,0)\) 没有任何周期点。
\end{enumerate}
因此，\(H\) 的任一非单位元素在该 \(2\)-进仿射作用下都不产生非平凡周期轨道。
\end{theorem}
\begin{proof}
若 \(k\ge 1\)，则定理 \ref{thm:xi-time-part9g-affine-semigroup-unique-fixed-point} 已给出唯一不动点
\[
x_{r,k}=(1-B^k)^{-1}\operatorname{ev}_B(r).
\]
记 \(\Phi:=\Psi(r,k)\)。由不动点方程
\[
x_{r,k}=B^k x_{r,k}+\operatorname{ev}_B(r)
\]
可得对任意 \(x\in T\)
\[
\Phi(x)-x_{r,k}=B^k(x-x_{r,k}).
\]
归纳即得
\[
\Phi^n(x)-x_{r,k}=B^{kn}(x-x_{r,k})
\qquad (n\ge 1).
\]
若 \(x\) 为周期点，则存在 \(n\ge 1\) 使 \(\Phi^n(x)=x\)，故
\[
(1-B^{kn})(x-x_{r,k})=0.
\]
而
\[
1-B^{kn}=1-2^{Lkn}\in \ZZ_2^\times
\]
仍为单位，因此必有 \(x=x_{r,k}\)。故唯一周期点就是唯一不动点。

若 \(k=0\)，则
\[
\Psi(r,0)(x)=x+\operatorname{ev}_B(r).
\]
若 \((r,0)\neq e\)，则由表示 \(\Psi\) 的忠实性可知
\[
\operatorname{ev}_B(r)\neq 0.
\]
若存在周期点 \(x\)，则对某个 \(n\ge 1\) 有
\[
x+n\operatorname{ev}_B(r)=x,
\]
即
\[
n\operatorname{ev}_B(r)=0.
\]
但 \(T\cong \ZZ_2^2\) 是无挠 \(\ZZ_2\)-模，故这不可能发生。于是 \(\Psi(r,0)\) 没有任何周期点。
\end{proof}

\noindent
由此，有限局部化账本与 \(2\)-进 Gödel--Tate 仿射表示两条主线都达到分类级终点：在代数侧，单对象与有限直和对象的全部态射已被局部化乘法与稀疏矩阵完全穷尽；在动力学侧，仿射表示的外部对称只剩下对地址线 \(L_v\) 恒等的线性中心化子，而一切非单位仿射元都不产生任何非平凡周期轨道。于是，素数账本的线性可见性与 Gödel 半直积的 \(2\)-进动力学刚性在同一框架中同时闭合为严格的终端分类。

\endinput


\endinput

\paragraph{六模拓扑相图、精确暗块、boundary 奇偶盲性与纯鸽巢截断}
在最大非可缩纤维的模 \(6\) 比例极限、审计偶窗最小退化扇区的奇偶同伦相型、window-\(6\) 的 boundary sheet-parity 中心分裂，以及黄金分支在 Fibonacci 分母上的 \(\mathsf T\)/\(D^\ast\) 严格奇偶双相都已固定之后，还可进一步抽取出下列四条跨链后果。它们分别刻画精确暗块与顶端非可缩纤维的非耦合共存、根字典对 boundary 奇偶三格的严格盲性、审计偶窗最弱退化扇区在阿贝尔化中的精确 Fibonacci 份额，以及最小退化 Fibonacci 律一旦试图脱离审计窗口便会遭遇的首个纯计数截断。

\begin{corollary}[精确暗块与拓扑主纤维的非耦合共存]\label{cor:xi-time-part69c-darkblock-topfiber-decoupling}
沿用推论 \ref{cor:xi-binfold-zero-block-collision-gap-decoupling} 与定理 \ref{thm:xi-time-part62db-maxfiber-homotopy-phase-mod6} 的记号，记
\[
D_m:=\max_{x\in X_m}d_m(x),
\qquad
\widetilde D_m:=\max\{d_m(x):x\in X_m,\ |K_x|\not\simeq *\},
\qquad
\bar d_m:=\frac{2^m}{F_{m+2}}.
\]
若
\[
m\equiv 4\pmod 6,
\]
则同时成立
\[
|\mathcal Z_m|\ge F_{(m+2)/2},
\qquad
\widetilde D_m=D_m,
\]
并且
\[
\liminf_{\substack{m\to\infty\\ m\equiv 4\!\!\!\!\pmod 6}}
\frac{D_m}{\bar d_m}
\ge
1+C_\varphi.
\]
因此，半阶指数规模的精确暗块与不发生拓扑折损的顶端非可缩纤维，可以在同一分辨率相上严格共存。
\end{corollary}
\begin{proof}
当 \(m\equiv 4\pmod 6\) 时，\(m+2\equiv 0\pmod 6\)。故由推论 \ref{cor:xi-binfold-zero-block-collision-gap-decoupling}，
\[
|\mathcal Z_m|\ge F_{(m+2)/2}.
\]
同一同余类上，定理 \ref{thm:xi-time-part62db-maxfiber-homotopy-phase-mod6} 直接给出
\[
\widetilde D_m=D_m.
\]
另一方面，结论 \ref{con:xi-time-part9m-cphi-forces-collision-gap-and-peak-bias} 已证明
\[
\liminf_{m\to\infty}\frac{M_m}{\overline d_m}\ge 1+C_\varphi.
\]
将其中记号 \(M_m,\overline d_m\) 分别改写为 \(D_m,\bar d_m\)，再限制到子序列 \(m\equiv 4\pmod 6\)，便得到第三式。证毕。
\end{proof}

\begin{theorem}[window-\texorpdfstring{$6$}{6} 根字典的 boundary 盲性余维恰为 \texorpdfstring{$3$}{3}]\label{thm:xi-time-part69c-window6-root-boundary-blind-codim3}
沿用定理 \ref{thm:xi-window6-cyclic-kernel-bc3-root-datum} 与定理 \ref{thm:xi-window6-center-exact-splitting-bdry-cyclic} 的记号。通过定理 \ref{thm:xi-foldbin-gauge-representation-distribution-law-m6-spectrum} 的直积分解，识别
\[
\bigl(G_6^{\mathrm{bin}}\bigr)^{\mathrm{ab}}
\cong
(\ZZ/2\ZZ)^{X_6}.
\]
令
\[
\rho:
\bigl(G_6^{\mathrm{bin}}\bigr)^{\mathrm{ab}}
\cong
(\ZZ/2\ZZ)^{X_6}
\longrightarrow
(\ZZ/2\ZZ)^{X_6^{\mathrm{cyc}}}
\]
为忘去 \(X_6^{\mathrm{bdry}}\) 三个坐标的自然投影，并记
\[
\iota_{\mathrm{ab}}:
Z\!\bigl(G_6^{\mathrm{bin}}\bigr)
\hookrightarrow
\bigl(G_6^{\mathrm{bin}}\bigr)^{\mathrm{ab}}
\]
为自然交换化嵌入。则
\[
\ker\rho
\cong
(\ZZ/2\ZZ)^3.
\]
更精确地，
\[
\ker\rho
=
\iota_{\mathrm{ab}}(Z_{\partial}),
\qquad
Z_{\partial}
=
\langle \tau_w:\ w\in X_6^{\mathrm{bdry}}\rangle
\subset Z\!\bigl(G_6^{\mathrm{bin}}\bigr).
\]
因此，任何仅经由 \(X_6^{\mathrm{cyc}}\) 上的 \(B_3/C_3\) 根字典读取的可观测量，都必然对这一规范三维 boundary 奇偶子格失明。
\end{theorem}
\begin{proof}
由定理 \ref{thm:xi-foldbin-gauge-representation-distribution-law-m6-spectrum}，
\[
G_6^{\mathrm{bin}}
\cong
\prod_{w\in X_6}\mathfrak S_{d_6(w)}.
\]
由于 \(m=6\) 时全部纤维大小只可能为 \(2,3,4\)，而
\[
\mathfrak S_2^{\mathrm{ab}}\cong
\mathfrak S_3^{\mathrm{ab}}\cong
\mathfrak S_4^{\mathrm{ab}}
\cong
\ZZ/2\ZZ,
\]
故交换化给出规范识别
\[
\bigl(G_6^{\mathrm{bin}}\bigr)^{\mathrm{ab}}
\cong
\prod_{w\in X_6}\mathfrak S_{d_6(w)}^{\mathrm{ab}}
\cong
(\ZZ/2\ZZ)^{X_6}.
\]
再由定理 \ref{thm:xi-window6-cyclic-kernel-bc3-root-datum}，
\[
X_6=X_6^{\mathrm{cyc}}\sqcup X_6^{\mathrm{bdry}},
\qquad
|X_6^{\mathrm{cyc}}|=18,
\qquad
|X_6^{\mathrm{bdry}}|=3.
\]
因此忘去 boundary 三坐标的投影 \(\rho\) 的核恰为
\[
(\ZZ/2\ZZ)^{X_6^{\mathrm{bdry}}}
\cong
(\ZZ/2\ZZ)^3.
\]

另一方面，定理 \ref{thm:xi-window6-center-exact-splitting-bdry-cyclic} 已经把
\[
Z_{\partial}
=
\langle \tau_w:\ w\in X_6^{\mathrm{bdry}}\rangle
\cong
(\ZZ/2\ZZ)^{X_6^{\mathrm{bdry}}}
\]
识别为 \(Z(G_6^{\mathrm{bin}})\) 的规范直和因子。故
\[
\ker\rho=\iota_{\mathrm{ab}}(Z_{\partial})\cong (\ZZ/2\ZZ)^3.
\]
凡仅通过 \(X_6^{\mathrm{cyc}}\) 的根字典读取的可观测量，都必经由 \(\rho\) 因子化，因而在 \(\ker\rho=\iota_{\mathrm{ab}}(Z_{\partial})\) 上必为常值。证毕。
\end{proof}

\begin{corollary}[审计偶窗最弱退化扇区的奇偶份额恰为 \texorpdfstring{$F_m/F_{m+2}$}{F_m/F_{m+2}}]\label{cor:xi-time-part69c-audited-even-minsector-abel-share}
对审计偶窗
\[
m\in\{6,8,10,12\},
\]
记
\[
\mathcal S_{m,*}:=\mathcal S_{m,d_{\min}(m)},
\qquad
\Lambda_m^{\min}:=(\ZZ/2\ZZ)^{\mathcal S_{m,*}}
\subset
\bigl(G_m^{\mathrm{bin}}\bigr)^{\mathrm{ab}},
\]
其中 \(\Lambda_m^{\min}\) 由定理 \ref{thm:xi-audited-even-abel-center-fibonacci-splitting} 的规范直和分裂给出。则
\[
\operatorname{rank}_{\FF_2}\Lambda_m^{\min}=F_m,
\qquad
\operatorname{rank}_{\FF_2}\!\bigl(G_m^{\mathrm{bin}}\bigr)^{\mathrm{ab}}=F_{m+2},
\]
从而
\[
\frac{\operatorname{rank}_{\FF_2}\Lambda_m^{\min}}
{\operatorname{rank}_{\FF_2}\!\bigl(G_m^{\mathrm{bin}}\bigr)^{\mathrm{ab}}}
=
\frac{F_m}{F_{m+2}}.
\]
具体地，
\[
\frac{8}{21},\qquad
\frac{21}{55},\qquad
\frac{55}{144},\qquad
\frac{144}{377}
\]
分别对应 \(m=6,8,10,12\)。
\end{corollary}
\begin{proof}
由定理 \ref{thm:xi-audited-even-abel-center-fibonacci-splitting}，
\[
\bigl(G_m^{\mathrm{bin}}\bigr)^{\mathrm{ab}}
\cong
(\ZZ/2\ZZ)^{\mathcal S_{m,*}}
\oplus
(\ZZ/2\ZZ)^{X_m\setminus\mathcal S_{m,*}},
\]
并且
\[
|\mathcal S_{m,*}|=F_m.
\]
故
\[
\operatorname{rank}_{\FF_2}\Lambda_m^{\min}=F_m.
\]
另一方面，审计偶窗上由定理 \ref{thm:xi-foldbin-audited-even-center-collapse-abel-full-rank}，
\[
\bigl(G_m^{\mathrm{bin}}\bigr)^{\mathrm{ab}}
\cong
(\ZZ/2\ZZ)^{|X_m|}.
\]
再由稳定类型计数律 \(|X_m|=F_{m+2}\)，得到
\[
\operatorname{rank}_{\FF_2}\!\bigl(G_m^{\mathrm{bin}}\bigr)^{\mathrm{ab}}=F_{m+2}.
\]
于是所示比值公式成立，四个具体分数只需代入 \(m=6,8,10,12\) 即得。证毕。
\end{proof}

\begin{theorem}[审计偶窗最小退化 Fibonacci 外推的首个纯计数截断]\label{thm:xi-time-part69c-minsector-fib-extension-pigeonhole-cutoff}
\leanverified{paper\_xi\_time\_part69c\_minsector\_fib\_extension\_pigeonhole\_cutoff}
把审计恒等式
\[
d_{\min}(m)=F_{m/2}
\qquad
\bigl(m\in\{6,8,10,12\}\bigr)
\]
视为全部偶窗上的延拓候选。则纯计数必要条件已经强制
\[
2^m\ge F_{m+2}F_{m/2}.
\]
该条件在
\[
m=22
\]
时仍成立，但在
\[
m=24
\]
时首次失效；更精确地，
\[
F_{24}F_{11}=46368\cdot 89=4126752<2^{22}=4194304,
\]
而
\[
F_{26}F_{12}=121393\cdot 144=17480592>2^{24}=16777216.
\]
因此，若试图把审计偶窗上的最小退化 Fibonacci 律无条件外推到全部偶窗，则其首个纯鸽巢截断最迟发生在 \(m=24\)。
\end{theorem}
\begin{proof}
若对某个偶数 \(m\) 仍有
\[
d_{\min}(m)=F_{m/2},
\]
则每条纤维都至少含有 \(F_{m/2}\) 个微态。由稳定类型计数律 \(|X_m|=F_{m+2}\) 与恒等式
\[
2^m=\sum_{x\in X_m}d_m(x),
\]
立即得到必要条件
\[
2^m\ge |X_m|\,d_{\min}(m)=F_{m+2}F_{m/2}.
\]
对 \(m=22\) 与 \(m=24\) 直接代入，便得到
\[
F_{24}F_{11}=46368\cdot 89=4126752<2^{22}=4194304,
\]
\[
F_{26}F_{12}=121393\cdot 144=17480592>2^{24}=16777216.
\]
因此这一纯计数必要条件在 \(m=22\) 仍未破裂，而在 \(m=24\) 已经失效；于是任何无条件外推都不可能延续到 \(m=24\) 之后。证毕。
\end{proof}

\noindent
由此，最大非可缩纤维的模 \(6\) 相图、半阶精确暗块、window-\(6\) 根字典对 boundary 奇偶的三维盲核，以及黄金分支在 Fibonacci 分母上的双相输运常数，便不再构成彼此平行的孤立事实：前两者说明同一同余相上可以同时出现大规模谱消失与无折损拓扑主纤维，中者表明可见的根字典天然遗漏一条规范的中心奇偶方向，而纯鸽巢截断又说明最小退化 Fibonacci 律不能脱离审计窗口作无条件延展。因此，该折叠系统的可见图样同时受同余类拓扑、精确暗块与侧向采样异常三类彼此不可约机制共同支配。

\endinput

\paragraph{二原子经验谱、临界容量反演、Bernoulli 残差比与奇素数缺陷尾谱}
在 bin-fold 的两态一致渐近、critical-scale 容量极限曲线、fold-gauge 常数的 Bernoulli 递归剥离以及输出位势的奇素数单位根扭曲四阶展开已经建立之后，还可继续抽取出下列七条后果。它们分别刻画稳定层均匀经验谱的严格两点弱极限、连续容量极限曲线的分布二阶导数、目标成功率下的最优辅助比特预算、由 \((\log C_m)\) 序列恢复 Bernoulli 塔时的 Richardson 残差比、奇素数模 RH 缺陷的可和性与其 \(\ell^\alpha\) 临界指数，以及最小非平凡 prime-twist 谱半径对 \(\lambda\) 的严格二次接近律与最终单调逼近。

\begin{theorem}[bin-fold 纤维谱经自然缩放后的稳定层两点弱极限]\label{thm:xi-time-part70-binfold-uniform-empirical-two-point-law}
\leanverified{paper\_xi\_time\_part70\_binfold\_uniform\_empirical\_two\_point\_law}
定义尺度化多重度
\[
Z_m^{\mathrm{bin}}(w):=\left(\frac{\varphi}{2}\right)^m d_m^{\mathrm{bin}}(w)
\qquad (w\in X_m),
\]
以及稳定层上的均匀经验测度
\[
\nu_m:=\frac{1}{|X_m|}\sum_{w\in X_m}\delta_{Z_m^{\mathrm{bin}}(w)}.
\]
再记
\[
A_0^{(m)}:=\{w\in X_m:\ w_m=0\},
\qquad
A_1^{(m)}:=\{w\in X_m:\ w_m=1\}.
\]
则有弱收敛
\[
\nu_m\xrightarrow{\mathrm w}\nu_\infty
:=
\frac{1}{\varphi}\,\delta_1+\frac{1}{\varphi^2}\,\delta_{\varphi^{-1}}.
\]
等价地，对任意有界连续函数 \(f\)，都有
\[
\lim_{m\to\infty}\frac{1}{|X_m|}\sum_{w\in X_m}f\!\bigl(Z_m^{\mathrm{bin}}(w)\bigr)
=
\frac{1}{\varphi}f(1)+\frac{1}{\varphi^2}f(\varphi^{-1}).
\]
\end{theorem}
\begin{proof}
由定理 \ref{thm:fold-bin-two-state-asymptotic}，存在 \(\varepsilon_m\to 0\) 使得对全部 \(w\in X_m\) 一致有
\[
\bigl|Z_m^{\mathrm{bin}}(w)-\varphi^{-w_m}\bigr|\le \varepsilon_m.
\]
因此对充分大的 \(m\)，所有 \(Z_m^{\mathrm{bin}}(w)\) 都落在某个固定紧区间内。任取有界连续函数 \(f\)，其在该紧区间上一致连续。于是若定义
\[
\eta_{m,0}:=\sup_{w\in A_0^{(m)}}\bigl|f(Z_m^{\mathrm{bin}}(w))-f(1)\bigr|,
\qquad
\eta_{m,1}:=\sup_{w\in A_1^{(m)}}\bigl|f(Z_m^{\mathrm{bin}}(w))-f(\varphi^{-1})\bigr|,
\]
则 \(\eta_{m,0},\eta_{m,1}\to 0\)。从而
\[
\frac{1}{|X_m|}\sum_{w\in X_m}f\!\bigl(Z_m^{\mathrm{bin}}(w)\bigr)
=
\frac{|A_0^{(m)}|}{|X_m|}\bigl(f(1)+o(1)\bigr)
+
\frac{|A_1^{(m)}|}{|X_m|}\bigl(f(\varphi^{-1})+o(1)\bigr).
\]
再由稳定词计数
\[
|A_0^{(m)}|=F_{m+1},
\qquad
|A_1^{(m)}|=F_m,
\qquad
|X_m|=F_{m+2},
\]
以及 Fibonacci 比值极限
\[
\frac{F_{m+1}}{F_{m+2}}\to \frac{1}{\varphi},
\qquad
\frac{F_m}{F_{m+2}}\to \frac{1}{\varphi^2},
\]
便得到所示极限。证毕。
\end{proof}

\begin{theorem}[连续容量极限曲线的分布二阶导数恰为两原子测度]\label{thm:xi-time-part70-continuous-capacity-two-atom-second-derivative}
定义尺度化连续容量
\[
\widehat C_m(s):=
2^{-m}\,\mathcal C_m^{\mathrm{bin,cont}}\!\left(
s\left(\frac{2}{\varphi}\right)^m
\right)
\qquad (s\ge 0),
\]
并记其极限函数
\[
L(s):=\frac{\varphi}{\sqrt5}\min\{1,s\}+\frac{1}{\sqrt5}\min\{\varphi^{-1},s\}.
\]
则在分布意义下有精确恒等式
\[
-D^2L
=
\frac{1}{\sqrt5}\,\delta_{\varphi^{-1}}
+
\frac{\varphi}{\sqrt5}\,\delta_1.
\]
换言之，极限容量曲线只有两个不可消去的折点，其位置分别位于 \(s=\varphi^{-1}\) 与 \(s=1\)，对应原子质量分别为 \(1/\sqrt5\) 与 \(\varphi/\sqrt5\)。
\end{theorem}
\begin{proof}
这正是定理 \ref{thm:xi-foldbin-underresolution-two-atomic-curvature-measure} 中曲率测度公式的直接重写。若将
\[
L(s)=
\begin{cases}
\dfrac{\varphi^2}{\sqrt5}\,s, & 0\le s\le \varphi^{-1},\\[6pt]
\dfrac{\varphi}{\sqrt5}\,s+\dfrac{1}{\varphi\sqrt5}, & \varphi^{-1}\le s\le 1,\\[8pt]
1, & s\ge 1
\end{cases}
\]
逐段求导，则 \(L'\) 在 \(s=\varphi^{-1}\) 处的跃降为
\[
\frac{\varphi^2}{\sqrt5}-\frac{\varphi}{\sqrt5}=\frac{1}{\sqrt5},
\]
在 \(s=1\) 处的跃降为
\[
\frac{\varphi}{\sqrt5}-0=\frac{\varphi}{\sqrt5}.
\]
因此分布导数满足
\[
-D^2L
=
\frac{1}{\sqrt5}\,\delta_{\varphi^{-1}}
+
\frac{\varphi}{\sqrt5}\,\delta_1.
\]
证毕。
\end{proof}

\begin{theorem}[目标成功率下最优辅助比特预算的严格分段阈值]\label{thm:xi-time-part70-target-success-optimal-budget-threshold}
沿用定理 \ref{thm:xi-foldbin-dyadic-capacity-critical-window-limit} 的记号。对任意目标成功率 \(\sigma\in(0,1)\)，定义最小预算
\[
B_m^\star(\sigma):=
\min\bigl\{B\in\mathbb N:\ \mathrm{Succ}_m^{\mathrm{bin}}(B)\ge \sigma\bigr\}.
\]
再定义连续阈值常数
\[
s_\sigma:=
\begin{cases}
\dfrac{\sqrt5}{\varphi^2}\,\sigma,
& 0<\sigma\le \dfrac{\varphi}{\sqrt5},\\[8pt]
\dfrac{\sqrt5\,\sigma-\varphi^{-1}}{\varphi},
& \dfrac{\varphi}{\sqrt5}\le \sigma<1.
\end{cases}
\]
则有
\[
B_m^\star(\sigma)
=
m\log_2\!\left(\frac{2}{\varphi}\right)+\log_2 s_\sigma+O(1).
\]
更一般地，若整数序列 \((B_m)_{m\ge 1}\) 满足
\[
2^{B_m}\left(\frac{\varphi}{2}\right)^m\longrightarrow s\in(0,\infty),
\]
则
\[
\mathrm{Succ}_m^{\mathrm{bin}}(B_m)\longrightarrow L(s),
\]
其中 \(L\) 即定理 \ref{thm:xi-time-part70-continuous-capacity-two-atom-second-derivative} 的极限曲线；因此 \(s_\sigma\) 正是方程 \(L(s)=\sigma\) 的唯一解。
\end{theorem}
\begin{proof}
第二个结论正是定理 \ref{thm:xi-foldbin-dyadic-capacity-critical-window-limit} 与推论 \ref{cor:xi-foldbin-critical-scale-success-curve-inversion} 的内容，其中极限成功曲线即
\[
L(s)=
\begin{cases}
\dfrac{\varphi^2}{\sqrt5}\,s, & 0\le s\le \varphi^{-1},\\[6pt]
\dfrac{\varphi}{\sqrt5}\,s+\dfrac{1}{\varphi\sqrt5}, & \varphi^{-1}\le s\le 1,\\[8pt]
1, & s\ge 1.
\end{cases}
\]
因此 \(s_\sigma=L^{-1}(\sigma)\) 的显式公式恰为上式两段线性方程的解。

现证明 \(B_m^\star(\sigma)\) 的渐近公式。由于 \(\sigma\in(0,1)\)，故 \(s_\sigma\in(0,1)\)。取
\[
0<\varepsilon<\min\{s_\sigma,1-s_\sigma\}.
\]
由 \(L\) 在 \([0,1]\) 上严格递增可知
\[
L(s_\sigma-\varepsilon)<\sigma<L(s_\sigma+\varepsilon).
\]
再由 critical-scale 极限，分别取
\[
\widetilde B_m^-:=\left\lfloor \log_2\!\Bigl((s_\sigma-\varepsilon)\Bigl(\frac{2}{\varphi}\Bigr)^m\Bigr)\right\rfloor,
\qquad
\widetilde B_m^+:=\left\lceil \log_2\!\Bigl((s_\sigma+\varepsilon)\Bigl(\frac{2}{\varphi}\Bigr)^m\Bigr)\right\rceil,
\]
则当 \(m\) 充分大时有
\[
\mathrm{Succ}_m^{\mathrm{bin}}(\widetilde B_m^-)<\sigma,
\qquad
\mathrm{Succ}_m^{\mathrm{bin}}(\widetilde B_m^+)\ge \sigma.
\]
由于 \(\mathrm{Succ}_m^{\mathrm{bin}}(B)\) 关于 \(B\) 单调不减，而 \(B_m^\star(\sigma)\) 按定义是达到成功率阈值的最小整数预算，遂得
\[
\widetilde B_m^-+1\le B_m^\star(\sigma)\le \widetilde B_m^+.
\]
把上式展开后便得到
\[
\log_2\!\Bigl((s_\sigma-\varepsilon)\Bigl(\frac{2}{\varphi}\Bigr)^m\Bigr)
\le
B_m^\star(\sigma)
\le
\log_2\!\Bigl((s_\sigma+\varepsilon)\Bigl(\frac{2}{\varphi}\Bigr)^m\Bigr)+1.
\]
由于 \(\varepsilon\) 固定，端点与
\[
m\log_2\!\left(\frac{2}{\varphi}\right)+\log_2 s_\sigma
\]
之差都是有界量，故结论成立。证毕。
\end{proof}

\begin{theorem}[由 \texorpdfstring{$(\log C_m)$}{(log Cm)} 序列恢复 Bernoulli 塔时的 Richardson 残差比极限]\label{thm:xi-time-part70-logcm-richardson-residual-ratio}
记
\[
q:=\frac{\varphi}{2}\in(0,1),
\qquad
E_m:=\log C_m-\log(2\pi),
\]
\[
a_r:=
\frac{B_{2r}}{r(2r-1)}
\bigl(\varphi^{-1}+\varphi^{2r-3}\bigr)
\qquad (r\ge 1).
\]
对每个固定整数 \(R\ge 0\)，定义 \(R\) 阶 Richardson 残差
\[
\mathcal R_m^{(R)}
:=
E_m-\sum_{r=1}^{R}a_r q^{(2r-1)m}.
\]
则有严格比值极限
\[
\frac{\mathcal R_{m+1}^{(R)}}{\mathcal R_m^{(R)}}
\longrightarrow
q^{2R+1}
=
\left(\frac{\varphi}{2}\right)^{2R+1}.
\]
换言之，一旦减去前 \(R\) 层 Bernoulli 模态，剩余误差的相邻比值必然锁定到奇数次黄金幂 \((\varphi/2)^{2R+1}\)。
\end{theorem}
\begin{proof}
对任意固定 \(R\ge 0\)，将定理 \ref{thm:fold-bin-gauge-constant-stirling-bernoulli-hierarchy} 应用于 \(R+1\) 阶，得到
\[
E_m
=
\sum_{r=1}^{R+1}a_r q^{(2r-1)m}
+O\!\bigl(q^{(2R+3)m}\bigr).
\]
因此
\[
\mathcal R_m^{(R)}
=
a_{R+1}q^{(2R+1)m}
+O\!\bigl(q^{(2R+3)m}\bigr).
\]
又因为偶 Bernoulli 数满足 \(B_{2r}\neq 0\)（\(r\ge 1\)），而 \(\varphi^{-1}+\varphi^{2r-3}>0\)，故
\[
a_{R+1}\neq 0.
\]
于是
\[
\mathcal R_m^{(R)}
=
a_{R+1}q^{(2R+1)m}\bigl(1+O(q^{2m})\bigr),
\]
从而
\[
\frac{\mathcal R_{m+1}^{(R)}}{\mathcal R_m^{(R)}}
=
\frac{a_{R+1}q^{(2R+1)(m+1)}\bigl(1+O(q^{2m})\bigr)}
{a_{R+1}q^{(2R+1)m}\bigl(1+O(q^{2m})\bigr)}
\longrightarrow
q^{2R+1}.
\]
证毕。
\end{proof}

\begin{theorem}[奇素数模 RH 缺陷总质量的可和性]\label{thm:xi-time-part70-odd-prime-rh-defect-summable-mass}
记
\[
\lambda:=\varphi^2,
\qquad
\rho_{p,1}:=\rho\!\bigl(M(e^{2\pi i/p})\bigr)
\qquad
(p\ge 3\ \text{为奇素数}),
\]
并定义扭曲指数与其缺陷
\[
\theta_{p,1}:=\frac{\log \rho_{p,1}}{\log\lambda},
\qquad
\delta_p:=1-\theta_{p,1}.
\]
则有显式展开
\[
\delta_p
=
\frac{4\pi^2(6\sqrt5-5)}{250\log(\varphi^2)}\,\frac{1}{p^2}
+O\!\left(\frac{1}{p^4}\right).
\]
从而
\[
\sum_{\substack{p\ge 3\\ p\ \text{为奇素数}}}\delta_p<\infty,
\qquad
\sum_{\substack{p>P\\ p\ \text{为奇素数}}}\delta_p=O(P^{-1}).
\]
因此，尽管奇素数方向上的平方根门槛失效在频率上呈余有限出现，其缺陷幅度却形成可和尾谱。
\end{theorem}
\begin{proof}
由定理 \ref{thm:xi-output-potential-odd-prime-rh-obstruction-index-to-one}，
\[
\theta_{p,1}
=
1-\frac{6\sqrt5-5}{250\log\lambda}\Bigl(\frac{2\pi}{p}\Bigr)^2
+\frac{7+24\sqrt5}{75000\log\lambda}\Bigl(\frac{2\pi}{p}\Bigr)^4
+O(p^{-6}).
\]
将 \(\lambda=\varphi^2\) 代入并移项，便得到
\[
\delta_p
=
\frac{4\pi^2(6\sqrt5-5)}{250\log(\varphi^2)}\,\frac{1}{p^2}
+O(p^{-4}).
\]
特别地，存在常数 \(c_1,c_2>0\) 与奇素数阈值 \(p_0\)，使得对所有奇素数 \(p\ge p_0\) 都有
\[
\frac{c_1}{p^2}\le \delta_p\le \frac{c_2}{p^2}.
\]
于是
\[
\sum_{\substack{p\ge p_0\\ p\ \text{为奇素数}}}\delta_p
\ll
\sum_{n\ge p_0}\frac{1}{n^2}
<\infty,
\]
从而总和收敛。对尾和，同样有
\[
\sum_{\substack{p>P\\ p\ \text{为奇素数}}}\delta_p
\ll
\sum_{n>P}\frac{1}{n^2}
=
O(P^{-1}).
\]
证毕。
\end{proof}

\begin{corollary}[奇素数 RH 缺陷尾谱的 \texorpdfstring{$\ell^\alpha$}{ell alpha} 临界指数恰为 \texorpdfstring{$1/2$}{1/2}]\label{cor:xi-time-part70-odd-prime-rh-defect-ell-alpha-threshold}
沿用定理 \ref{thm:xi-time-part70-odd-prime-rh-defect-summable-mass} 的记号。则对任意 \(\alpha>0\)，有
\[
(\delta_p)_{p\ \text{为奇素数}}\in \ell^\alpha
\qquad\Longleftrightarrow\qquad
\alpha>\frac12.
\]
\end{corollary}
\begin{proof}
由定理 \ref{thm:xi-time-part70-odd-prime-rh-defect-summable-mass}，存在常数 \(c_1,c_2>0\) 与 \(p_0\) 使得对全部奇素数 \(p\ge p_0\) 都有
\[
c_1p^{-2}\le \delta_p\le c_2p^{-2}.
\]
于是
\[
\delta_p^\alpha\asymp p^{-2\alpha}
\qquad (p\to\infty,\ p\ \text{为奇素数}).
\]
而素数级数满足
\[
\sum_{p}p^{-s}<\infty
\qquad\Longleftrightarrow\qquad
s>1.
\]
因此
\[
\sum_{\substack{p\ge 3\\ p\ \text{为奇素数}}}\delta_p^\alpha<\infty
\qquad\Longleftrightarrow\qquad
\sum_{\substack{p\ge 3\\ p\ \text{为奇素数}}}p^{-2\alpha}<\infty
\qquad\Longleftrightarrow\qquad
2\alpha>1,
\]
即所求结论。证毕。
\end{proof}

\begin{theorem}[奇素数最小非平凡单位根扭曲的谱半径二次接近律与最终严格单调逼近]\label{thm:xi-time-part70-prime-twist-spectral-radius-quadratic-approach}
沿用推论 \ref{cor:real-input-40-collision-twist-fourth} 的记号，并记
\[
\lambda:=\rho(1)=\varphi^2,
\qquad
\rho_p:=\rho\!\left(e^{2\pi i/p}\right)
\qquad
(p\ge 3\ \text{为奇素数}).
\]
则成立：
\begin{enumerate}
  \item 有精确二次接近律
  \[
  \boxed{\;
  p^2\left(1-\frac{\rho_p}{\lambda}\right)
  \longrightarrow
  \frac{6\sqrt5-5}{250}\,(2\pi)^2.\;}
  \]
  \item 存在奇素数阈值 \(p_0\)，使得对任意奇素数 \(p<q\) 且 \(p,q\ge p_0\)，都有
  \[
  \boxed{\;
  \rho_p<\rho_q<\lambda.\;}
  \]
\end{enumerate}
换言之，最小非平凡 prime twist 的谱半径不仅以严格 \(p^{-2}\) 主项从下方逼近 \(\lambda\)，而且在充分大素数区间上按模数严格单调恢复。
\end{theorem}
\begin{proof}
由推论 \ref{cor:real-input-40-collision-twist-fourth}，当 \(t\to 0\) 时有
\[
\log\frac{\rho(t)}{\lambda}
=
-A t^2+B t^4+O(t^6),
\]
其中
\[
A:=\frac{6\sqrt5-5}{250}>0,
\qquad
B:=\frac{7+24\sqrt5}{75000}.
\]
将其指数化，得到
\[
\frac{\rho(t)}{\lambda}
=
\exp\!\bigl(-A t^2+B t^4+O(t^6)\bigr)
=
1-A t^2+O(t^4).
\]
对 \(t=t_p:=2\pi/p\) 代入，即得
\[
1-\frac{\rho_p}{\lambda}
=
A\left(\frac{2\pi}{p}\right)^2+O(p^{-4}),
\]
从而
\[
p^2\left(1-\frac{\rho_p}{\lambda}\right)
\longrightarrow
A(2\pi)^2
=
\frac{6\sqrt5-5}{250}\,(2\pi)^2,
\]
这就证明了第一条。

再令
\[
R(t):=\frac{\rho(t)}{\lambda}.
\]
由上式可得
\[
\log R(t)=-A t^2+B t^4+O(t^6).
\]
对其求导，得到
\[
\frac{R'(t)}{R(t)}
=
-2At+4Bt^3+O(t^5)
=
-2At\bigl(1+O(t^2)\bigr).
\]
因此存在 \(t_0>0\)，使得当 \(0<t\le t_0\) 时，
\[
\frac{R'(t)}{R(t)}<0.
\]
又因 \(R(t)>0\)，可知 \(R'(t)<0\)；故 \(R\) 在 \((0,t_0]\) 上严格递减。另一方面，由
\[
R(t)=1-A t^2+O(t^4)
\]
可知在同一小邻域内 \(R(t)<1\)。

现取奇素数阈值 \(p_0\) 使得
\[
\frac{2\pi}{p_0}\le t_0.
\]
若 \(p<q\) 且 \(p,q\ge p_0\)，则
\[
t_p=\frac{2\pi}{p}>\frac{2\pi}{q}=t_q,
\]
而 \(R\) 在 \((0,t_0]\) 上严格递减，遂得
\[
R(t_p)<R(t_q)<1.
\]
两边同乘以 \(\lambda\)，即
\[
\rho_p<\rho_q<\lambda.
\]
第二条得证。证毕。
\end{proof}

\noindent
上述七条结果表明，二原子塌缩并不只停留在 bin-fold 的输出实验或连续容量曲线上，而会同时外显为稳定层均匀经验谱的严格弱极限、目标成功率反演中的唯一分段阈值、\((\log C_m)\) 的奇次黄金幂残差比，以及奇素数单位根扭曲缺陷的可和尾谱与谱半径恢复律。因而，\(\{\varphi^{-1},1\}\) 这两个黄金比例尺度既支配欠分辨容量几何，也支配 Bernoulli 递归剥离后的剩余模态；而在奇素数单位根扭曲方向上，坏模数的频率与缺陷的振幅不仅被严格分离为“余有限出现”与“总质量可和”两种不同层级，而且最小非平凡扭曲的谱半径还会按严格 \(p^{-2}\) 主项最终单调地从下方逼近 \(\lambda\)。于是，稳定层二点律、资源阈值反演、Richardson 诊断与奇素数 RH 缺陷最终被压缩进同一条黄金尺度主线。

\paragraph{稳定层二态统计极限与规范体积修正的一阶账本}
在稳定层均匀经验谱的两点弱极限、统一层矩函数以及 fold-gauge 常数恢复链已经建立之后，还可把二态一致渐近对均匀层平均量与规范体积修正的支配压缩为下列四条直接后果。它们分别给出缩放退化度的量化两点律、全部缩放矩与幂和首项常数、均匀平均对数退化度的二态展开，以及 \(g_m-\kappa_m+1\) 的精确首项。

\begin{theorem}[均匀稳定类型下缩放退化度的量化两点极限律]\label{thm:xi-time-part70a-uniform-scaled-degeneracy-quantitative}
令
\[
u_m(w):=\frac{1}{|X_m|}\qquad (w\in X_m),
\]
并取 \(W_m\sim u_m\)。定义缩放退化度随机变量
\[
Y_m:=\left(\frac{\varphi}{2}\right)^m d_m^{\mathrm{bin}}(W_m).
\]
则对任意有界 Lipschitz 函数 \(f:\RR\to\RR\)，都有
\[
\mathbb E[f(Y_m)]
=
\frac{F_{m+1}}{F_{m+2}}\,f(1)
+
\frac{F_m}{F_{m+2}}\,f(\varphi^{-1})
+O\!\left(\left(\frac{\varphi}{2}\right)^m\right),
\]
其中隐含常数仅依赖于 \(f\) 的 Lipschitz 常数。特别地，\(Y_m\) 依分布收敛到两点随机变量
\[
Y_\infty=
\begin{cases}
1,&\text{概率 } \varphi^{-1},\\[3pt]
\varphi^{-1},&\text{概率 } \varphi^{-2}.
\end{cases}
\]
\end{theorem}
\begin{proof}
由定理 \ref{thm:fold-bin-two-state-asymptotic}，存在常数 \(C>0\)，使得对充分大的 \(m\) 与全部 \(w\in X_m\) 都一致有
\[
\left|
\left(\frac{\varphi}{2}\right)^m d_m^{\mathrm{bin}}(w)-\varphi^{-w_m}
\right|
\le C\left(\frac{\varphi}{2}\right)^m.
\]
因此
\[
\bigl|Y_m-\varphi^{-W_m(m)}\bigr|
\le C\left(\frac{\varphi}{2}\right)^m
\]
几乎处处成立。若记 \(L_f\) 为 \(f\) 的 Lipschitz 常数，则
\[
\bigl|f(Y_m)-f(\varphi^{-W_m(m)})\bigr|
\le
L_f C\left(\frac{\varphi}{2}\right)^m.
\]
对均匀分布取期望，得到
\[
\mathbb E[f(Y_m)]
=
\mathbb E[f(\varphi^{-W_m(m)})]
+O\!\left(\left(\frac{\varphi}{2}\right)^m\right).
\]
另一方面，
\[
\mathbb E[f(\varphi^{-W_m(m)})]
=
\frac{F_{m+1}}{F_{m+2}}\,f(1)
+
\frac{F_m}{F_{m+2}}\,f(\varphi^{-1}),
\]
因为
\[
\mathbb P(W_m(m)=0)=\frac{F_{m+1}}{F_{m+2}},
\qquad
\mathbb P(W_m(m)=1)=\frac{F_m}{F_{m+2}}.
\]
这就得到所示定量公式。再用
\[
\frac{F_{m+1}}{F_{m+2}}\to\varphi^{-1},
\qquad
\frac{F_m}{F_{m+2}}\to\varphi^{-2},
\]
即可得到所述两点极限分布。证毕。
\end{proof}

\begin{corollary}[统一缩放矩极限与 \texorpdfstring{$S_q^{\mathrm{bin}}(m)$}{Sqbin(m)} 的首项常数]\label{cor:xi-time-part70a-scaled-moments-leading-constant}
对任意实数 \(q\ge 0\)，都有
\[
\lim_{m\to\infty}
\frac{1}{|X_m|}
\sum_{w\in X_m}
\left(
\left(\frac{\varphi}{2}\right)^m d_m^{\mathrm{bin}}(w)
\right)^q
=
\frac{1}{\varphi}+\varphi^{-q-2}.
\]
等价地，若记
\[
S_q^{\mathrm{bin}}(m):=\sum_{w\in X_m}\bigl(d_m^{\mathrm{bin}}(w)\bigr)^q,
\]
则有精确首项渐近
\[
\boxed{\;
S_q^{\mathrm{bin}}(m)
\sim
\frac{\varphi+\varphi^{-q}}{\sqrt5}
\left(\frac{2^q}{\varphi^{q-1}}\right)^m.\;}
\]
\end{corollary}
\begin{proof}
对固定 \(q\ge 0\)，由上一条定理作用于函数 \(f_q(y):=y^q\) 即可。事实上，由定理 \ref{thm:xi-time-part70a-uniform-scaled-degeneracy-quantitative} 的证明，对充分大的 \(m\) 与几乎处处的样本点都有
\[
Y_m\in\left[\frac{\varphi^{-1}}{2},2\right],
\]
故 \(f_q\) 在该固定紧区间上 Lipschitz。于是
\[
\frac{1}{|X_m|}
\sum_{w\in X_m}
\left(
\left(\frac{\varphi}{2}\right)^m d_m^{\mathrm{bin}}(w)
\right)^q
\to
\frac{1}{\varphi}+\varphi^{-q-2}.
\]
将上式改写为
\[
S_q^{\mathrm{bin}}(m)
=
|X_m|
\left(\frac{2}{\varphi}\right)^{mq}
\left(\frac{1}{\varphi}+\varphi^{-q-2}+o(1)\right),
\]
再用 Binet 公式
\[
|X_m|=F_{m+2}\sim \frac{\varphi^{m+2}}{\sqrt5},
\]
便得到
\[
S_q^{\mathrm{bin}}(m)
\sim
\frac{\varphi^{m+2}}{\sqrt5}
\left(\frac{2}{\varphi}\right)^{mq}
\left(\varphi^{-1}+\varphi^{-q-2}\right)
=
\frac{\varphi+\varphi^{-q}}{\sqrt5}
\left(\frac{2^q}{\varphi^{q-1}}\right)^m.
\]
证毕。
\end{proof}

\begin{theorem}[均匀稳定层平均对数退化度的二态展开]\label{thm:xi-time-part70a-uniform-average-logdeg-two-state}
定义均匀稳定层平均对数退化度
\[
\overline{\ell}_m
:=
\frac{1}{|X_m|}\sum_{w\in X_m}\log d_m^{\mathrm{bin}}(w).
\]
则有
\[
\boxed{\;
\overline{\ell}_m
=
m\log\!\Bigl(\frac{2}{\varphi}\Bigr)
-\frac{F_m}{F_{m+2}}\log\varphi
+O\!\left(\left(\frac{\varphi}{2}\right)^m\right).\;}
\]
因而
\[
\boxed{\;
\overline{\ell}_m
=
m\log\!\Bigl(\frac{2}{\varphi}\Bigr)
-\varphi^{-2}\log\varphi
+O\!\left(\left(\frac{\varphi}{2}\right)^m\right).\;}
\]
\end{theorem}
\begin{proof}
由定理 \ref{thm:fold-bin-two-state-asymptotic} 的一致对数展开，
\[
\log d_m^{\mathrm{bin}}(w)
=
m\log\!\Bigl(\frac{2}{\varphi}\Bigr)
-w_m\log\varphi
+O\!\left(\left(\frac{\varphi}{2}\right)^m\right)
\qquad (w\in X_m),
\]
且误差在 \(w\) 上一致。对全部 \(w\in X_m\) 求平均，得到
\[
\overline{\ell}_m
=
m\log\!\Bigl(\frac{2}{\varphi}\Bigr)
-\frac{1}{|X_m|}\sum_{w\in X_m}w_m\,\log\varphi
+O\!\left(\left(\frac{\varphi}{2}\right)^m\right).
\]
而
\[
\sum_{w\in X_m}w_m=|A_1|=F_m,
\qquad
|X_m|=F_{m+2},
\]
故
\[
\frac{1}{|X_m|}\sum_{w\in X_m}w_m=\frac{F_m}{F_{m+2}}.
\]
代回即得第一式。再用 Binet 公式给出的比值极限
\[
\frac{F_m}{F_{m+2}}=\varphi^{-2}+O(\varphi^{-2m}),
\]
便得到第二式。证毕。
\end{proof}
\leanverified{paper\_xi\_time\_part70a\_uniform\_average\_logdeg\_two\_state}

\leanverified{paper\_xi\_time\_part70a\_gauge\_fiber\_gap\_first\_order}
\begin{theorem}[规范体积与纤维信息差的一阶首项]\label{thm:xi-time-part70a-gauge-fiber-gap-first-order}
定义
\[
\mu_m(w):=\frac{d_m^{\mathrm{bin}}(w)}{2^m},
\qquad
\kappa_m:=\sum_{w\in X_m}\mu_m(w)\log d_m^{\mathrm{bin}}(w),
\qquad
g_m:=2^{-m}\log|G_m^{\mathrm{bin}}|,
\]
并沿用推论 \ref{cor:fold-bin-recover-2pi} 的规范常数
\[
C_m
:=
\exp\!\left(
\frac{2^{m+1}}{|X_m|}\bigl(g_m-\kappa_m+1\bigr)-\overline{\ell}_m
\right).
\]
则有
\[
\boxed{\;
\frac{2^{m+1}}{|X_m|}\bigl(g_m-\kappa_m+1\bigr)
=
m\log\!\Bigl(\frac{2}{\varphi}\Bigr)
-\frac{F_m}{F_{m+2}}\log\varphi
+\log(2\pi)
+O\!\left(\left(\frac{\varphi}{2}\right)^m\right).\;}
\]
因此
\[
\boxed{\;
g_m-\kappa_m+1
=
\frac{\varphi^2}{2\sqrt5}
\Bigl(
m\log\!\Bigl(\frac{2}{\varphi}\Bigr)
-\varphi^{-2}\log\varphi
+\log(2\pi)
\Bigr)
\left(\frac{\varphi}{2}\right)^m
+O\!\left(\left(\frac{\varphi}{2}\right)^{2m}\right).\;}
\]
\end{theorem}
\begin{proof}
由 \(C_m\) 的定义可得精确恒等式
\[
\frac{2^{m+1}}{|X_m|}\bigl(g_m-\kappa_m+1\bigr)=\log C_m+\overline{\ell}_m.
\]
另一方面，推论 \ref{cor:fold-bin-recover-2pi} 给出
\[
\log C_m=\log(2\pi)+O\!\left(\left(\frac{\varphi}{2}\right)^m\right),
\]
而定理 \ref{thm:xi-time-part70a-uniform-average-logdeg-two-state} 已证明
\[
\overline{\ell}_m
=
m\log\!\Bigl(\frac{2}{\varphi}\Bigr)
-\frac{F_m}{F_{m+2}}\log\varphi
+O\!\left(\left(\frac{\varphi}{2}\right)^m\right).
\]
相加便得到第一条公式。

再由 Binet 公式，
\[
\frac{|X_m|}{2^{m+1}}
=
\frac{F_{m+2}}{2^{m+1}}
=
\frac{\varphi^2}{2\sqrt5}\left(\frac{\varphi}{2}\right)^m
+O\!\bigl((2\varphi)^{-m}\bigr),
\]
且
\[
\frac{F_m}{F_{m+2}}=\varphi^{-2}+O(\varphi^{-2m}).
\]
用这两式乘第一条公式，得到
\[
g_m-\kappa_m+1
=
\frac{\varphi^2}{2\sqrt5}
\Bigl(
m\log\!\Bigl(\frac{2}{\varphi}\Bigr)
-\varphi^{-2}\log\varphi
+\log(2\pi)
\Bigr)
\left(\frac{\varphi}{2}\right)^m
+O\!\bigl(m(2\varphi)^{-m}\bigr)
+O\!\left(\left(\frac{\varphi}{2}\right)^{2m}\right).
\]
由于
\[
\frac{1}{2\varphi}<\left(\frac{\varphi}{2}\right)^2,
\]
故
\[
m(2\varphi)^{-m}=O\!\left(\left(\frac{\varphi}{2}\right)^{2m}\right).
\]
于是第二条公式成立。证毕。
\end{proof}

\noindent
由此，bin-fold 的两态一致渐近不但控制稳定层均匀经验谱与全部缩放矩，而且把均匀平均的对数退化度与规范体积修正的首项常数同时固定在同一末位二态账本之上。于是，均匀层二点律、幂和首项、平均对数势与规范体积差额最终被收束为同一条可显式求值的黄金比例主线。

\endinput


\paragraph{bin-fold 退化度的全实矩闭式与 Bernoulli 塔的逆矩驱动}
在稳定层均匀经验谱的两点极限、全部非负缩放矩的首项常数以及 fold-gauge 常数的 Stirling--Bernoulli 层级已经建立之后，还可把 bin-fold 退化度的实参数矩谱与 Bernoulli 模态系数进一步压缩为下列两条直接接口。

\begin{theorem}[bin-fold 退化度的全实矩闭式]\label{thm:xi-time-part70aa-binfold-all-real-moments-closed}
对每个 \(m\ge 1\)，记
\[
A_0^{(m)}:=\{w\in X_m:\ w_m=0\},
\qquad
A_1^{(m)}:=\{w\in X_m:\ w_m=1\},
\]
并定义稳定层均匀实矩
\[
M_s^{\mathrm{bin}}(m)
:=
\frac{1}{|X_m|}
\sum_{w\in X_m}\bigl(d_m^{\mathrm{bin}}(w)\bigr)^s
\qquad (s\in\RR).
\]
则对任意固定实数 \(s\)，都有渐近公式
\[
\boxed{\;
M_s^{\mathrm{bin}}(m)
=
\left(\frac{2}{\varphi}\right)^{sm}
\left(
\varphi^{-1}+\varphi^{-s-2}
\right)
+
O_s\!\left(
\left(\frac{2}{\varphi}\right)^{sm}
\left(\frac{\varphi}{2}\right)^m
\right).\;}
\]
等价地，
\[
\boxed{\;
\lim_{m\to\infty}
\left(\frac{\varphi}{2}\right)^{sm}
M_s^{\mathrm{bin}}(m)
=
\varphi^{-1}+\varphi^{-s-2}.\;}
\]
\leanverified{paper\_xi\_time\_part70aa\_binfold\_all\_real\_moments\_closed}
\end{theorem}
\begin{proof}
记
\[
\varepsilon_m:=\left(\frac{\varphi}{2}\right)^m.
\]
由定理 \ref{thm:fold-bin-two-state-asymptotic}，存在常数 \(C>0\)，使对充分大的 \(m\) 与全部 \(w\in X_m\) 都有
\[
\left(\frac{\varphi}{2}\right)^m d_m^{\mathrm{bin}}(w)
=
\begin{cases}
1+O(\varepsilon_m), & w\in A_0^{(m)},\\[4pt]
\varphi^{-1}+O(\varepsilon_m), & w\in A_1^{(m)},
\end{cases}
\]
且误差在 \(w\) 上一致。由于两块主值 \(1\) 与 \(\varphi^{-1}\) 都严格远离 \(0\)，对任意固定 \(s\in\RR\)，幂函数 \(x\mapsto x^s\) 在这两个主值附近都可作一致一阶展开，从而
\[
\bigl(d_m^{\mathrm{bin}}(w)\bigr)^s
=
\left(\frac{2}{\varphi}\right)^{sm}
\begin{cases}
1+O_s(\varepsilon_m), & w\in A_0^{(m)},\\[4pt]
\varphi^{-s}+O_s(\varepsilon_m), & w\in A_1^{(m)},
\end{cases}
\]
其中隐含常数仅依赖于 \(s\)。

对两块分别求和并除以 \(|X_m|=F_{m+2}\)，得到
\[
M_s^{\mathrm{bin}}(m)
=
\left(\frac{2}{\varphi}\right)^{sm}
\left(
\frac{F_{m+1}+\varphi^{-s}F_m}{F_{m+2}}
+
O_s(\varepsilon_m)
\right),
\]
因为
\[
|A_0^{(m)}|=F_{m+1},
\qquad
|A_1^{(m)}|=F_m.
\]
再用 Binet 公式给出的比值渐近
\[
\frac{F_{m+1}}{F_{m+2}}=\varphi^{-1}+O\!\left(\varphi^{-2m}\right),
\qquad
\frac{F_m}{F_{m+2}}=\varphi^{-2}+O\!\left(\varphi^{-2m}\right),
\]
并注意 \(\varphi^{-2m}=O(\varepsilon_m)\)，便得到
\[
M_s^{\mathrm{bin}}(m)
=
\left(\frac{2}{\varphi}\right)^{sm}
\left(
\varphi^{-1}+\varphi^{-s-2}
+
O_s(\varepsilon_m)
\right).
\]
这就给出所示渐近式；第二条极限式只是将主项重新整理后的等价写法。证毕。
\end{proof}

\begin{corollary}[Bernoulli 级数系数的逆矩解释]\label{cor:xi-time-part70aa-bernoulli-coefficient-inverse-moment-interpretation}
对每个整数 \(r\ge 1\)，定义奇数阶逆矩
\[
I_{2r-1}(m)
:=
\frac{1}{|X_m|}
\sum_{w\in X_m}\bigl(d_m^{\mathrm{bin}}(w)\bigr)^{-(2r-1)}.
\]
则有
\[
\boxed{\;
I_{2r-1}(m)
=
\left(\frac{\varphi}{2}\right)^{(2r-1)m}
\left(
\varphi^{-1}+\varphi^{2r-3}
\right)
+
O_r\!\left(
\left(\frac{\varphi}{2}\right)^{2rm}
\right).\;}
\]
从而
\[
\boxed{\;
\lim_{m\to\infty}
\left(\frac{2}{\varphi}\right)^{(2r-1)m}
I_{2r-1}(m)
=
\varphi^{-1}+\varphi^{2r-3}.\;}
\]
另一方面，对任意固定 \(R\ge 1\)，定理 \ref{thm:fold-bin-gauge-constant-stirling-bernoulli-hierarchy} 给出
\[
\log C_m
=
\log(2\pi)
+
\sum_{r=1}^{R}
\frac{B_{2r}}{r(2r-1)}
\bigl(\varphi^{-1}+\varphi^{2r-3}\bigr)
\left(\frac{\varphi}{2}\right)^{(2r-1)m}
+
O\!\left(\left(\frac{\varphi}{2}\right)^{(2R+1)m}\right).
\]
因此，第 \(r\) 级 Bernoulli 模态系数恰可写为
\[
\frac{B_{2r}}{r(2r-1)}
\cdot
\lim_{m\to\infty}
\left(\frac{2}{\varphi}\right)^{(2r-1)m}
I_{2r-1}(m).
\]
\end{corollary}
\begin{proof}
在定理 \ref{thm:xi-time-part70aa-binfold-all-real-moments-closed} 中取
\[
s=-(2r-1)
\]
即可得到
\[
I_{2r-1}(m)=M_{-(2r-1)}^{\mathrm{bin}}(m),
\]
并由
\[
\varphi^{-(-2r+1)-2}=\varphi^{2r-3}
\]
得到所示主项常数。误差项则满足
\[
\left(\frac{2}{\varphi}\right)^{-(2r-1)m}
\left(\frac{\varphi}{2}\right)^m
=
\left(\frac{\varphi}{2}\right)^{2rm}.
\]
最后，把该极限常数与定理 \ref{thm:fold-bin-gauge-constant-stirling-bernoulli-hierarchy} 中对应的第 \(r\) 级系数逐项比较，即得最后一式。证毕。
\end{proof}

\noindent
这表明，fold-gauge 常数序列中的 Bernoulli 塔并不是脱离 bin-fold 退化度谱而附加的解析修正；相反，它正由稳定层退化度的奇数阶逆矩常数逐级驱动。

\endinput


\paragraph{末位二态分层、逆阈值几何与 Bernoulli 塔的 Mellin 编码}
在稳定层均匀经验谱的两点极限、escort 温度族的末位 Bernoulli 塌缩、中心临界容量骨架的双折点闭式以及实参数矩谱与 Bernoulli 模态系数之间的桥接已经建立之后，还可把阈值/采样读出、路径/相位重加权与 Stirling--Bernoulli 层级进一步压缩为同一条末位二态主线。下列结果分别刻画最终严格分层、分位二原子塌缩、归一化逆阈值几何、escort 熵谱的完整常数项，以及容量曲线对 Bernoulli 塔几何权重的 Mellin 编码。

\begin{theorem}[末位相位的最终严格分层]\label{thm:xi-time-part70ab-terminal-bit-eventual-strict-stratification}
\leanverified{paper\_xi\_time\_part70ab\_terminal\_bit\_eventual\_strict\_stratification}
沿用记号
\[
d_m(w):=d_m^{\mathrm{bin}}(w),\qquad
A_i^{(m)}:=\{w\in X_m:\ w_m=i\}\quad(i=0,1),
\qquad
A_m:=\left(\frac{2}{\varphi}\right)^m.
\]
则存在整数 \(m_0\)，使得对一切 \(m\ge m_0\)，都有严格分离
\[
\min_{w\in A_0^{(m)}} d_m(w)
>
\max_{v\in A_1^{(m)}} d_m(v).
\]
因此在充分高分辨率下，纤维多重度的全序被末位 \(w_m\) 强制划分为两层。
\end{theorem}
\begin{proof}
由定理 \ref{thm:fold-bin-two-state-asymptotic}，存在绝对常数 \(C>0\)，使对充分大的 \(m\) 与全部 \(x\in X_m\) 都有
\[
\bigl|d_m(x)-\varphi^{-x_m}A_m\bigr|\le C.
\]
若 \(w\in A_0^{(m)}\) 且 \(v\in A_1^{(m)}\)，则
\[
d_m(w)-d_m(v)
\ge
A_m-C-\bigl(\varphi^{-1}A_m+C\bigr)
=
\varphi^{-2}A_m-2C.
\]
由于 \(A_m\to\infty\)，右端对充分大的 \(m\) 为正，结论成立。
\end{proof}

\begin{corollary}[极值纤维的末位定位与最小尺度失配率]\label{cor:xi-time-part70ab-extremal-fibers-terminal-bit-localization}
\leanverified{paper\_xi\_time\_part70ab\_extremal\_fibers\_terminal\_bit\_localization}
定义
\[
d_{\min}(m):=\min_{w\in X_m}d_m(w),
\qquad
M_m:=\max_{w\in X_m}d_m(w).
\]
则对充分大的 \(m\)，有
\[
\arg\min_{w\in X_m} d_m(w)\subseteq A_1^{(m)},
\qquad
\arg\max_{w\in X_m} d_m(w)\subseteq A_0^{(m)}.
\]
并且
\[
d_{\min}(m)=\varphi^{-1}A_m+O(1),
\qquad
M_m=A_m+O(1),
\qquad
\frac{M_m}{d_{\min}(m)}\longrightarrow \varphi.
\]
此外，沿偶数分辨率 \(m\to\infty\) 有
\[
\frac{d_{\min}(m)}{F_{m/2}}\longrightarrow 0.
\]
\end{corollary}
\begin{proof}
由定理 \ref{thm:xi-time-part70ab-terminal-bit-eventual-strict-stratification}，对充分大的 \(m\)，最小值只能取自 \(A_1^{(m)}\)，最大值只能取自 \(A_0^{(m)}\)。另一方面，由定理 \ref{thm:fold-bin-two-state-asymptotic} 的一致渐近，
\[
d_m(w)=\varphi^{-1}A_m+O(1)\quad(w\in A_1^{(m)}),
\qquad
d_m(w)=A_m+O(1)\quad(w\in A_0^{(m)}),
\]
且误差在各块上统一，故
\[
d_{\min}(m)=\varphi^{-1}A_m+O(1),
\qquad
M_m=A_m+O(1).
\]
于是 \(M_m/d_{\min}(m)\to \varphi\)。

再沿偶数分辨率比较 \(F_{m/2}\)。由 Binet 公式，
\[
F_{m/2}\sim \frac{\varphi^{m/2}}{\sqrt5},
\qquad
d_{\min}(m)\sim \varphi^{-1}\left(\frac{2}{\varphi}\right)^m,
\]
从而
\[
\frac{d_{\min}(m)}{F_{m/2}}
\sim
\sqrt5\,\varphi^{-1}
\left(\frac{2}{\varphi^{3/2}}\right)^m.
\]
由于 \(\varphi^3=2+\sqrt5>4\)，故 \(\varphi^{3/2}>2\)，于是上式趋于 \(0\)。
\end{proof}

\begin{theorem}[中心投影的分位二原子塌缩]\label{thm:xi-time-part70ab-uniform-quantile-two-atom-collapse}
\leanverified{paper\_xi\_time\_part70ab\_uniform\_quantile\_two\_atom\_collapse}
设 \(U_m\) 服从 \(X_m\) 上的均匀分布，并定义尺度化随机变量
\[
Y_m:=\left(\frac{\varphi}{2}\right)^m d_m(U_m).
\]
再定义左连续分位函数
\[
Q_m(\alpha):=\inf\{t\ge 0:\ \mathbf P(Y_m\le t)\ge \alpha\},
\qquad
0<\alpha<1.
\]
则对每个
\[
\alpha\in(0,1)\setminus\{\varphi^{-2}\},
\]
都有极限
\[
Q_m(\alpha)\longrightarrow Q(\alpha),
\qquad
Q(\alpha)=
\begin{cases}
\varphi^{-1}, & 0<\alpha<\varphi^{-2},\\[2mm]
1, & \varphi^{-2}<\alpha<1.
\end{cases}
\]
并且在临界分位 \(\alpha=\varphi^{-2}\) 处出现奇偶振荡：
\[
Q_{2n}(\varphi^{-2})\longrightarrow 1,
\qquad
Q_{2n+1}(\varphi^{-2})\longrightarrow \varphi^{-1}.
\]
\end{theorem}
\begin{proof}
取
\[
0<\varepsilon<\frac{1-\varphi^{-1}}{3}.
\]
由定理 \ref{thm:fold-bin-two-state-asymptotic}，对充分大的 \(m\) 与全部 \(w\in X_m\) 都有
\[
\left|\left(\frac{\varphi}{2}\right)^m d_m(w)-\varphi^{-w_m}\right|
\le
C\left(\frac{\varphi}{2}\right)^m<\varepsilon.
\]
因此当 \(m\) 充分大时，
\[
Y_m\in[\varphi^{-1}-\varepsilon,\varphi^{-1}+\varepsilon]
\quad\text{在 }A_1^{(m)}\text{ 上},
\]
\[
Y_m\in[1-\varepsilon,1+\varepsilon]
\quad\text{在 }A_0^{(m)}\text{ 上}.
\]
故分位值只由块质量
\[
\mathbf P\bigl(Y_m\le \varphi^{-1}+\varepsilon\bigr)
=
\frac{|A_1^{(m)}|}{|X_m|}
=
\frac{F_m}{F_{m+2}}
\]
决定。若 \(\alpha<\varphi^{-2}\) 或 \(\alpha>\varphi^{-2}\)，由
\[
\frac{F_m}{F_{m+2}}\longrightarrow \varphi^{-2}
\]
知对充分大的 \(m\)，比较关系 \(\alpha\lessgtr F_m/F_{m+2}\) 不再改变，于是 \(Q_m(\alpha)\) 分别落在左原子簇与右原子簇中，令 \(\varepsilon\to 0\) 即得所示极限。

最后，由 Binet 公式可直接计算
\[
\frac{F_m}{F_{m+2}}-\varphi^{-2}
=
\frac{(-1)^{m+1}(1-\varphi^{-4})\varphi^{-m}}{\sqrt5\,F_{m+2}},
\]
其符号恰为 \((-1)^{m+1}\)。因此当 \(m\) 为奇数时，
\[
\frac{F_m}{F_{m+2}}>\varphi^{-2},
\]
故 \(Q_m(\varphi^{-2})\) 落在左原子簇；当 \(m\) 为偶数时，
\[
\frac{F_m}{F_{m+2}}<\varphi^{-2},
\]
故 \(Q_m(\varphi^{-2})\) 落在右原子簇。令 \(\varepsilon\to 0\) 即得两条子序列极限。
\end{proof}

\begin{theorem}[escort 温度族的分位相图]\label{thm:xi-time-part70ab-escort-quantile-phase-diagram}
对任意 \(q\ge 0\)，令
\[
W_m^{(q)}\sim \pi_{m,q},
\qquad
Y_m^{(q)}:=\left(\frac{\varphi}{2}\right)^m d_m\bigl(W_m^{(q)}\bigr),
\qquad
\theta(q):=\frac{1}{1+\varphi^{q+1}}.
\]
并定义左连续分位函数
\[
Q_{m,q}(\alpha):=\inf\{t\ge 0:\ \mathbf P(Y_m^{(q)}\le t)\ge \alpha\},
\qquad
0<\alpha<1.
\]
则对每个固定 \(q\ge 0\) 与每个
\[
\alpha\in(0,1)\setminus\{\theta(q)\},
\]
都有极限
\[
Q_{m,q}(\alpha)\longrightarrow Q_q(\alpha),
\qquad
Q_q(\alpha)=
\begin{cases}
\varphi^{-1}, & 0<\alpha<\theta(q),\\[2mm]
1, & \theta(q)<\alpha<1.
\end{cases}
\]
\end{theorem}
\begin{proof}
同样取
\[
0<\varepsilon<\frac{1-\varphi^{-1}}{3}.
\]
由定理 \ref{thm:fold-bin-two-state-asymptotic}，对充分大的 \(m\) 与全部 \(w\in X_m\) 都有
\[
\left|\left(\frac{\varphi}{2}\right)^m d_m(w)-\varphi^{-w_m}\right|<\varepsilon.
\]
因此当 \(m\) 充分大时，
\[
Y_m^{(q)}\in[\varphi^{-1}-\varepsilon,\varphi^{-1}+\varepsilon]
\quad\text{在事件 }W_m^{(q)}(m)=1\text{ 上},
\]
\[
Y_m^{(q)}\in[1-\varepsilon,1+\varepsilon]
\quad\text{在事件 }W_m^{(q)}(m)=0\text{ 上}.
\]
另一方面，由命题 \ref{prop:fold-bin-escort-lastbit}，
\[
\mathbf P\bigl(W_m^{(q)}(m)=1\bigr)=\theta(q)+O\!\left(\left(\frac{\varphi}{2}\right)^m\right).
\]
于是当 \(\alpha<\theta(q)\) 时，充分大的 \(m\) 上有 \(\alpha<\mathbf P(W_m^{(q)}(m)=1)\)，故 \(Q_{m,q}(\alpha)\) 落在左原子簇；当 \(\alpha>\theta(q)\) 时，则 \(Q_{m,q}(\alpha)\) 落在右原子簇。令 \(\varepsilon\to 0\) 即得结论。
\end{proof}

\begin{theorem}[中心容量逆问题的闭式解]\label{thm:xi-time-part70ab-center-capacity-inverse-threshold-closed}
\leanverified{paper\_xi\_time\_part70ab\_center\_capacity\_inverse\_threshold\_closed}
定义归一化连续容量
\[
\Psi_m(\tau):=
2^{-m}\,\mathcal C_m^{\mathrm{bin,cont}}(\tau A_m),
\qquad
\tau\ge 0,
\]
并定义达到水平 \(\rho\) 的最小归一化阈值
\[
\tau_\ast(\rho):=
\inf\left\{\tau\ge 0:\ \liminf_{m\to\infty}\Psi_m(\tau)\ge \rho\right\},
\qquad
0\le \rho\le 1.
\]
则有严格闭式
\[
\tau_\ast(\rho)=
\begin{cases}
\displaystyle \frac{\sqrt5}{\varphi^2}\,\rho,
& 0\le \rho\le \dfrac{\varphi}{\sqrt5},\\[3mm]
\displaystyle \frac{\sqrt5\,\rho-\varphi^{-1}}{\varphi},
& \dfrac{\varphi}{\sqrt5}\le \rho\le 1.
\end{cases}
\]
\end{theorem}
\begin{proof}
由定理 \ref{thm:xi-time-part70b-center-capacity-two-kink-skeleton}，
\[
\Psi_m(\tau)\longrightarrow \Psi(\tau):=
\frac{\varphi}{\sqrt5}\min\{1,\tau\}
+
\frac{1}{\sqrt5}\min\{\varphi^{-1},\tau\}.
\]
因此
\[
\tau_\ast(\rho)=\inf\{\tau\ge 0:\ \Psi(\tau)\ge \rho\}.
\]
而
\[
\Psi(\tau)=
\begin{cases}
\dfrac{\varphi^2}{\sqrt5}\tau, & 0\le \tau\le \varphi^{-1},\\[6pt]
\dfrac{\varphi}{\sqrt5}\tau+\dfrac{1}{\varphi\sqrt5}, & \varphi^{-1}\le \tau\le 1,\\[8pt]
1, & \tau\ge 1,
\end{cases}
\]
故广义逆由两段线性方程直接求得，即为所示公式。
\end{proof}

\begin{theorem}[escort 响应的逆温度阈值公式]\label{thm:xi-time-part70ab-escort-response-inverse-threshold-closed}
对固定 \(q\ge 0\)，定义 escort 归一化响应
\[
\Psi_{m,q}^{\mathrm{esc}}(\tau):=
\mathbb E_{\pi_{m,q}}
\left[
\min\!\left(1,\frac{\tau A_m}{d_m(W_m^{(q)})}\right)
\right],
\qquad
\tau\ge 0.
\]
再定义达到水平 \(\rho\) 的最小归一化阈值
\[
\tau_{q,\ast}(\rho):=
\inf\left\{\tau\ge 0:\ \liminf_{m\to\infty}\Psi_{m,q}^{\mathrm{esc}}(\tau)\ge \rho\right\},
\qquad
0\le \rho\le 1.
\]
记
\[
\rho_c(q):=\theta(q)+\varphi^{-1}\bigl(1-\theta(q)\bigr).
\]
则有闭式
\[
\tau_{q,\ast}(\rho)=
\begin{cases}
\displaystyle \frac{\rho}{1+(\varphi-1)\theta(q)},
& 0\le \rho\le \rho_c(q),\\[3mm]
\displaystyle \frac{\rho-\theta(q)}{1-\theta(q)},
& \rho_c(q)\le \rho\le 1.
\end{cases}
\]
\end{theorem}
\begin{proof}
由定理 \ref{thm:xi-time-part70b-escort-two-kink-response}，
\[
\Psi_{m,q}^{\mathrm{esc}}(\tau)\longrightarrow
\Psi_q^{\mathrm{esc}}(\tau):=
\bigl(1-\theta(q)\bigr)\min\{1,\tau\}
+
\theta(q)\min\{1,\varphi\tau\}.
\]
因此
\[
\tau_{q,\ast}(\rho)=\inf\{\tau\ge 0:\ \Psi_q^{\mathrm{esc}}(\tau)\ge \rho\}.
\]
把 \(\Psi_q^{\mathrm{esc}}\) 在 \([0,\varphi^{-1}]\) 与 \([\varphi^{-1},1]\) 上分段展开，得到
\[
\Psi_q^{\mathrm{esc}}(\tau)=
\begin{cases}
\bigl(1+(\varphi-1)\theta(q)\bigr)\tau,
& 0\le \tau\le \varphi^{-1},\\[6pt]
\theta(q)+\bigl(1-\theta(q)\bigr)\tau,
& \varphi^{-1}\le \tau\le 1,\\[6pt]
1, & \tau\ge 1.
\end{cases}
\]
其连接点函数值为
\[
\Psi_q^{\mathrm{esc}}(\varphi^{-1})
=
\theta(q)+\varphi^{-1}\bigl(1-\theta(q)\bigr)
=
\rho_c(q).
\]
故广义逆由上述两段线性函数分别求反即可。
\end{proof}

\begin{theorem}[escort 熵谱的完整常数项]\label{thm:xi-time-part70ab-escort-renyi-entropy-full-constant}
对固定 \(q\ge 0\)、\(a>0\) 且 \(a\neq 1\)，定义 escort Rényi 熵
\[
H_a(\pi_{m,q})
:=
\frac{1}{1-a}\log\sum_{w\in X_m}\pi_{m,q}(w)^a.
\]
则有渐近展开
\[
H_a(\pi_{m,q})
=
m\log\varphi
-\frac12\log 5
+
\frac{1}{1-a}
\log\!\Bigl(
\varphi^{1-a}\bigl(1-\theta(q)\bigr)^a+\theta(q)^a
\Bigr)
+o(1).
\]
当 \(a\to 1\) 时，得到 Shannon 熵常数项
\[
H(\pi_{m,q})
=
m\log\varphi
-\frac12\log 5
+
\bigl(1-\theta(q)\bigr)\log\varphi
+
h\bigl(\theta(q)\bigr)
+o(1),
\]
其中
\[
h(t):=-t\log t-(1-t)\log(1-t).
\]
\end{theorem}
\begin{proof}
记
\[
S_t(m):=\sum_{w\in X_m} d_m(w)^t.
\]
由定理 \ref{thm:xi-time-part70aa-binfold-all-real-moments-closed} 与 Binet 公式，对每个固定实数 \(t\) 都有
\[
S_t(m)
=
\frac{\varphi+\varphi^{-t}}{\sqrt5}
\bigl(2^t\varphi^{1-t}\bigr)^m
\bigl(1+o(1)\bigr).
\]
因此
\[
\sum_{w\in X_m}\pi_{m,q}(w)^a
=
\frac{S_{aq}(m)}{S_q(m)^a}
=
5^{(a-1)/2}\,
\varphi^{(1-a)m}\,
\frac{\varphi+\varphi^{-aq}}{(\varphi+\varphi^{-q})^a}
\bigl(1+o(1)\bigr).
\]
取对数并除以 \(1-a\)，得到
\[
H_a(\pi_{m,q})
=
m\log\varphi
-\frac12\log 5
+
\frac{1}{1-a}
\log\!\frac{\varphi+\varphi^{-aq}}{(\varphi+\varphi^{-q})^a}
+o(1).
\]
再由
\[
\theta(q)=\frac{\varphi^{-q}}{\varphi+\varphi^{-q}},
\qquad
1-\theta(q)=\frac{\varphi}{\varphi+\varphi^{-q}},
\]
可将最后一项重写为
\[
\frac{\varphi+\varphi^{-aq}}{(\varphi+\varphi^{-q})^a}
=
\varphi^{1-a}\bigl(1-\theta(q)\bigr)^a+\theta(q)^a.
\]
这就给出第一条公式。

当 \(a\to 1\) 时，对
\[
F_q(a):=
\log\!\Bigl(
\varphi^{1-a}\bigl(1-\theta(q)\bigr)^a+\theta(q)^a
\Bigr)
\]
应用 l'Hospital 法则，得到
\[
\lim_{a\to 1}\frac{F_q(a)}{1-a}
=
-F_q'(1)
=
\bigl(1-\theta(q)\bigr)\log\varphi+h\bigl(\theta(q)\bigr).
\]
代回即得 Shannon 常数项。
\end{proof}

\begin{theorem}[容量极限曲线的负矩 Mellin 反演]\label{thm:xi-time-part70ab-capacity-limit-negative-moment-mellin}
定义中心容量极限
\[
m_\ast(\tau):=
\lim_{m\to\infty}
\frac{\sqrt5}{\varphi^2}\,
2^{-m}\,\mathcal C_m^{\mathrm{bin,cont}}(\tau A_m),
\qquad
\tau\ge 0.
\]
则对任意实数 \(s>0\)，成立严格恒等式
\[
s(s+1)\int_0^\infty \tau^{-s-2}\bigl(\tau-m_\ast(\tau)\bigr)\,\dd\tau
=
\varphi^{-1}+\varphi^{s-2}.
\]
\end{theorem}
\begin{proof}
由定理 \ref{thm:xi-time-part70b-center-capacity-two-kink-skeleton}，
\[
m_\ast(\tau)=
\varphi^{-1}\min\{1,\tau\}
+
\varphi^{-2}\min\{\varphi^{-1},\tau\}.
\]
因而若定义双原子概率测度
\[
\mu_\ast:=\varphi^{-2}\delta_{\varphi^{-1}}+\varphi^{-1}\delta_1,
\]
则有
\[
m_\ast(\tau)=\int \min\{y,\tau\}\,\dd\mu_\ast(y).
\]
对任意固定 \(y>0\) 与 \(s>0\)，直接计算得
\[
\int_0^\infty \tau^{-s-2}\bigl(\tau-\min\{y,\tau\}\bigr)\,\dd\tau
=
\int_y^\infty \tau^{-s-2}(\tau-y)\,\dd\tau
=
\frac{y^{-s}}{s(s+1)}.
\]
两边乘以 \(s(s+1)\) 并对 \(\mu_\ast\) 积分，即得
\[
s(s+1)\int_0^\infty \tau^{-s-2}\bigl(\tau-m_\ast(\tau)\bigr)\,\dd\tau
=
\int y^{-s}\,\dd\mu_\ast(y)
=
\varphi^{-1}+\varphi^{-2}\varphi^s
=
\varphi^{-1}+\varphi^{s-2}.
\]
证毕。
\end{proof}

\begin{corollary}[Stirling--Bernoulli 几何系数的容量解释]\label{cor:xi-time-part70ab-bernoulli-hierarchy-capacity-interpretation}
\leanverified{paper\_xi\_time\_part70ab\_bernoulli\_hierarchy\_capacity\_interpretation}
对任意整数 \(r\ge 1\)，都有
\[
\varphi^{-1}+\varphi^{2r-3}
=
(2r-1)(2r)\int_0^\infty \tau^{-2r-1}\bigl(\tau-m_\ast(\tau)\bigr)\,\dd\tau.
\]
从而 Stirling--Bernoulli 层级中第 \(r\) 级系数可写为
\[
\frac{B_{2r}}{r(2r-1)}
\bigl(\varphi^{-1}+\varphi^{2r-3}\bigr)
=
\frac{2B_{2r}}{r}
\int_0^\infty \tau^{-2r-1}\bigl(\tau-m_\ast(\tau)\bigr)\,\dd\tau.
\]
特别地，对任意固定 \(R\ge 1\)，定理 \ref{thm:fold-bin-gauge-constant-stirling-bernoulli-hierarchy} 的展开可重写为
\[
\log C_m
=
\log(2\pi)
+
\sum_{r=1}^{R}
\frac{2B_{2r}}{r}
\left(
\int_0^\infty \tau^{-2r-1}\bigl(\tau-m_\ast(\tau)\bigr)\,\dd\tau
\right)
\left(\frac{\varphi}{2}\right)^{(2r-1)m}
+
O\!\left(\left(\frac{\varphi}{2}\right)^{(2R+1)m}\right).
\]
\end{corollary}
\begin{proof}
在定理 \ref{thm:xi-time-part70ab-capacity-limit-negative-moment-mellin} 中取 \(s=2r-1\) 即得第一式。再把该恒等式代入定理 \ref{thm:fold-bin-gauge-constant-stirling-bernoulli-hierarchy} 的系数表达，便得到后两式。
\end{proof}

\noindent
由此，末位 \(0/1\) 的二态骨架不只控制纤维多重度的最终排序与均匀/escort 采样下的分位相图，还把中心容量曲线的广义逆、escort 温度响应的广义逆以及 \(\log C_m\) 的 Stirling--Bernoulli 几何权重全部还原为同一组双原子数据。于是，阈值/采样轴、路径/相位轴与 Bernoulli 塔的解析层级便在同一条末位二态主线上闭合。

\endinput

\paragraph{全息地址前缀商的精确熵律与最小判别比特}
在 H.160--H.161 的 faithful \(2^L\)-进地址表示、精确柱分离与圆柱--球识别已经固定之后，有限深度地址商的可见信息量还可进一步压缩为一条完全闭式的计数律。

\begin{theorem}[\texorpdfstring{$2^L$}{2^L}-进全息地址前缀商的精确计数与最小判别比特]\label{thm:xi-time-part70ac-hologram-prefix-quotient-entropy-bits}
\leanverified{paper\_xi\_time\_part70ac\_hologram\_prefix\_quotient\_entropy\_bits}
设 \(E\) 为有限事件集，
\[
\mathrm{code}:E\hookrightarrow \NN,
\qquad
M:=\max_{e\in E}\mathrm{code}(e),
\qquad
B:=2^L>M,
\]
并记
\[
T:=T_2(E_{\min}),
\qquad
v\in T\setminus 2T,
\]
其中
\[
X(\mathbf e):=\sum_{t\ge 1}\mathrm{code}(e_t)\,B^{t-1}v
\qquad
\bigl(\mathbf e=(e_t)_{t\ge 1}\in E^{\NN}\bigr).
\]
并定义
\[
\mathcal X_{E,v}:=X(E^{\NN}).
\]
对每个 \(n\ge 1\)，定义第 \(n\) 级地址商集
\[
\mathcal A_n
:=
\left\{
X(\mathbf e)\bmod B^nT:\ \mathbf e\in E^{\NN}
\right\}
\subseteq T/B^nT.
\]
再对每个长度 \(n\) 前缀
\[
\mathbf a=(a_1,\dots,a_n)\in E^n
\]
记
\[
[\mathbf a]
:=
\left\{
\mathbf e\in E^{\NN}:\ \mathrm{pref}_n(\mathbf e)=\mathbf a
\right\},
\qquad
x_{\mathbf a}
:=
\sum_{t=1}^{n}\mathrm{code}(a_t)\,B^{t-1}v.
\]
则成立：
\begin{enumerate}
  \item 地址商集满足精确计数公式
  \[
  |\mathcal A_n|=|E|^n.
  \]
  \item 若定义深度 \(n\) 的 clopen 地址柱代数
  \[
  \mathcal B_n
  :=
  \left\{
  \bigcup_{\mathbf a\in U}X([\mathbf a]):\ U\subseteq E^n
  \right\},
  \]
  则
  \[
  |\mathcal B_n|=2^{|E|^n}.
  \]
  \item 若 \(A\) 为有限集，且读出映射
  \[
  R:\mathcal X_{E,v}\longrightarrow A
  \]
  满足对任意不同前缀 \(\mathbf a\neq \mathbf b\in E^n\) 及任意
  \[
  x\in X([\mathbf a]),
  \qquad
  y\in X([\mathbf b]),
  \]
  都有 \(R(x)\neq R(y)\)，则
  \[
  |A|\ge |E|^n.
  \]
  因而区分全部长度 \(n\) 前缀所需的最小判别比特数恰为
  \[
  \left\lceil n\log_2|E|\right\rceil.
  \]
\end{enumerate}
\end{theorem}
\begin{proof}
固定 \(n\ge 1\)。定义映射
\[
\Phi_n:E^n\longrightarrow \mathcal A_n,
\qquad
\mathbf a\longmapsto x_{\mathbf a}+B^nT.
\]
先证满射。任取 \(\mathbf e=(e_t)_{t\ge 1}\in E^{\NN}\)，并记
\[
\mathbf a:=\mathrm{pref}_n(\mathbf e)=(e_1,\dots,e_n).
\]
则
\[
X(\mathbf e)-x_{\mathbf a}
=
\sum_{t>n}\mathrm{code}(e_t)\,B^{t-1}v
\in B^nT,
\]
故
\[
X(\mathbf e)\bmod B^nT=\Phi_n(\mathbf a).
\]
因此 \(\Phi_n\) 满射。

再证单射。若
\[
\Phi_n(\mathbf a)=\Phi_n(\mathbf b),
\]
则
\[
x_{\mathbf a}-x_{\mathbf b}\in B^nT.
\]
任取
\[
\mathbf e\in[\mathbf a],
\qquad
\mathbf f\in[\mathbf b].
\]
由结论 \ref{con:xi-time-part62dgb-hologram-cylinder-ball-isometry}，
\[
X(\mathbf e)\equiv x_{\mathbf a}\pmod{B^nT},
\qquad
X(\mathbf f)\equiv x_{\mathbf b}\pmod{B^nT},
\]
从而
\[
X(\mathbf e)\equiv X(\mathbf f)\pmod{B^nT}.
\]
再由同一结论中的同余--前缀等价式，只能有
\[
\mathrm{pref}_n(\mathbf e)=\mathrm{pref}_n(\mathbf f),
\]
亦即 \(\mathbf a=\mathbf b\)。故 \(\Phi_n\) 为双射，于是
\[
|\mathcal A_n|=|E^n|=|E|^n,
\]
第 \(1\) 条得证。

对第 \(2\) 条，由定理 \ref{thm:xi-time-part62d-hologram-image-cantor-homeomorphism}，映射
\[
X:E^{\NN}\longrightarrow \mathcal X_{E,v}
\]
是拓扑同胚。长度 \(n\) 的 cylinder 集族
\[
\{[\mathbf a]:\ \mathbf a\in E^n\}
\]
在 \(E^{\NN}\) 中构成一个由 \(|E|^n\) 个原子组成的 clopen 分割，故其像
\[
\{X([\mathbf a]):\ \mathbf a\in E^n\}
\]
在 \(\mathcal X_{E,v}\) 中同样构成一个由 \(|E|^n\) 个原子组成的 clopen 分割。于是由这些原子生成的 Boolean 代数正是全部原子并的集合族，故
\[
|\mathcal B_n|=2^{|E|^n}.
\]

再证第 \(3\) 条。对每个 \(\mathbf a\in E^n\)，取一点
\[
x_{\mathbf a}^{\ast}\in X([\mathbf a]),
\]
这是可能的，因为每个 cylinder 都非空。若 \(\mathbf a\neq \mathbf b\)，则按假设
\[
R(x_{\mathbf a}^{\ast})\neq R(x_{\mathbf b}^{\ast}).
\]
因此映射
\[
E^n\longrightarrow A,
\qquad
\mathbf a\longmapsto R(x_{\mathbf a}^{\ast})
\]
是单射，故
\[
|A|\ge |E|^n.
\]
这给出输出值个数的下界。

反向上，由定理 \ref{thm:xi-time-part62d-hologram-image-cantor-homeomorphism}，\(X\) 为双射，故前缀读出
\[
R_n:\mathcal X_{E,v}\longrightarrow E^n,
\qquad
R_n(x):=\mathrm{pref}_n(X^{-1}(x))
\]
良定义，并且恰取 \(|E|^n\) 个值。于是最小所需输出值个数正好等于 \(|E|^n\)。将有限输出值编码为二进制字，最小判别比特数遂为
\[
\left\lceil \log_2|E|^n\right\rceil
=
\left\lceil n\log_2|E|\right\rceil.
\]
证毕。
\end{proof}

\noindent
因此，深度 \(n\) 的地址可见层并非只在数量级上与 \(|E|^n\) 同阶，而是恰由 \(|E|^n\) 个前缀原子严格组成；相应有限值读出的最小复杂度也被精确锁定为 \(\left\lceil n\log_2|E|\right\rceil\) 比特。于是，时间推进所带来的地址精度提升被压缩为一条完全刚性的有限计数律。

\endinput

\paragraph{bin-fold 输出分布的全阶信息几何常数层闭式}
在输出分布的 Rényi 熵率塌缩、隐藏对数多重度的二原子 Laplace 极限、规范体积密度的两尺度展开以及碰撞矩的两态主项都已建立之后，还可把 bin-fold 推前分布 \(\mu_m\) 的常数层信息几何进一步压缩为下列四条闭式结论。

\begin{theorem}[输出分布的全阶 Rényi 熵具有精确常数漂移]\label{thm:xi-time-part70ad-binfold-renyi-entropy-constant-drift}
\leanverified{paper\_xi\_time\_part70ad\_binfold\_renyi\_entropy\_constant\_drift}
记
\[
\mu_m(w):=\frac{d_m^{\mathrm{bin}}(w)}{2^m}
\qquad (w\in X_m),
\]
并对 \(q>0\)、\(q\neq 1\) 定义
\[
H_q(\mu_m):=\frac{1}{1-q}\log\sum_{w\in X_m}\mu_m(w)^q.
\]
则对每个固定 \(q>0\)、\(q\neq 1\)，都有精确渐近
\[
\boxed{\;
H_q(\mu_m)
=
m\log\varphi
+
\frac{1}{1-q}\log\!\Bigl(\frac{\varphi+\varphi^{-q}}{\sqrt5}\Bigr)
+
O_q\!\Bigl(\Bigl(\frac{\varphi}{2}\Bigr)^m\Bigr).\;}
\]
特别地，
\[
\boxed{\;
\lim_{m\to\infty}\frac{1}{m}H_q(\mu_m)=\log\varphi
\qquad (q>0,\ q\neq 1).\;}
\]
\end{theorem}
\begin{proof}
记
\[
A_0^{(m)}:=\{w\in X_m:\ w_m=0\},
\qquad
A_1^{(m)}:=\{w\in X_m:\ w_m=1\}.
\]
由定理 \ref{thm:fold-bin-two-state-asymptotic} 的一致主项，对每个固定 \(q>0\) 都存在常数 \(C_q>0\)，使得对充分大的 \(m\) 与全部 \(w\in X_m\) 都有
\[
\bigl(d_m^{\mathrm{bin}}(w)\bigr)^q
=
\left(\frac{2}{\varphi}\right)^{qm}
\begin{cases}
1+O_q\!\bigl((\varphi/2)^m\bigr), & w\in A_0^{(m)},\\[4pt]
\varphi^{-q}+O_q\!\bigl((\varphi/2)^m\bigr), & w\in A_1^{(m)}.
\end{cases}
\]
故
\[
S_q^{\mathrm{bin}}(m):=\sum_{w\in X_m}\bigl(d_m^{\mathrm{bin}}(w)\bigr)^q
=
\left(\frac{2}{\varphi}\right)^{qm}
\Bigl(
F_{m+1}+\varphi^{-q}F_m
\Bigr)
\Bigl(1+O_q\!\bigl((\varphi/2)^m\bigr)\Bigr),
\]
因为 \(|A_0^{(m)}|=F_{m+1}\)、\(|A_1^{(m)}|=F_m\)。于是
\[
\sum_{w\in X_m}\mu_m(w)^q
=
2^{-qm}S_q^{\mathrm{bin}}(m)
=
\varphi^{-qm}
\Bigl(
F_{m+1}+\varphi^{-q}F_m
\Bigr)
\Bigl(1+O_q\!\bigl((\varphi/2)^m\bigr)\Bigr).
\]
再由 Binet 公式，
\[
F_{m+1}+\varphi^{-q}F_m
=
\frac{\varphi^m}{\sqrt5}\bigl(\varphi+\varphi^{-q}\bigr)+O(\varphi^{-m}),
\]
而 \(O(\varphi^{-m})=O((\varphi/2)^m\varphi^m)\)，故
\[
\sum_{w\in X_m}\mu_m(w)^q
=
\varphi^{(1-q)m}
\frac{\varphi+\varphi^{-q}}{\sqrt5}
\Bigl(1+O_q\!\bigl((\varphi/2)^m\bigr)\Bigr).
\]
右端主因子严格为正，取对数并除以 \(1-q\) 即得
\[
H_q(\mu_m)
=
m\log\varphi
+
\frac{1}{1-q}\log\!\Bigl(\frac{\varphi+\varphi^{-q}}{\sqrt5}\Bigr)
+
O_q\!\Bigl(\Bigl(\frac{\varphi}{2}\Bigr)^m\Bigr).
\]
再除以 \(m\) 并令 \(m\to\infty\)，便得到熵率极限。证毕。
\end{proof}

\leanverified{paper\_xi\_time\_part70ad\_hidden\_logmultiplicity\_exponential\_moments}
\begin{theorem}[中心化对数多重度的全指数矩极限与全部累积量闭式]\label{thm:xi-time-part70ad-hidden-logmultiplicity-exponential-moments}
定义
\[
Y_m(w):=\log d_m^{\mathrm{bin}}(w)-m\log\!\Bigl(\frac{2}{\varphi}\Bigr),
\qquad
W_m\sim \mu_m.
\]
则对任意固定复数 \(t\in\CC\)，都有
\[
\boxed{\;
\EE_{\mu_m}\!\bigl[e^{tY_m(W_m)}\bigr]
=
\frac{\varphi+\varphi^{-(t+1)}}{\sqrt5}
+
O_t\!\Bigl(\Bigl(\frac{\varphi}{2}\Bigr)^m\Bigr).\;}
\]
因此 \(Y_m(W_m)\) 在分布上并且在全部指数矩意义下收敛到两点随机变量
\[
Y_\infty=
\begin{cases}
0,& \text{概率 } \dfrac{\varphi}{\sqrt5},\\[6pt]
-\log\varphi,& \text{概率 } \dfrac{1}{\varphi\sqrt5}.
\end{cases}
\]
若记
\[
p:=\frac{1}{\varphi\sqrt5},
\qquad
\ell:=\log\varphi,
\]
则其极限累积量生成函数为
\[
K_\infty(t)=\log\bigl((1-p)+pe^{-t\ell}\bigr),
\]
并且前四阶极限累积量满足
\[
\boxed{\;
\kappa_1^\infty=-p\ell,\qquad
\kappa_2^\infty=p(1-p)\ell^2,\qquad
\kappa_3^\infty=-p(1-p)(1-2p)\ell^3,\qquad
\kappa_4^\infty=p(1-p)(1-6p+6p^2)\ell^4.\;}
\]
\end{theorem}
\begin{proof}
记 \(B_m:=W_m(m)\in\{0,1\}\)。由定理 \ref{thm:fold-bin-two-state-asymptotic} 的一致对数展开，
\[
Y_m(w)=-w_m\log\varphi+O\!\Bigl(\Bigl(\frac{\varphi}{2}\Bigr)^m\Bigr)
\qquad (w\in X_m),
\]
于是对任意固定 \(t\in\CC\)，都有一致展开
\[
e^{tY_m(w)}
=
\varphi^{-t w_m}
\Bigl(1+O_t\!\Bigl(\Bigl(\frac{\varphi}{2}\Bigr)^m\Bigr)\Bigr).
\]
对 \(\mu_m\) 取期望，并利用命题 \ref{prop:fold-bin-escort-lastbit} 在 \(q=1\) 处的末位质量公式
\[
\mu_m(B_m=0)=\frac{\varphi}{\sqrt5}+O\!\Bigl(\Bigl(\frac{\varphi}{2}\Bigr)^m\Bigr),
\qquad
\mu_m(B_m=1)=\frac{1}{\varphi\sqrt5}+O\!\Bigl(\Bigl(\frac{\varphi}{2}\Bigr)^m\Bigr),
\]
便得到
\[
\EE_{\mu_m}\!\bigl[e^{tY_m(W_m)}\bigr]
=
\frac{\varphi}{\sqrt5}
+
\frac{1}{\varphi\sqrt5}\varphi^{-t}
+
O_t\!\Bigl(\Bigl(\frac{\varphi}{2}\Bigr)^m\Bigr)
=
\frac{\varphi+\varphi^{-(t+1)}}{\sqrt5}
+
O_t\!\Bigl(\Bigl(\frac{\varphi}{2}\Bigr)^m\Bigr).
\]
这就得到指数矩极限。

由于 \(Y_m\) 一致有界，上式同时给出所有指数矩收敛，并推出 \(Y_m(W_m)\Rightarrow Y_\infty\)，其中
\[
\PP(Y_\infty=0)=1-p=\frac{\varphi}{\sqrt5},
\qquad
\PP(Y_\infty=-\ell)=p=\frac{1}{\varphi\sqrt5}.
\]
故极限累积量生成函数即为
\[
K_\infty(t)=\log\EE[e^{tY_\infty}]
=
\log\bigl((1-p)+pe^{-t\ell}\bigr).
\]
对 \(t=0\) 逐阶求导，便得到
\[
\kappa_1^\infty=-p\ell,
\qquad
\kappa_2^\infty=p(1-p)\ell^2,
\]
\[
\kappa_3^\infty=-p(1-p)(1-2p)\ell^3,
\qquad
\kappa_4^\infty=p(1-p)(1-6p+6p^2)\ell^4.
\]
证毕。
\end{proof}

\begin{theorem}[规范群熵的精确首项与常数漂移]\label{thm:xi-time-part70ad-gauge-entropy-leading-drift}
记
\[
G_m^{\mathrm{bin}}
\cong
\prod_{w\in X_m}S_{d_m^{\mathrm{bin}}(w)},
\qquad
H_m^{\mathrm{gauge}}:=\log_2|G_m^{\mathrm{bin}}|.
\]
则其自然对数版本满足
\[
\boxed{\;
\frac{1}{2^m}\log |G_m^{\mathrm{bin}}|
=
m\log\!\Bigl(\frac{2}{\varphi}\Bigr)
-1
-\frac{\log\varphi}{\varphi\sqrt5}
+o(1).\;}
\]
等价地，以 bits 计，
\[
\boxed{\;
H_m^{\mathrm{gauge}}
=
2^m\left[
m\log_2\!\Bigl(\frac{2}{\varphi}\Bigr)
-\frac{1}{\log 2}
\Bigl(
1+\frac{\log\varphi}{\varphi\sqrt5}
\Bigr)
+o(1)
\right].\;}
\]
\end{theorem}
\begin{proof}
由定理 \ref{thm:xi-fold-hidden-logmultiplicity-kappa-gauge-two-scale-explicit}，若记
\[
g_m:=2^{-m}\log|G_m^{\mathrm{bin}}|,
\]
则有
\[
g_m
=
m\log\!\Bigl(\frac{2}{\varphi}\Bigr)
-\Bigl(1+\frac{\log\varphi}{1+\varphi^2}\Bigr)
+O\!\Bigl(\Bigl(\frac{\varphi}{2}\Bigr)^m\Bigr).
\]
再用恒等式
\[
\frac{1}{1+\varphi^2}=\frac{1}{\varphi\sqrt5},
\]
即得
\[
\frac{1}{2^m}\log |G_m^{\mathrm{bin}}|
=
m\log\!\Bigl(\frac{2}{\varphi}\Bigr)
-1
-\frac{\log\varphi}{\varphi\sqrt5}
+o(1).
\]
最后两边乘以 \(2^m/\log 2\)，便得到 bits 形式。证毕。
\end{proof}

\begin{theorem}[bin-fold 输出分布的 Tsallis 熵同样具有闭式常数层]\label{thm:xi-time-part70ad-binfold-tsallis-entropy-asymptotics}
\leanverified{paper\_xi\_time\_part70ad\_binfold\_tsallis\_entropy\_asymptotics}
对 \(q>0\)、\(q\neq 1\)，定义 Tsallis 熵
\[
T_q(\mu_m):=\frac{1-\sum_{w\in X_m}\mu_m(w)^q}{q-1}.
\]
则有如下两段渐近：
\begin{enumerate}
  \item 若 \(q>1\)，则
  \[
  \boxed{\;
  T_q(\mu_m)
  =
  \frac{1}{q-1}
  -
  \frac{1}{q-1}\,
  \varphi^{(1-q)m}
  \frac{\varphi+\varphi^{-q}}{\sqrt5}
  +
  O_q\!\Bigl(\varphi^{(1-q)m}\Bigl(\frac{\varphi}{2}\Bigr)^m\Bigr).\;}
  \]
  特别地，
  \[
  \boxed{\;
  T_q(\mu_m)\longrightarrow \frac{1}{q-1}.\;}
  \]
  \item 若 \(0<q<1\)，则
  \[
  \boxed{\;
  T_q(\mu_m)
  =
  \frac{1}{1-q}\,
  \varphi^{(1-q)m}
  \frac{\varphi+\varphi^{-q}}{\sqrt5}
  \bigl(1+o(1)\bigr).\;}
  \]
  因而
  \[
  \boxed{\;
  \lim_{m\to\infty}\frac{1}{m}\log T_q(\mu_m)
  =
  (1-q)\log\varphi.\;}
  \]
\end{enumerate}
\end{theorem}
\begin{proof}
由定理 \ref{thm:xi-time-part70ad-binfold-renyi-entropy-constant-drift} 的证明中间式，
\[
\sum_{w\in X_m}\mu_m(w)^q
=
\varphi^{(1-q)m}
\frac{\varphi+\varphi^{-q}}{\sqrt5}
\Bigl(1+O_q\!\Bigl(\Bigl(\frac{\varphi}{2}\Bigr)^m\Bigr)\Bigr).
\]
若 \(q>1\)，则 \(1-q<0\)，上式趋于 \(0\)。代入定义
\[
T_q(\mu_m)
=
\frac{1}{q-1}\left(1-\sum_{w\in X_m}\mu_m(w)^q\right),
\]
便立即得到第一段公式与极限。

若 \(0<q<1\)，则 \(1-q>0\)，故上式指数发散。于是
\[
T_q(\mu_m)
=
\frac{\sum_w\mu_m(w)^q-1}{1-q}
=
\frac{1}{1-q}\,
\varphi^{(1-q)m}
\frac{\varphi+\varphi^{-q}}{\sqrt5}
\bigl(1+o(1)\bigr),
\]
因为被减去的常数 \(1\) 相对于主项可忽略。最后取对数并除以 \(m\)，即可得到指数率
\[
\lim_{m\to\infty}\frac{1}{m}\log T_q(\mu_m)
=(1-q)\log\varphi.
\]
证毕。
\end{proof}

\noindent
由此，bin-fold 输出分布的全阶信息几何在常数层也达到闭合：Rényi 熵不只具有统一熵率，还具有显式 \(O(1)\) 漂移；中心化对数多重度不只具有一阶均值极限，而是整个指数矩与全部有限阶累积量都收敛到同一个二原子模型；规范群熵的指数体量则带有可审计的首项斜率与常数偏移；而 Tsallis 熵在 \(q>1\) 与 \(0<q<1\) 两侧的饱和值与增长率也由同一碰撞矩常数统一给出。于是，两态末位机制、碰撞矩闭式与规范体积展开最终在同一信息几何层面完全合流。

\endinput


\paragraph{中心双折点容量骨架、escort 温度重加权与 dark band 半阶尺度}
在 bin-fold 的中心临界容量极限、escort 温度族末位塌缩、尾部单跳阶梯以及六倍窗零块下界都已建立之后，还可把阈值/采样轴、相位/路径信息轴与 Fourier 侧 dark band 进一步压缩为同一组闭合接口。其共同特征是：中心层容量骨架始终只含两个不可消去折点，而温度参数只改变同一双折点骨架上的权重，不改变骨架本身。

\begin{theorem}[中心临界容量骨架的双折点闭式]\label{thm:xi-time-part70b-center-capacity-two-kink-skeleton}
定义
\[
A_m:=\left(\frac{2}{\varphi}\right)^m,
\qquad
\Psi_m(\tau):=
2^{-m}\,\mathcal C_m^{\mathrm{bin,cont}}(\tau A_m)
\qquad (\tau\ge 0).
\]
则对每个固定 \(\tau\ge 0\)，都有
\[
\Psi_m(\tau)\longrightarrow \Psi(\tau)
:=
\frac{\varphi}{\sqrt5}\min\{1,\tau\}
+
\frac{1}{\sqrt5}\min\{\varphi^{-1},\tau\}.
\]
等价地，
\[
\Psi(\tau)=
\begin{cases}
\dfrac{\varphi^2}{\sqrt5}\,\tau, & 0\le \tau\le \varphi^{-1},\\[6pt]
\dfrac{\varphi}{\sqrt5}\,\tau+\dfrac{1}{\varphi\sqrt5}, & \varphi^{-1}\le \tau\le 1,\\[8pt]
1, & \tau\ge 1.
\end{cases}
\]
特别地，中心临界容量骨架只有两个不可消去的折点，位置分别位于
\[
\tau=\varphi^{-1},
\qquad
\tau=1.
\]
\end{theorem}
\begin{proof}
这是定理 \ref{thm:xi-foldbin-dyadic-capacity-critical-window-limit} 与定理 \ref{thm:xi-time-part70-continuous-capacity-two-atom-second-derivative} 的直接改写。分段表达式由极限函数逐段展开得到，而折点位置则由分布二阶导数
\[
-D^2\Psi
=
\frac{1}{\sqrt5}\,\delta_{\varphi^{-1}}
+
\frac{\varphi}{\sqrt5}\,\delta_1
\]
唯一确定。证毕。
\end{proof}

\begin{theorem}[escort 温度族对中心双折点骨架的重加权响应]\label{thm:xi-time-part70b-escort-two-kink-response}
对任意固定 \(q\ge 0\)，令 \(W_m^{(q)}\sim \pi_{m,q}\)，并定义
\[
\Psi_{m,q}^{\mathrm{esc}}(\tau)
:=
\mathbb E_{\pi_{m,q}}
\left[
\min\!\left(
1,\frac{\tau A_m}{d_m^{\mathrm{bin}}(W_m^{(q)})}
\right)
\right]
\qquad (\tau\ge 0).
\]
则对每个固定 \(\tau\ge 0\)，极限
\[
\Psi_q^{\mathrm{esc}}(\tau):=
\lim_{m\to\infty}\Psi_{m,q}^{\mathrm{esc}}(\tau)
\]
存在，并满足闭式
\[
\Psi_q^{\mathrm{esc}}(\tau)
=
\bigl(1-\theta(q)\bigr)\min\{1,\tau\}
+
\theta(q)\min\{1,\varphi\tau\},
\qquad
\theta(q):=\frac{1}{1+\varphi^{q+1}}.
\]
因此，escort 温度参数 \(q\) 并不改变中心容量骨架的两个折点位置，而只在同一双折点框架上重新分配两条线性支的权重。
\end{theorem}
\begin{proof}
由定理 \ref{thm:fold-bin-two-state-asymptotic}，存在常数 \(C>0\)，使对充分大的 \(m\) 与全部 \(w\in X_m\) 都有
\[
\left|
\frac{d_m^{\mathrm{bin}}(w)}{A_m}-\varphi^{-w_m}
\right|
\le C\left(\frac{\varphi}{2}\right)^m.
\]
故对充分大的 \(m\)，上述比值一律落在紧区间
\[
\left[\frac{\varphi^{-1}}{2},\,2\right]
\]
内。对固定 \(\tau\ge 0\)，函数
\[
x\longmapsto \min\!\left(1,\frac{\tau}{x}\right)
\]
在该紧区间上 Lipschitz，因而存在常数 \(C_\tau>0\) 使
\[
\min\!\left(1,\frac{\tau A_m}{d_m^{\mathrm{bin}}(w)}\right)
=
\min\!\bigl(1,\tau\varphi^{w_m}\bigr)
+O\!\left(\left(\frac{\varphi}{2}\right)^m\right)
\]
在 \(w\in X_m\) 上一致成立。对 \(\pi_{m,q}\) 取期望，得到
\[
\Psi_{m,q}^{\mathrm{esc}}(\tau)
=
\mathbb E_{\pi_{m,q}}\!\bigl[\min(1,\tau\varphi^{W_m^{(q)}(m)})\bigr]
+O\!\left(\left(\frac{\varphi}{2}\right)^m\right).
\]
另一方面，由命题 \ref{prop:fold-bin-escort-lastbit}，
\[
\mathbb P\bigl(W_m^{(q)}(m)=1\bigr)
=
\theta(q)+O\!\left(\left(\frac{\varphi}{2}\right)^m\right),
\]
\[
\mathbb P\bigl(W_m^{(q)}(m)=0\bigr)
=
1-\theta(q)+O\!\left(\left(\frac{\varphi}{2}\right)^m\right).
\]
于是
\[
\Psi_{m,q}^{\mathrm{esc}}(\tau)
=
\bigl(1-\theta(q)\bigr)\min\{1,\tau\}
+
\theta(q)\min\{1,\varphi\tau\}
+O\!\left(\left(\frac{\varphi}{2}\right)^m\right).
\]
令 \(m\to\infty\)，即得所示极限公式。证毕。
\end{proof}

\begin{corollary}[中心双折点骨架上的温度斜率可辨识性]\label{cor:xi-time-part70b-escort-slope-identifiability}
沿用定理 \ref{thm:xi-time-part70b-escort-two-kink-response} 的记号。则在中间区间
\[
\tau\in (\varphi^{-1},1)
\]
上有
\[
\Psi_q^{\mathrm{esc}}(\tau)
=
\bigl(1-\theta(q)\bigr)\tau+\theta(q),
\]
从而导数恒为
\[
\frac{\dd}{\dd\tau}\Psi_q^{\mathrm{esc}}(\tau)
=
\sigma_q
:=
1-\theta(q)\in(0,1).
\]
并且仅由该斜率 \(\sigma_q\) 即可唯一恢复温度参数 \(q\)：
\[
q=\log_{\varphi}\!\left(\frac{\sigma_q}{1-\sigma_q}\right)-1.
\]
\end{corollary}
\begin{proof}
当 \(\tau\in(\varphi^{-1},1)\) 时，
\[
\min\{1,\tau\}=\tau,
\qquad
\min\{1,\varphi\tau\}=1,
\]
故定理 \ref{thm:xi-time-part70b-escort-two-kink-response} 中的极限曲线退化为
\[
\Psi_q^{\mathrm{esc}}(\tau)
=
\bigl(1-\theta(q)\bigr)\tau+\theta(q).
\]
于是导数恒为 \(1-\theta(q)\)。再由
\[
\theta(q)=\frac{1}{1+\varphi^{q+1}}
\]
可得
\[
\frac{1-\theta(q)}{\theta(q)}=\varphi^{q+1}.
\]
代入 \(\sigma_q=1-\theta(q)\) 并整理，即得所示反演公式。证毕。
\end{proof}

\begin{corollary}[固定末位层的阈值层纹由单跳尾和谱完全恢复]\label{cor:xi-time-part70b-lastbit-tail-jump-recovery}
沿用定理 \ref{thm:xi-foldbin-lastbit-tail-spectrum-single-jump-staircase} 的记号，并把其中的末位指标 \(i\) 记为 \(b\in\{0,1\}\)。定义
\[
\delta_{m,b}(V):=
D_{m,b}(V)
=
\#\bigl\{t\in \mathcal T_{m,b}: t\le 2^m-1-V\bigr\}.
\]
则对一切相容的 \(V\) 都有精确跳跃公式
\[
\delta_{m,b}(V)-\delta_{m,b}(V+1)
=
\mathbf 1_{\{\,2^m-1-V\in\mathcal T_{m,b}\,\}}.
\]
因而合法尾和集可由跳点完全恢复：
\[
\mathcal T_{m,b}
=
\bigl\{2^m-1-V:\ \delta_{m,b}(V)-\delta_{m,b}(V+1)=1\bigr\}.
\]
特别地，在固定末位相位内，每一次向下跳变都恰对应一个合法尾和被阈值切除。
\end{corollary}
\begin{proof}
这只是定理 \ref{thm:xi-foldbin-lastbit-tail-spectrum-single-jump-staircase} 的盒装恒等式与其逆向重写。证毕。
\end{proof}

\begin{theorem}[六倍窗 exact dark band 的半阶指数尺度]\label{thm:xi-time-part70b-dark-band-half-order-scale}
\leanverified{paper\_xi\_time\_part70b\_dark\_band\_half\_order\_scale}
沿用定理 \ref{thm:fold-zero-density-sparse} 与推论 \ref{cor:fold-zero-half-index-multiple6} 的记号，记
\[
\mathcal Z_m:=\{k\in \ZZ/(F_{m+2}\ZZ):\ \widehat d_m(k)=0\},
\]
并令 \(\tau_{\mathrm{div}}(n)\) 表示 \(n\) 的约数个数函数。若
\[
m+2\equiv 0\pmod 6,
\]
则有双边界
\[
F_{(m+2)/2}
\le
|\mathcal Z_m|
\le
\tau_{\mathrm{div}}(m+2)\,F_{(m+2)/2}.
\]
从而沿该六倍窗子序列，
\[
\log|\mathcal Z_m|
=
\frac{m}{2}\log\varphi+O(\log m),
\qquad
\frac{1}{m}\log|\mathcal Z_m|
\longrightarrow
\frac12\log\varphi.
\]
换言之，exact dark band 在六倍窗上具有 \(\varphi^{m/2}\) 级的半阶指数尺度。
\end{theorem}
\begin{proof}
下界正是推论 \ref{cor:fold-zero-half-index-multiple6} 的内容。上界则由定理 \ref{thm:fold-zero-density-sparse} 给出：
\[
|\mathcal Z_m|
\le
\tau_{\mathrm{div}}(m+2)\,F_{\lfloor (m+2)/2\rfloor}
=
\tau_{\mathrm{div}}(m+2)\,F_{(m+2)/2}.
\]
另一方面，Binet 公式给出
\[
\log F_n=n\log\varphi+O(1),
\]
而初等约数界
\[
\tau_{\mathrm{div}}(n)\le 2\sqrt n
\qquad (n\ge 1)
\]
蕴含
\[
\log\tau_{\mathrm{div}}(m+2)=O(\log m).
\]
把这些估计代入上下界，便得到
\[
\log|\mathcal Z_m|
=
\frac{m}{2}\log\varphi+O(\log m),
\]
并进一步推出极限
\[
\frac{1}{m}\log|\mathcal Z_m|\to \frac12\log\varphi.
\]
证毕。
\end{proof}

\noindent
上述结果把中心读出、escort 重加权、阈值层纹与 Fourier 侧零块压缩到同一条双折点--跳点主线上：在阈值/采样轴上，中心容量骨架只允许两个折点 \(\varphi^{-1}\) 与 \(1\)；在相位/路径信息轴上，温度参数 \(q\) 仅通过 \(\theta(q)\) 重新分配两支权重；在固定末位相位内，层纹由合法尾和集的单跳删除精确生成；而六倍窗上的 exact dark band 则在 Fourier 侧取得 \(\varphi^{m/2}\) 级的半阶指数尺度。因而，相应可视化中的连续着色纹理与局部波动若不改变这些离散骨架，便不能改变中心容量、尾和跳点结构与 dark-band 的主尺度。

\endinput


\paragraph{谱零点的高阶 \texorpdfstring{$\ZZ_2$}{Z2} 轨道抵消}
在谱零点的单坐标判别、余类并结构与六倍窗 half-order dark band 已经建立之后，还可进一步把单个按位翻转对合的相消推广为一个多坐标 \(\ZZ_2\)-轨道抵消机制。其结果是：一旦某个频率 \(k\) 触发若干独立坐标上的 \(-1\) 相位翻转，则整个 \(\Omega_m\) 会按自由的 \((\ZZ/2\ZZ)^r\)-轨道分块，而每个轨道上的相位和都逐块归零。

\begin{theorem}[谱零点的高阶 \texorpdfstring{$\ZZ_2$}{Z2} 轨道抵消]\label{thm:xi-time-part70c-spectrum-zero-z2-orbit-cancellation}
\leanverified{paper\_xi\_time\_part70c\_spectrum\_zero\_z2\_orbit\_cancellation}
沿用定理 \ref{thm:fold-spectrum-zero-criterion} 与推论 \ref{cor:fold-spectrum-zero-involution} 的记号，记
\[
M:=F_{m+2},
\qquad
\widehat d_m(k)=\sum_{\omega\in\Omega_m}e^{-2\pi i kN(\omega)/M}.
\]
对任意 \(k\in \ZZ/(M\ZZ)\)，定义
\[
J(k):=
\left\{
j\in\{1,\dots,m\}:\ 2kF_{j+1}\equiv M\pmod{2M}
\right\},
\]
并令
\[
\Phi_k(\omega):=e^{-2\pi i kN(\omega)/M}
\qquad
(\omega\in \Omega_m).
\]
对每个 \(j\in J(k)\)，记
\[
\iota_j:\Omega_m\to \Omega_m,
\qquad
\iota_j(\omega):=\omega\oplus e_j.
\]
再定义
\[
G_k:=\langle \iota_j:\ j\in J(k)\rangle.
\]
则成立：
\begin{enumerate}
  \item \(G_k\cong (\ZZ/2\ZZ)^{|J(k)|}\)，且其在 \(\Omega_m\) 上的作用自由；
  \item 对任意
  \[
  g=\prod_{j\in S}\iota_j\in G_k,
  \]
  都有
  \[
  \Phi_k(g\omega)=(-1)^{|S|}\Phi_k(\omega)
  \qquad
  (\omega\in\Omega_m);
  \]
  \item 若 \(J(k)\neq \varnothing\)，则对任意 \(G_k\)-轨道 \(\mathcal O\subset \Omega_m\)，都有
  \[
  \sum_{\omega\in\mathcal O}\Phi_k(\omega)=0;
  \]
  特别地，
  \[
  \widehat d_m(k)=0
  \qquad\Longleftrightarrow\qquad
  J(k)\neq \varnothing.
  \]
\end{enumerate}
\end{theorem}
\begin{proof}
由定理 \ref{thm:fold-spectrum-zero-criterion}，
\[
\widehat d_m(k)=0
\qquad\Longleftrightarrow\qquad
\exists\,j\in\{1,\dots,m\}\ \text{使}\ 2kF_{j+1}\equiv M\pmod{2M},
\]
这正等价于 \(J(k)\neq \varnothing\)。

对每个 \(j\in J(k)\)，推论 \ref{cor:fold-spectrum-zero-involution} 已给出：\(\iota_j\) 是 \(\Omega_m\) 上无固定点的对合，且
\[
\Phi_k(\iota_j\omega)=-\Phi_k(\omega).
\]
若 \(j,j'\in J(k)\) 且 \(j\neq j'\)，则 \(\iota_j,\iota_{j'}\) 只翻转不同坐标，因此彼此交换，且各自平方为恒等。于是
\[
G_k\cong (\ZZ/2\ZZ)^{|J(k)|}.
\]
任一非平凡元素 \(g=\prod_{j\in S}\iota_j\) 都会翻转至少一个坐标，故不可能固定任何 \(\omega\in\Omega_m\)；从而 \(G_k\) 在 \(\Omega_m\) 上自由作用。

再对
\[
g=\prod_{j\in S}\iota_j
\]
重复使用 \(\Phi_k(\iota_j\omega)=-\Phi_k(\omega)\)，便得
\[
\Phi_k(g\omega)=(-1)^{|S|}\Phi_k(\omega).
\]
因此当 \(J(k)\neq \varnothing\) 时，对任意轨道
\[
\mathcal O=G_k\cdot \omega
\]
都有
\[
\sum_{\eta\in\mathcal O}\Phi_k(\eta)
=
\Phi_k(\omega)\sum_{S\subseteq J(k)}(-1)^{|S|}
=
\Phi_k(\omega)(1-1)^{|J(k)|}
=0.
\]
把 \(\Omega_m\) 分解为 \(G_k\)-轨道的并，即得
\[
\widehat d_m(k)=\sum_{\omega\in\Omega_m}\Phi_k(\omega)=0.
\]
与第一步的零点判别合并，即得最后的充要条件。证毕。
\end{proof}

\endinput


\paragraph{六倍窗零点密度的精确指数率与 KL--最大纤维耦合}
在谱零点的高阶 \texorpdfstring{$(\ZZ_2)^r$}{(Z2)r} 轨道抵消已经给出逐点见证之后，还可继续沿频域--信息主线抽取出下列两条结果。它们分别把六倍窗子序列上的归一化零点密度锁定在严格的半维指数率 \(\varphi^{-m/2+o(m)}\)，并把输出分布相对于均匀分布的 KL 缺口与最大纤维超额压缩为同一条确定性控制不等式。

\begin{theorem}[六倍窗上的零点密度具有精确指数率]\label{thm:xi-time-part70c-zero-density-exact-ratio-exponent-multiple6}
\leanverified{paper\_xi\_time\_part70c\_zero\_density\_exact\_ratio\_exponent\_multiple6}
若
\[
m+2\equiv 0\pmod 6,
\qquad
d:=\frac{m+2}{2},
\qquad
M:=F_{m+2},
\qquad
\mathcal Z_m:=\{k\in \ZZ/(M\ZZ):\ \widehat d_m(k)=0\},
\]
则有极限
\[
\lim_{\substack{m\to\infty\\ m+2\equiv 0\ (6)}}
\frac1m\log\frac{|\mathcal Z_m|}{F_{m+2}}
=
-\frac12\log\varphi.
\]
等价地，在六倍窗子序列上，
\[
\frac{|\mathcal Z_m|}{F_{m+2}}
=
\varphi^{-m/2+o(m)}.
\]
\end{theorem}
\begin{proof}
下界方面，由推论 \ref{cor:fold-zero-half-index-multiple6}，
\[
|\mathcal Z_m|
\ge
F_d.
\]
因此
\[
\frac{|\mathcal Z_m|}{F_{m+2}}
\ge
\frac{F_d}{F_{m+2}}.
\]
由 Binet 公式 \(F_n=\varphi^{n+o(n)}\) 得
\[
\liminf_{\substack{m\to\infty\\ m+2\equiv 0\ (6)}}
\frac1m\log\frac{|\mathcal Z_m|}{F_{m+2}}
\ge
\lim_{m\to\infty}
\frac1m\log\frac{F_d}{F_{m+2}}
=
-\frac12\log\varphi.
\]

上界方面，由定理 \ref{thm:fold-zero-density-sparse}，
\[
\frac{|\mathcal Z_m|}{F_{m+2}}
\le
\tau(m+2)\frac{F_{\lfloor (m+2)/2\rfloor}}{F_{m+2}}.
\]
又由于
\[
\tau(n)=e^{o(n)},
\qquad
F_{\lfloor (m+2)/2\rfloor}= \varphi^{m/2+o(m)},
\qquad
F_{m+2}=\varphi^{m+o(m)},
\]
故
\[
\limsup_{\substack{m\to\infty\\ m+2\equiv 0\ (6)}}
\frac1m\log\frac{|\mathcal Z_m|}{F_{m+2}}
\le
-\frac12\log\varphi.
\]
上下界相合，即得结论。
\end{proof}

\begin{theorem}[最大纤维超额由 KL 缺口确定性控制]\label{thm:xi-time-part70c-kl-gap-controls-maxfiber-excess}
设
\[
p_m(r):=\frac{d_m(r)}{2^m},
\qquad
M:=F_{m+2},
\qquad
u(r):=\frac1M,
\]
并记
\[
M_m:=\max_r d_m(r),
\qquad
d_m^{\mathrm{avg}}:=\frac{2^m}{M}.
\]
则有确定性耦合不等式
\[
\boxed{\;
M_m-d_m^{\mathrm{avg}}
\le
2^m\,\|p_m-u\|_{\mathrm{TV}}
\le
2^m\sqrt{\frac12\,D_{\mathrm{KL}}(p_m\|u)}.\;}
\]
特别地，一旦 KL 缺口被控制，最大纤维相对于平均纤维的超额便被同步控制。
\leanverified{paper\_xi\_time\_part70c\_kl\_gap\_controls\_maxfiber\_excess}
\end{theorem}
\begin{proof}
首先由
\[
M_m=2^m\max_r p_m(r),
\qquad
d_m^{\mathrm{avg}}=\frac{2^m}{M}=2^m u(r),
\]
可得
\[
M_m-d_m^{\mathrm{avg}}
=
2^m\left(\max_r p_m(r)-\frac1M\right)
\le
2^m\max_r|p_m(r)-u(r)|.
\]
另一方面，对任意概率分布 \(p,u\) 都有
\[
\max_r|p(r)-u(r)|
\le
\frac12\sum_r|p(r)-u(r)|
=
\|p-u\|_{\mathrm{TV}},
\]
故
\[
M_m-d_m^{\mathrm{avg}}
\le
2^m\,\|p_m-u\|_{\mathrm{TV}}.
\]
最后应用 Pinsker 不等式
\[
\|p_m-u\|_{\mathrm{TV}}
\le
\sqrt{\frac12\,D_{\mathrm{KL}}(p_m\|u)},
\]
便得到所示结论。证毕。
\end{proof}

\noindent
由此，谱零点结构在三个层级上完成闭合：在单频层级，高阶 \texorpdfstring{$(\ZZ_2)^r$}{(Z2)r} 轨道抵消给出逐点见证；在密度层级，六倍窗上的归一化零点比例被锁定在 \(\varphi^{-m/2+o(m)}\)；在信息层级，KL 缺口与最大纤维超额被耦合为同一条确定性不等式。于是，轨道抵消、半维稀疏率与信息偏置便被压缩进同一条完全可审计的频域主线。

\endinput


\paragraph{末位终板刚性、临界预言机振荡与 Riesz 能量常数}
在末位相位的最终严格分层、bin-fold 规范群的最终满秩化、critical-scale 预言机容量骨架、escort 温度族的双折点响应以及六倍窗 exact dark band 的半阶指数尺度已经确立之后，还可把终板几何、商结构障碍、临界预算振荡与二阶能量常数进一步压缩为同一条后续链。在本段中统一记
\[
d_m(w):=d_m^{\mathrm{bin}}(w),
\qquad
\lambda:=\frac{2}{\varphi},
\qquad
\Omega_m:=\{0,1\}^m.
\]
下列结果分别刻画最小扇区的末位实现、大窗口上的全局非商化刚性、异常签名维数的 Fibonacci 饱和、临界比特率下的不可压平簇集带、escort 温度族的临界响应带、中心纤维场的精确二阶常数、碰撞概率与非平凡 Fourier 能量的精确主项，以及最大纤维超额的显式主系数。

\begin{theorem}[最小扇区的末位实现与终板逼近]\label{thm:xi-time-part70d-minsector-terminal-sheet-realization}
定义
\[
d_{\min}(m):=\min_{v\in X_m}d_m(v),
\qquad
A_1^{(m)}:=\{w\in X_m:\ w_m=1\},
\qquad
S_{m,\min}:=\{w\in X_m:\ d_m(w)=d_{\min}(m)\}.
\]
若“子覆盖同构”恒等式
\[
|S_{m,\min}|=|X_{m-2}|
\]
沿某个趋于无穷的偶数序列成立，则沿同一序列有
\[
\frac{|S_{m,\min}\triangle A_1^{(m)}|}{|X_m|}\longrightarrow 0.
\]
\end{theorem}
\leanverified{paper\_xi\_time\_part70d\_minsector\_terminal\_sheet\_realization}
\begin{proof}
由定理 \ref{thm:xi-time-part70ab-terminal-bit-eventual-strict-stratification}，对充分大的 \(m\) 有
\[
S_{m,\min}\subseteq A_1^{(m)}.
\]
另一方面，在假设下
\[
\frac{|S_{m,\min}|}{|X_m|}
=
\frac{|X_{m-2}|}{|X_m|}
=
\frac{F_m}{F_{m+2}}
\longrightarrow
\varphi^{-2},
\]
这也正是推论 \ref{cor:fold-bin-minsector-density-phi-minus2} 的内容。又由稳定词计数，
\[
|A_1^{(m)}|=F_m,
\qquad
|X_m|=F_{m+2},
\]
故
\[
\frac{|A_1^{(m)}|}{|X_m|}=\frac{F_m}{F_{m+2}}\longrightarrow \varphi^{-2}.
\]
由于 \(S_{m,\min}\subseteq A_1^{(m)}\) 且两者相对密度极限相同，遂有
\[
\frac{|A_1^{(m)}\setminus S_{m,\min}|}{|X_m|}\longrightarrow 0.
\]
而此时
\[
S_{m,\min}\triangle A_1^{(m)}=A_1^{(m)}\setminus S_{m,\min},
\]
故结论成立。
\end{proof}

\begin{theorem}[大窗口上的全局非商化刚性]\label{thm:xi-time-part70d-global-nonquotient-rigidity}
存在整数 \(m_0\)，使得对一切 \(m\ge m_0\)，等价关系
\[
\omega\sim_m\eta
\qquad\Longleftrightarrow\qquad
\Fold_m^{\mathrm{bin}}(\omega)=\Fold_m^{\mathrm{bin}}(\eta)
\]
既不能由任何有限群在 \(\Omega_m\) 上的自由作用之轨道划分实现，也不能来自任何有限阿贝尔群上的线性陪集划分。更准确地，不存在有限阿贝尔群 \(A_m\)、子群 \(H_m\le A_m\) 及双射
\[
\iota_m:\Omega_m\longrightarrow A_m
\]
使得
\[
\omega\sim_m\eta
\qquad\Longleftrightarrow\qquad
\iota_m(\omega)-\iota_m(\eta)\in H_m.
\]
\end{theorem}
\leanverified{paper\_xi\_time\_part70d\_global\_nonquotient\_rigidity}
\begin{proof}
记
\[
A_0^{(m)}:=\{w\in X_m:\ w_m=0\},
\qquad
A_1^{(m)}:=\{w\in X_m:\ w_m=1\}.
\]
由定理 \ref{thm:xi-time-part70ab-terminal-bit-eventual-strict-stratification}，存在 \(m_0\) 使得当 \(m\ge m_0\) 时，
\[
\min_{w\in A_0^{(m)}}d_m(w)>\max_{v\in A_1^{(m)}}d_m(v).
\]
由于 \(A_0^{(m)}\) 与 \(A_1^{(m)}\) 都非空，故对每个 \(m\ge m_0\)，纤维大小在同一窗口内至少取两个不同整数值。

若 \(\sim_m\) 来自某个有限群 \(G_m\) 的自由作用，则每个轨道大小都恒等于 \(|G_m|\)，因而所有纤维必同势；矛盾。

若 \(\sim_m\) 来自上式所述的线性陪集划分，则每个等价类都是某个陪集的原像，其大小恒等于 \(|H_m|\)，故所有纤维同样必同势；再度矛盾。于是两类实现均不可能存在。
\end{proof}

\begin{theorem}[异常签名维数的渐近饱和]\label{thm:xi-time-part70d-anomaly-signature-saturation}
\leanverified{paper\_xi\_time\_part70d\_anomaly\_signature\_saturation}
记
\[
G_m^{\mathrm{bin}}:=\prod_{w\in X_m}\mathfrak S_{d_m(w)}.
\]
则存在整数 \(m_0\)，使得对一切 \(m\ge m_0\)，都有
\[
\bigl(G_m^{\mathrm{bin}}\bigr)^{\mathrm{ab}}
\cong
(\ZZ/2\ZZ)^{|X_m|}
\cong
(\ZZ/2\ZZ)^{F_{m+2}}.
\]
因而任何保持“纤维独立置换”语义、并在 \(\FF_2\) 上完备分离 \(\bigl(G_m^{\mathrm{bin}}\bigr)^{\mathrm{ab}}\) 全部坐标的异常见证系统，其最小比特数恰为
\[
F_{m+2}.
\]
\end{theorem}
\begin{proof}
这正是定理 \ref{thm:xi-time-part65-binfold-center-vanishing-signature-stability} 的重述。该定理已证明：对充分大的 \(m\)，有
\[
\bigl(G_m^{\mathrm{bin}}\bigr)^{\mathrm{ab}}
\cong
(\ZZ/2\ZZ)^{F_{m+2}},
\qquad
\operatorname{rank}_{\FF_2}\!\bigl((G_m^{\mathrm{bin}})^{\mathrm{ab}}\bigr)=F_{m+2},
\]
并且纤维奇偶荷的最小完备秩在同一范围内严格等于 \(F_{m+2}\)。这恰好给出所述异常签名维数结论。
\end{proof}

\begin{theorem}[均匀微态输入下的临界比特率相变与簇集区间]\label{thm:xi-time-part70d-critical-bitrate-phase-cluster}
记
\[
\alpha:=\log_2\!\left(\frac{2}{\varphi}\right),
\qquad
\mathrm{Succ}_m(B):=\mathrm{Succ}_m^{\mathrm{bin}}(B).
\]
则成立：
\begin{enumerate}
  \item 若 \(\rho<\alpha\)，则
  \[
  \mathrm{Succ}_m(\lfloor \rho m\rfloor)\longrightarrow 0.
  \]
  \item 若 \(\rho>\alpha\)，则
  \[
  \mathrm{Succ}_m(\lfloor \rho m\rfloor)\longrightarrow 1.
  \]
  \item 对任意固定 \(\beta\in\RR\)，令
  \[
  B_m:=\lfloor \alpha m+\beta\rfloor.
  \]
  则序列 \(\bigl(\mathrm{Succ}_m(B_m)\bigr)_{m\ge 1}\) 不收敛，其全部簇集点恰组成闭区间
  \[
  \left[\frac{\varphi^2}{2\sqrt5},\,1\right].
  \]
\end{enumerate}
\end{theorem}
\begin{proof}
记 \(\lambda:=2/\varphi\)，则 \(\alpha=\log_2\lambda\)。

若 \(\rho<\alpha\)，则由平凡上界
\[
C_m^{\mathrm{bin}}(B)\le 2^B|X_m|
\]
得到
\[
\mathrm{Succ}_m(\lfloor \rho m\rfloor)
\le
\frac{2^{\lfloor \rho m\rfloor}F_{m+2}}{2^m}
=
\frac{2^{\lfloor \rho m\rfloor}}{\lambda^m}\cdot \frac{F_{m+2}}{\varphi^m}
\longrightarrow 0,
\]
因为 \(2^{\lfloor \rho m\rfloor}/\lambda^m\to 0\) 且 \(F_{m+2}/\varphi^m\to \varphi^2/\sqrt5\)。

若 \(\rho>\alpha\)，则 \(2^{\lfloor \rho m\rfloor}/\lambda^m\to+\infty\)。由推论 \ref{cor:xi-time-part70ab-extremal-fibers-terminal-bit-localization}，
\[
M_m:=\max_{w\in X_m}d_m(w)=\lambda^m+O(1).
\]
故对充分大的 \(m\) 有 \(2^{\lfloor \rho m\rfloor}>M_m\)，从而所有纤维都完全通过，即
\[
\mathrm{Succ}_m(\lfloor \rho m\rfloor)=1.
\]

现取临界序列 \(B_m=\lfloor \alpha m+\beta\rfloor\)。定义
\[
\theta_m:=\alpha m+\beta-\lfloor \alpha m+\beta\rfloor\in[0,1),
\qquad
\tau_m:=\frac{2^{B_m}}{\lambda^m}=2^{-\theta_m}\in[1/2,1].
\]
若 \(\log_2\varphi\in\QQ\)，则存在整数 \(p,q\ge 1\) 使 \(\varphi^q=2^p\in\QQ\)，这与 \(\varphi^q\notin\QQ\) 矛盾。故 \(\alpha=1-\log_2\varphi\) 为无理数，从而 \((\theta_m)\) 在 \([0,1]\) 上稠密，\((\tau_m)\) 的簇集点恰为整个区间 \([1/2,1]\)。

任取 \(\tau\in[1/2,1]\) 并取子列 \(m_j\) 使 \(\tau_{m_j}\to\tau\)。由定理 \ref{thm:xi-foldbin-dyadic-capacity-critical-window-limit}，
\[
\mathrm{Succ}_{m_j}(B_{m_j})
\longrightarrow
\Psi(\tau),
\]
其中
\[
\Psi(\tau)=
\begin{cases}
\dfrac{\varphi^2}{\sqrt5}\,\tau, & 0\le \tau\le \varphi^{-1},\\[6pt]
\dfrac{\varphi}{\sqrt5}\,\tau+\dfrac{1}{\varphi\sqrt5}, & \varphi^{-1}\le \tau\le 1,\\[8pt]
1, & \tau\ge 1.
\end{cases}
\]
因此 \(\bigl(\mathrm{Succ}_m(B_m)\bigr)\) 的全部簇集点恰为连续单调函数 \(\Psi\) 在 \([1/2,1]\) 上的像，即
\[
\Psi([1/2,1])
=
\left[\Psi(1/2),\,\Psi(1)\right]
=
\left[\frac{\varphi^2}{2\sqrt5},\,1\right].
\]
该区间非单点，故序列不收敛。
\end{proof}

\begin{theorem}[escort 温度族下的临界响应带]\label{thm:xi-time-part70d-escort-critical-response-band}
对任意 \(q\ge 0\)，定义 escort 加权的阈值响应
\[
\mathrm{Succ}_{m}^{(q)}(B)
:=
\sum_{w\in X_m}\pi_{m,q}(w)\min\!\left(1,\frac{2^B}{d_m(w)}\right),
\qquad
\theta(q):=\frac{1}{1+\varphi^{q+1}}.
\]
则同一临界比特率
\[
\alpha=\log_2\!\left(\frac{2}{\varphi}\right)
\]
给出 sharp transition：
\begin{enumerate}
  \item 若 \(\rho<\alpha\)，则
  \[
  \mathrm{Succ}_{m}^{(q)}(\lfloor \rho m\rfloor)\longrightarrow 0.
  \]
  \item 若 \(\rho>\alpha\)，则
  \[
  \mathrm{Succ}_{m}^{(q)}(\lfloor \rho m\rfloor)\longrightarrow 1.
  \]
  \item 若对固定 \(\beta\in\RR\) 令
  \[
  B_m=\lfloor \alpha m+\beta\rfloor,
  \]
  则序列 \(\bigl(\mathrm{Succ}_{m}^{(q)}(B_m)\bigr)_{m\ge 1}\) 不收敛，其全部簇集点恰为
  \[
  \left[\frac{1+(\varphi-1)\theta(q)}{2},\,1\right].
  \]
\end{enumerate}
\end{theorem}
\begin{proof}
若 \(\rho<\alpha\)，则 \(2^{\lfloor \rho m\rfloor}/\lambda^m\to 0\)。由定理 \ref{thm:fold-bin-two-state-asymptotic}，存在常数 \(C>0\) 使对充分大的 \(m\) 与全部 \(w\in X_m\) 都有
\[
d_m(w)\ge \lambda^m\left(\varphi^{-1}-C\Bigl(\frac{\varphi}{2}\Bigr)^m\right).
\]
于是
\[
\mathrm{Succ}_{m}^{(q)}(\lfloor \rho m\rfloor)
\le
\frac{2^{\lfloor \rho m\rfloor}}{\lambda^m\left(\varphi^{-1}-C(\varphi/2)^m\right)}
\longrightarrow 0.
\]

若 \(\rho>\alpha\)，则记
\[
M_m:=\max_{w\in X_m}d_m(w).
\]
由推论 \ref{cor:xi-time-part70ab-extremal-fibers-terminal-bit-localization}，\(M_m=\lambda^m+O(1)\)。故对充分大的 \(m\) 有 \(2^{\lfloor \rho m\rfloor}>M_m\)，于是每个纤维都完全通过，从而
\[
\mathrm{Succ}_{m}^{(q)}(\lfloor \rho m\rfloor)=1.
\]

现令 \(B_m=\lfloor \alpha m+\beta\rfloor\)，并仍记
\[
\tau_m:=\frac{2^{B_m}}{\lambda^m}\in[1/2,1].
\]
由上一条定理的证明，\((\tau_m)\) 的全部簇集点恰为 \([1/2,1]\)。

记
\[
I:=[1/2,1],
\qquad
J:=\left[\frac{\varphi^{-1}}{2},\,2\right].
\]
函数
\[
f(\tau,x):=\min\!\left(1,\frac{\tau}{x}\right)
\]
在紧集 \(I\times J\) 上联合 Lipschitz。再由定理 \ref{thm:fold-bin-two-state-asymptotic}，
\[
\left|\frac{d_m(w)}{\lambda^m}-\varphi^{-w_m}\right|
\le
C\Bigl(\frac{\varphi}{2}\Bigr)^m
\qquad (w\in X_m),
\]
故对充分大的 \(m\) 与全部 \(w\in X_m\) 都有
\[
\min\!\left(1,\frac{2^{B_m}}{d_m(w)}\right)
=
f\!\left(\tau_m,\frac{d_m(w)}{\lambda^m}\right)
=
f\!\left(\tau_m,\varphi^{-w_m}\right)
+O\!\left(\left(\frac{\varphi}{2}\right)^m\right),
\]
且误差对 \(w\) 与 \(\tau_m\in I\) 一致。对 \(\pi_{m,q}\) 取期望，并用命题 \ref{prop:fold-bin-escort-lastbit}
\[
\mathbf P_{\pi_{m,q}}(w_m=1)=\theta(q)+O\!\left(\left(\frac{\varphi}{2}\right)^m\right),
\qquad
\mathbf P_{\pi_{m,q}}(w_m=0)=1-\theta(q)+O\!\left(\left(\frac{\varphi}{2}\right)^m\right),
\]
得到
\[
\mathrm{Succ}_{m}^{(q)}(B_m)
=
\bigl(1-\theta(q)\bigr)\min\{1,\tau_m\}
+
\theta(q)\min\{1,\varphi\tau_m\}
+o(1).
\]
任取 \(\tau\in[1/2,1]\) 及子列 \(m_j\) 使 \(\tau_{m_j}\to\tau\)，便有
\[
\mathrm{Succ}_{m_j}^{(q)}(B_{m_j})
\longrightarrow
\Psi_q^{\mathrm{esc}}(\tau),
\qquad
\Psi_q^{\mathrm{esc}}(\tau)
:=
\bigl(1-\theta(q)\bigr)\min\{1,\tau\}
+
\theta(q)\min\{1,\varphi\tau\}.
\]
故全部簇集点恰为 \(\Psi_q^{\mathrm{esc}}([1/2,1])\)。由于 \(\Psi_q^{\mathrm{esc}}\) 在 \([1/2,1]\) 上连续单调，故该像为区间
\[
\left[\Psi_q^{\mathrm{esc}}(1/2),\,\Psi_q^{\mathrm{esc}}(1)\right].
\]
又因 \(1/2<\varphi^{-1}\)，故
\[
\Psi_q^{\mathrm{esc}}(1/2)
=
\frac{1-\theta(q)}{2}+\frac{\varphi\theta(q)}{2}
=
\frac{1+(\varphi-1)\theta(q)}{2},
\qquad
\Psi_q^{\mathrm{esc}}(1)=1.
\]
于是所求簇集区间即为
\[
\left[\frac{1+(\varphi-1)\theta(q)}{2},\,1\right].
\]
其长度严格为正，故序列不收敛。
\end{proof}

\begin{theorem}[中心纤维场的精确二阶常数]\label{thm:xi-time-part70d-central-fiber-second-moment}
\leanverified{paper\_xi\_time\_part70d\_central\_fiber\_second\_moment}
令 \(U_m\sim \mathrm{Unif}(X_m)\)，并定义尺度化纤维场
\[
Y_m:=\left(\frac{\varphi}{2}\right)^m d_m(U_m).
\]
则有极限
\[
\mathbb E(Y_m^2)\longrightarrow 2\varphi^{-2},
\qquad
\operatorname{Var}(Y_m)\longrightarrow \varphi^{-7}.
\]
\end{theorem}
\begin{proof}
由推论 \ref{cor:xi-time-part70a-scaled-moments-leading-constant}，当 \(q=1,2\) 时分别有
\[
\mathbb E(Y_m)
\longrightarrow
\frac{1}{\varphi}+\varphi^{-3},
\qquad
\mathbb E(Y_m^2)
\longrightarrow
\frac{1}{\varphi}+\varphi^{-4}.
\]
利用恒等式 \(\varphi^2=\varphi+1\) 与 \(\sqrt5=2\varphi-1\)，可将上式化简为
\[
\frac{1}{\varphi}+\varphi^{-4}=2\varphi^{-2},
\qquad
\frac{1}{\varphi}+\varphi^{-3}=\frac{\sqrt5}{\varphi^2}.
\]
因此
\[
\operatorname{Var}(Y_m)
=
\mathbb E(Y_m^2)-\bigl(\mathbb E(Y_m)\bigr)^2
\longrightarrow
2\varphi^{-2}-\frac{5}{\varphi^4}.
\]
再用 \(\varphi^2=\varphi+1\) 化简，得到
\[
2\varphi^{-2}-\frac{5}{\varphi^4}=\varphi^{-7}.
\]
证毕。
\end{proof}

\begin{theorem}[碰撞概率与非平凡 Fourier 能量的精确主项]\label{thm:xi-time-part70d-collision-riesz-main-terms}
沿用定理 \ref{thm:xi-time-part70d-central-fiber-second-moment} 的记号，并沿用定理 \ref{thm:xi-time-part9p-chi2-ranktwo-resonant-floor} 的 Fourier 记号。记
\[
\mathsf{Col}_m:=
\frac{1}{4^m}\sum_{w\in X_m}d_m(w)^2.
\]
则有
\[
\mathsf{Col}_m\sim \frac{2}{\sqrt5}\,\varphi^{-m},
\qquad
F_{m+2}\,\mathsf{Col}_m\longrightarrow \frac{2\varphi^2}{5}.
\]
进一步，非平凡 Fourier 能量满足
\[
4^{-m}\sum_{k\ne 0}\bigl|\widehat d_m(k)\bigr|^2
\longrightarrow
\frac{1}{5\varphi^3}.
\]
\end{theorem}
\begin{proof}
由 \(U_m\sim\mathrm{Unif}(X_m)\) 的定义，
\[
\mathsf{Col}_m
=
\frac{|X_m|}{4^m}\,\mathbb E\!\bigl(d_m(U_m)^2\bigr).
\]
又因为
\[
d_m(U_m)=\lambda^m Y_m,
\qquad
\lambda=\frac{2}{\varphi},
\]
故
\[
\mathsf{Col}_m
=
|X_m|\,\varphi^{-2m}\,\mathbb E(Y_m^2).
\]
由 \(|X_m|=F_{m+2}\) 与定理 \ref{thm:xi-time-part70d-central-fiber-second-moment}，
\[
\mathsf{Col}_m
=
F_{m+2}\varphi^{-2m}\bigl(2\varphi^{-2}+o(1)\bigr).
\]
再用 Binet 公式
\[
F_{m+2}\sim \frac{\varphi^{m+2}}{\sqrt5},
\]
便得
\[
\mathsf{Col}_m
\sim
\frac{\varphi^{m+2}}{\sqrt5}\cdot \varphi^{-2m}\cdot 2\varphi^{-2}
=
\frac{2}{\sqrt5}\,\varphi^{-m}.
\]
两边再乘以 \(F_{m+2}\) 即得
\[
F_{m+2}\,\mathsf{Col}_m\longrightarrow \frac{2\varphi^2}{5}.
\]

另一方面，由定理 \ref{thm:xi-time-part9p-chi2-ranktwo-resonant-floor} 的第一条精确恒等式，
\[
F_{m+2}\,\mathsf{Col}_m-1
=
4^{-m}\sum_{k\ne 0}\bigl|\widehat d_m(k)\bigr|^2.
\]
取极限并用
\[
\frac{2\varphi^2}{5}-1=\frac{1}{5\varphi^3},
\]
便得到
\[
4^{-m}\sum_{k\ne 0}\bigl|\widehat d_m(k)\bigr|^2
\longrightarrow
\frac{1}{5\varphi^3}.
\]
证毕。
\end{proof}

\begin{corollary}[最大纤维超额的精确主项]\label{cor:xi-time-part70d-maxfiber-excess-exact-main-term}
记
\[
\bar d_m:=\frac{2^m}{F_{m+2}},
\qquad
M_m:=\max_{w\in X_m}d_m(w).
\]
则有
\[
M_m-\bar d_m
=
\varphi^{-4}\left(\frac{2}{\varphi}\right)^m
+o\!\left(\left(\frac{2}{\varphi}\right)^m\right).
\]
\end{corollary}
\begin{proof}
由推论 \ref{cor:xi-time-part70ab-extremal-fibers-terminal-bit-localization}，
\[
M_m=\left(\frac{2}{\varphi}\right)^m+O(1).
\]
另一方面，由 Binet 公式，
\[
\bar d_m
=
\frac{2^m}{F_{m+2}}
=
\frac{\sqrt5}{\varphi^2}\left(\frac{2}{\varphi}\right)^m
+o\!\left(\left(\frac{2}{\varphi}\right)^m\right).
\]
由于 \(O(1)=o((2/\varphi)^m)\)，相减即得
\[
M_m-\bar d_m
=
\left(1-\frac{\sqrt5}{\varphi^2}\right)\left(\frac{2}{\varphi}\right)^m
+o\!\left(\left(\frac{2}{\varphi}\right)^m\right).
\]
再用 \(\sqrt5=2\varphi-1\) 与 \(\varphi-1=\varphi^{-1}\) 化简：
\[
1-\frac{\sqrt5}{\varphi^2}
=
1-\frac{2\varphi-1}{\varphi^2}
=
\frac{(\varphi-1)^2}{\varphi^2}
=
\varphi^{-4}.
\]
故结论成立。
\end{proof}

\noindent
由此，末位 \(1\) 终板不再只是最小纤维的容纳层，而在“子覆盖同构”外推下渐近地与整个最小扇区重合；与此同时，充分大窗口上的折叠关系也被彻底排除出自由商与线性陪集两类等势模型。再向临界资源侧推进，dyadic 预算与 escort 温度都在同一比特率 \(\log_2(2/\varphi)\) 处出现不可压平的簇集带，而其下边界由黄金比例与温度参数共同给出显式常数。最后，中心纤维场的二阶矩、碰撞概率、非平凡 Fourier 能量与最大纤维超额都锁定到闭式黄金常数，从而终板刚性、资源相变与频域能量三条先前分立的链被压缩为同一条可审计的 Fibonacci 后续主线。

\endinput

\paragraph{window-\texorpdfstring{$6$}{6} 抽象轨道型的最小循环实现与几何障碍的范畴分离}
在 window-\(6\) 的纤维多值性、自由作用障碍与线性陪集障碍都已确立之后，还可进一步把“抽象轨道大小多重集”与“真实立方体上的几何实现”彻底分开：就纯 \(G\)-集合而言，直方图 \(2^8,3^4,4^9\) 实际上已经能由最小的 \(12\) 阶循环群实现；因此 window-\(6\) 的真正障碍并不位于 Burnside 型轨道计数层，而位于把该轨道型忠实落到 \(\Omega_6\) 上时所必须满足的自由性与线性几何兼容条件。

\leanverified{paper\_xi\_time\_part70da\_window6\_abstract\_orbittype\_minimal\_c12}
\begin{theorem}[window-\texorpdfstring{$6$}{6} 的抽象轨道型可由最小阶循环群 \texorpdfstring{$C_{12}$}{C12} 实现]\label{thm:xi-time-part70da-window6-abstract-orbittype-minimal-c12}
设
\[
\mathfrak m:=2^8\,3^4\,4^9
\]
表示轨道大小多重集，其中大小为 \(2,3,4\) 的轨道个数分别为 \(8,4,9\)。则：
\begin{enumerate}
  \item 若某有限群 \(G\) 的某个有限 \(G\)-集合具有轨道型 \(\mathfrak m\)，则必有
  \[
  12\mid |G|.
  \]
  \item 该下界是锐的。更精确地，循环群 \(C_{12}\) 已足够：若对每个 \(d\mid 12\) 记 \(C_d\le C_{12}\) 为唯一的阶为 \(d\) 的子群，则
  \[
  Y:=8(C_{12}/C_6)\ \sqcup\ 4(C_{12}/C_4)\ \sqcup\ 9(C_{12}/C_3)
  \]
  是一个 \(64\) 点 \(C_{12}\)-集合，其轨道大小多重集恰为 \(\mathfrak m\)。
\end{enumerate}
因此，就抽象轨道算术而言，实现 window-\(6\) 直方图所需的最小群阶恰为
\[
12.
\]
\end{theorem}
\begin{proof}
若有限群 \(G\) 的某个轨道分解中出现大小为 \(2,3,4\) 的轨道，则每个轨道大小都必须整除 \(|G|\)。因此
\[
2\mid |G|,
\qquad
3\mid |G|,
\qquad
4\mid |G|,
\]
从而
\[
\operatorname{lcm}(2,3,4)=12\mid |G|.
\]
这就给出第一条与最小群阶的下界。

现取 \(G=C_{12}\)。由于循环群对每个除数都只有唯一子群，故存在唯一的阶 \(6,4,3\) 子群 \(C_6,C_4,C_3\)。相应陪集空间的基数分别为
\[
|C_{12}/C_6|=2,
\qquad
|C_{12}/C_4|=3,
\qquad
|C_{12}/C_3|=4.
\]
因此上式定义的离散并
\[
Y:=8(C_{12}/C_6)\ \sqcup\ 4(C_{12}/C_4)\ \sqcup\ 9(C_{12}/C_3)
\]
确为一个 \(C_{12}\)-集合，并且其轨道大小多重集恰为
\[
2^8\,3^4\,4^9.
\]
总点数为
\[
8\cdot 2+4\cdot 3+9\cdot 4=16+12+36=64.
\]
故 \(C_{12}\) 已实现该抽象轨道型。结合上界即知最小群阶恰为 \(12\)。证毕。
\end{proof}

\begin{corollary}[window-\texorpdfstring{$6$}{6} 的压缩障碍是几何性的，而非纯轨道计数性的]\label{cor:xi-time-part70da-window6-obstruction-geometric-not-burnside}
对 \(\Omega_6=\{0,1\}^6\) 上的 bin-fold 纤维划分，成立如下分层：
\begin{enumerate}
  \item 在抽象有限 \(G\)-集合范畴中，轨道型
  \[
  2^8\,3^4\,4^9
  \]
  已由定理 \ref{thm:xi-time-part70da-window6-abstract-orbittype-minimal-c12} 中的 \(C_{12}\)-集合实现。
  \item 在真实立方体 \(\Omega_6\) 上，该纤维划分不可能来自任何有限群的自由作用轨道分解。
  \item 同样地，该纤维划分也不可能来自任何有限阿贝尔群上的线性陪集划分。
\end{enumerate}
因此，window-\(6\) 的压缩障碍并不是 Burnside 型轨道大小不可实现性，而是一个严格更强的几何--线性兼容性障碍。
\end{corollary}
\begin{proof}
第一条正是定理 \ref{thm:xi-time-part70da-window6-abstract-orbittype-minimal-c12}。第二、三条则由推论 \ref{cor:foldbin6-degeneracy-multivalued-nonfree} 直接给出：window-\(6\) 的纤维直方图
\[
2:8,\qquad 3:4,\qquad 4:9
\]
不能来自任何自由有限群作用，也不能来自任何有限阿贝尔群上的线性陪集划分。将三条结论并列，便得到所述范畴分层。证毕。
\end{proof}

\endinput

\paragraph{整振幅、仿射压力与 Stokes 归一化的单核闭合}
在稳定层两态一致渐近、bin-fold 退化度的全实矩闭式、相对稳定层均匀基线的 Rényi 散度极限，以及 fold-gauge 常数的 Stirling--Bernoulli 奇层展开均已建立之后，还可把这些此前分散出现的接口压缩为一条复参数解析骨架：同一个整函数振幅同时控制 bin-fold 的复矩主项、输出分布的全轴压力、Rényi 极限曲线的规范延拓以及 \(\log C_m\) 尾部的 Bernoulli 采样。由此，一切归一化无关的有限窗对数观测都退化为该整振幅在负奇整数点上的读出，而剩余自由度严格只保留一个乘法常数。

\begin{theorem}[bin-fold 全复矩谱的整振幅律]\label{thm:xi-time-part70e-binfold-entire-amplitude-law}
\leanverified{paper\_xi\_time\_part70e\_binfold\_entire\_amplitude\_law}
记
\[
\varrho:=\frac{\varphi}{2}\in(0,1),
\qquad
d_m(w):=d_m^{\mathrm{bin}}(w),
\qquad
S_m(z):=\sum_{w\in X_m} d_m(w)^z
\qquad (z\in\CC),
\]
其中 \(d_m(w)^z:=\exp(z\log d_m(w))\) 取实正数上的主支。则对任意紧集 \(K\subset\CC\)，都有一致渐近
\[
S_m(z)
=
\Bigl(\frac{2}{\varphi}\Bigr)^{zm}
\Bigl(F_{m+1}+\varphi^{-z}F_m\Bigr)
\Bigl(1+O_K(\varrho^m)\Bigr)
\qquad (z\in K).
\]
从而
\[
\mathcal A_m(z):=
\varphi^{-m}\Bigl(\frac{\varphi}{2}\Bigr)^{zm}S_m(z)
\]
在 \(\CC\) 的每个紧集上一致收敛到
\[
\mathcal A(z):=\frac{\varphi+\varphi^{-z}}{\sqrt5},
\]
因此 \(\mathcal A\) 为整函数。
\end{theorem}
\begin{proof}
由定理 \ref{thm:fold-bin-two-state-asymptotic}，存在 \(\varepsilon_m(w)\) 满足
\[
d_m(w)=\varphi^{-w_m}\Bigl(\frac{2}{\varphi}\Bigr)^m\bigl(1+\varepsilon_m(w)\bigr),
\qquad
\sup_{w\in X_m}|\varepsilon_m(w)|=O(\varrho^m).
\]
固定紧集 \(K\subset\CC\)。对充分大的 \(m\) 有 \(\sup_{w}|\varepsilon_m(w)|<1/2\)，因而
\[
\log\bigl(1+\varepsilon_m(w)\bigr)=O(\varrho^m)
\]
在 \(w\in X_m\) 上一致成立。由于 \(K\) 有界，遂有
\[
(1+\varepsilon_m(w))^z
=
\exp\!\bigl(z\log(1+\varepsilon_m(w))\bigr)
=
1+O_K(\varrho^m)
\]
在 \(z\in K\) 与 \(w\in X_m\) 上一致成立。于是
\[
d_m(w)^z
=
\Bigl(\frac{2}{\varphi}\Bigr)^{zm}\varphi^{-zw_m}
\Bigl(1+O_K(\varrho^m)\Bigr).
\]
把 \(X_m\) 按末位分拆为
\[
A_0^{(m)}:=\{w\in X_m:w_m=0\},
\qquad
A_1^{(m)}:=\{w\in X_m:w_m=1\},
\]
并用
\[
|A_0^{(m)}|=F_{m+1},
\qquad
|A_1^{(m)}|=F_m,
\]
便得到
\[
S_m(z)
=
\Bigl(\frac{2}{\varphi}\Bigr)^{zm}
\Bigl(F_{m+1}+\varphi^{-z}F_m\Bigr)
\Bigl(1+O_K(\varrho^m)\Bigr).
\]
再乘以 \(\varphi^{-m}(\varphi/2)^{zm}\)，并用 Binet 渐近
\[
\varphi^{-m}F_{m+1}\to \frac{\varphi}{\sqrt5},
\qquad
\varphi^{-m}F_m\to \frac{1}{\sqrt5},
\]
即可得到
\[
\mathcal A_m(z)\to \frac{\varphi+\varphi^{-z}}{\sqrt5}
\]
且收敛在每个紧集上一致。极限函数由显式公式显然为整函数。证毕。
\end{proof}

\begin{corollary}[全实轴仿射压力与单谱塌缩]\label{cor:xi-time-part70e-affine-pressure-monofractal}
\leanverified{paper\_xi\_time\_part70e\_affine\_pressure\_monofractal}
记
\[
\mu_m(w):=\frac{d_m^{\mathrm{bin}}(w)}{2^m},
\qquad
\Xi_m(t):=\sum_{w\in X_m}\mu_m(w)^t
\qquad (t\in\RR).
\]
则极限
\[
\tau_\mu(t):=\lim_{m\to\infty}\frac1m\log \Xi_m(t)
\]
对全部 \(t\in\RR\) 存在，并且满足闭式
\[
\tau_\mu(t)=(1-t)\log\varphi.
\]
因而 \(\tau_\mu\) 在全实轴上仿射，其 Legendre--Fenchel 变换退化为单点谱。等价地，局部信息密度满足一致塌缩
\[
-\frac1m\log \mu_m(w)=\log\varphi+O(m^{-1})
\qquad (w\in X_m),
\]
其中误差在 \(w\) 上一致；并且相对于
\[
\log|X_m|=(m+2)\log\varphi-\frac12\log 5+o(1)
\]
的归一化局部维数亦有
\[
-\frac{\log\mu_m(w)}{\log|X_m|}=1+O(m^{-1})
\qquad (w\in X_m),
\]
其中误差同样在 \(w\) 上一致。
\end{corollary}
\begin{proof}
由定理 \ref{thm:xi-time-part70e-binfold-entire-amplitude-law} 在实轴上的结论，
\[
\Xi_m(t)=2^{-tm}S_m(t)
=
\varphi^{-tm}\Bigl(F_{m+1}+\varphi^{-t}F_m\Bigr)\Bigl(1+O_t(\varrho^m)\Bigr).
\]
又 \(F_{m+1}+\varphi^{-t}F_m=\Theta_t(\varphi^m)\)，故
\[
\frac1m\log\Xi_m(t)\to (1-t)\log\varphi.
\]
这就得到 \(\tau_\mu(t)\) 的显式公式。由于 \(\tau_\mu\) 在全轴上仿射，其 Legendre--Fenchel 变换只能在唯一斜率处非平凡，故相应多重分形谱退化为单点。

另一方面，由定理 \ref{thm:fold-bin-two-state-asymptotic}，
\[
\mu_m(w)=\frac{d_m^{\mathrm{bin}}(w)}{2^m}
=
\varphi^{-m-w_m}\bigl(1+O(\varrho^m)\bigr)
\]
在 \(w\in X_m\) 上一致成立。因此
\[
-\log\mu_m(w)=m\log\varphi+w_m\log\varphi+O(\varrho^m)
\]
一致成立，两边除以 \(m\) 即得
\[
-\frac1m\log\mu_m(w)=\log\varphi+O(m^{-1})
\]
在 \(w\) 上一致。再用 \(|X_m|=F_{m+2}\) 的 Binet 渐近
\[
\log|X_m|=(m+2)\log\varphi-\frac12\log5+o(1),
\]
即可得到最后一式。证毕。
\end{proof}

\begin{theorem}[规范化 Rényi 极限曲线与负奇点 Bernoulli 采样]\label{thm:xi-time-part70e-renyi-curve-negative-odd-sampling}
\leanverified{paper\_xi\_time\_part70e\_renyi\_curve\_negative\_odd\_sampling}
对 \(\alpha>0\)，记 \(D_\alpha^\infty\) 为定理 \ref{thm:fold-bin-renyi-divergence-limit} 中的极限 Rényi 散度，并定义规范化曲线
\[
\mathcal R(\alpha):=
5^{(\alpha-1)/2}\varphi^{2-2\alpha}
\exp\!\bigl((\alpha-1)D_\alpha^\infty\bigr).
\]
则对全部 \(\alpha>0\) 都有严格恒等式
\[
\mathcal R(\alpha)=\frac{\varphi+\varphi^{-\alpha}}{\sqrt5}=\mathcal A(\alpha),
\]
其中 \(\mathcal A\) 为定理 \ref{thm:xi-time-part70e-binfold-entire-amplitude-law} 的整振幅。因而 \(\mathcal R\) 由正实轴数据唯一延拓为整函数，并与 \(\mathcal A\) 完全一致。进一步，对任意固定整数 \(R\ge 1\)，有
\[
\log C_m-\log(2\pi)
=
\frac{\sqrt5}{\varphi^2}
\sum_{r=1}^{R}
\frac{B_{2r}}{r(2r-1)}
\mathcal R\bigl(-(2r-1)\bigr)\,
\varrho^{(2r-1)m}
+
O\!\bigl(\varrho^{(2R+1)m}\bigr).
\]
换言之，fold-gauge 常数的尾部模态恰由规范化 Rényi 极限曲线在负奇整数点的取值经 Bernoulli 变换生成。
\end{theorem}
\begin{proof}
由定理 \ref{thm:fold-bin-renyi-divergence-limit}，当 \(\alpha\neq 1\) 时有
\[
D_\alpha^\infty
=
\frac{(2\alpha-2)\log\varphi+\log(\varphi+\varphi^{-\alpha})-\alpha\log\sqrt5}{\alpha-1}.
\]
故
\[
\exp\!\bigl((\alpha-1)D_\alpha^\infty\bigr)
=
\frac{\varphi^{2\alpha-2}}{5^{\alpha/2}}
\bigl(\varphi+\varphi^{-\alpha}\bigr)
=
\frac{\varphi^{2\alpha-2}}{5^{(\alpha-1)/2}}
\cdot
\frac{\varphi+\varphi^{-\alpha}}{\sqrt5}.
\]
两边再乘以 \(5^{(\alpha-1)/2}\varphi^{2-2\alpha}\)，便得到
\[
\mathcal R(\alpha)=\frac{\varphi+\varphi^{-\alpha}}{\sqrt5}.
\]
而当 \(\alpha=1\) 时，左端按定义等于 \(1\)，右端则满足
\[
\frac{\varphi+\varphi^{-1}}{\sqrt5}=1,
\]
故同样成立。于是 \(\mathcal R\) 在正实轴上与显式函数 \((\varphi+\varphi^{-z})/\sqrt5\) 一致，因此其唯一整延拓正是定理 \ref{thm:xi-time-part70e-binfold-entire-amplitude-law} 中的 \(\mathcal A\)。

另一方面，由定理 \ref{thm:fold-bin-gauge-constant-stirling-bernoulli-hierarchy}，
\[
\log C_m-\log(2\pi)
=
\sum_{r=1}^{R}
\frac{B_{2r}}{r(2r-1)}
\bigl(\varphi^{-1}+\varphi^{2r-3}\bigr)\varrho^{(2r-1)m}
+O\!\bigl(\varrho^{(2R+1)m}\bigr).
\]
而由刚刚得到的闭式，对每个 \(r\ge 1\) 都有
\[
\frac{\sqrt5}{\varphi^2}\mathcal R\bigl(-(2r-1)\bigr)
=
\frac{\sqrt5}{\varphi^2}\cdot
\frac{\varphi+\varphi^{2r-1}}{\sqrt5}
=
\varphi^{-1}+\varphi^{2r-3}.
\]
代回即可得到所示采样公式。证毕。
\end{proof}

\begin{theorem}[零和有限窗对数观测的热力学全决定性]\label{thm:xi-time-part70e-zero-sum-window-thermo-determinacy}
\leanverified{paper\_xi\_time\_part70e\_zero\_sum\_window\_thermo\_determinacy}
令移位算子 \(T\) 作用于序列 \(f=(f_m)_{m\ge 1}\) 为
\[
(Tf)_m:=f_{m+1}.
\]
取任意多项式
\[
Q(\lambda)=\sum_{j=0}^{J}c_j\lambda^j
\]
满足零和条件
\[
Q(1)=\sum_{j=0}^{J}c_j=0.
\]
则对任意固定整数 \(R\ge 1\)，有
\[
Q(T)\log C_m
=
\frac{\sqrt5}{\varphi^2}
\sum_{r=1}^{R}
\frac{B_{2r}}{r(2r-1)}
\mathcal R\bigl(-(2r-1)\bigr)\,
Q\!\bigl(\varrho^{2r-1}\bigr)\,
\varrho^{(2r-1)m}
+
O\!\bigl(\varrho^{(2R+1)m}\bigr).
\]
因此，一切有限窗、零和的对数观测量，例如
\[
\log\frac{C_{m+1}}{C_m},
\qquad
\log C_{m+2}-2\log C_{m+1}+\log C_m,
\]
在渐近层面都完全由 \(\mathcal R\) 决定；其中规范常数 \(\log(2\pi)\) 严格消失。
\end{theorem}
\begin{proof}
由定理 \ref{thm:xi-time-part70e-renyi-curve-negative-odd-sampling}，
\[
\log C_m
=
\log(2\pi)
+
\frac{\sqrt5}{\varphi^2}
\sum_{r=1}^{R}
\frac{B_{2r}}{r(2r-1)}
\mathcal R\bigl(-(2r-1)\bigr)\,
\varrho^{(2r-1)m}
+
O\!\bigl(\varrho^{(2R+1)m}\bigr).
\]
对上式施加算子 \(Q(T)\)。首先
\[
Q(T)\log(2\pi)=Q(1)\log(2\pi)=0.
\]
其次，对任意 \(r\ge 1\) 都有
\[
T\bigl(\varrho^{(2r-1)m}\bigr)=\varrho^{2r-1}\varrho^{(2r-1)m},
\]
故
\[
Q(T)\bigl(\varrho^{(2r-1)m}\bigr)
=
Q\!\bigl(\varrho^{2r-1}\bigr)\varrho^{(2r-1)m}.
\]
逐项代回即可得到所示公式。证毕。
\end{proof}

\begin{theorem}[热力学 germ 的一维 Stokes 归一化刚性]\label{thm:xi-time-part70e-stokes-normalization-rigidity}
设 \((G_m)_{m\ge 1}\) 为一列正数，并假定存在常数 \(c\in\RR\)，使得对每个固定整数 \(R\ge 1\) 都有
\[
\log G_m
=
c+
\frac{\sqrt5}{\varphi^2}
\sum_{r=1}^{R}
\frac{B_{2r}}{r(2r-1)}
\mathcal R\bigl(-(2r-1)\bigr)\,
\varrho^{(2r-1)m}
+
O\!\bigl(\varrho^{(2R+1)m}\bigr).
\]
则
\[
\log G_m-\log C_m
=
c-\log(2\pi)+O\!\bigl(\varrho^{(2R+1)m}\bigr)
\qquad (m\to\infty),
\]
从而
\[
\frac{G_m}{C_m}\longrightarrow e^{\,c-\log(2\pi)}.
\]
亦即：一旦整函数 \(\mathcal R=\mathcal A\) 被固定，整个 fold-gauge germ 在全部归一化无关方向上都已刚性化；剩余自由度严格只保留一个乘法常数，而本文中的 \(2\pi\) 正是该唯一自由度的规范取值。
\end{theorem}
\begin{proof}
由定理 \ref{thm:xi-time-part70e-renyi-curve-negative-odd-sampling}，
\[
\log C_m
=
\log(2\pi)+
\frac{\sqrt5}{\varphi^2}
\sum_{r=1}^{R}
\frac{B_{2r}}{r(2r-1)}
\mathcal R\bigl(-(2r-1)\bigr)\,
\varrho^{(2r-1)m}
+
O\!\bigl(\varrho^{(2R+1)m}\bigr).
\]
将其与 \(\log G_m\) 的假设展开相减，级数主项逐项抵消，即得
\[
\log G_m-\log C_m
=
c-\log(2\pi)+O\!\bigl(\varrho^{(2R+1)m}\bigr).
\]
取 \(R=1\) 并令 \(m\to\infty\)，得到
\[
\log G_m-\log C_m\to c-\log(2\pi).
\]
指数化即得所示比值极限。证毕。
\end{proof}

\begin{conjecture}[连续射影谱的精确闭式与整振幅插值]\label{conj:xi-time-part70e-projective-pressure-entire-amplitude}
沿用附录 \ref{app:pom-projective-pressure-operator} 的射影转移算子 \(\mathcal L_t\)（\(t\ge 1\)）。则其主谱半径满足精确闭式
\[
\rho(\mathcal L_t)=2^t\varphi^{1-t}
\qquad (t\ge 1),
\]
等价地，连续归一化压力
\[
\log\rho(\mathcal L_t)-t\log 2=(1-t)\log\varphi
\]
在整个半轴上恰为一条仿射直线。进一步，存在一种自然归一化，使得主特征数据在 \(t\)-轴上的解析延拓恰给出定理 \ref{thm:xi-time-part70e-binfold-entire-amplitude-law} 的整振幅
\[
\mathcal A(t)=\frac{\varphi+\varphi^{-t}}{\sqrt5}.
\]
若该猜想成立，则定理 \ref{thm:xi-time-part70e-renyi-curve-negative-odd-sampling} 中的尾部采样公式便不再只是 bin-fold moment asymptotics 的结果，而可被提升为连续射影主谱在负奇整数点的 Bernoulli 采样律。
\end{conjecture}
\begin{remark}
该猜想与现有证据链完全同向。首先，定理 \ref{thm:pom-real-q-pressure-analytic-interpolation} 已把整数 moment-kernel 塔嵌入一条实解析压力曲线；定理 \ref{thm:pom-projective-operator-degenerates-to-moment-kernel} 与定理 \ref{thm:xi-projective-moment-tower-anchors-normalized-pressure} 则说明，整数节点上的主谱数据正是离散碰撞谱 \(r_q\)。其次，命题 \ref{prop:fold-bin-power-sum-asymptotic} 已经给出
\[
r_q=2^q\varphi^{1-q}
\qquad (q\in\ZZ_{\ge 2}),
\]
而定理 \ref{thm:xi-time-part70e-binfold-entire-amplitude-law} 与定理 \ref{thm:xi-time-part70e-renyi-curve-negative-odd-sampling} 又表明，同一条离散谱线在复参数方向上伴随一个刚性的整振幅 \(\mathcal A\)，并且该振幅的负奇采样恰重建 \(\log C_m\) 的 Bernoulli 奇层。因而，从离散 moment-kernel、Rényi 极限曲线与 fold-gauge 尾部采样三侧同时观察，全部可见证据都指向同一连续谱公式。
\end{remark}

\noindent
由此得到的图景是：本文框架中的 \(\pi\)-机制并不承载于压力函数的曲率，而承载于一条全轴仿射的热力学骨架之上的整函数振幅及其负奇采样。换言之，指数率本身已经被压平为单一直线；真正保留归一化信息的，是该直线上方的整振幅 \(\mathcal A=\mathcal R\)，以及由 Bernoulli 变换读出的唯一 Stokes 型常数 \(2\pi\)。于是，bin-fold 的复矩谱、Rényi 极限曲线、有限窗消模器与 fold-gauge 常数 germ 不再是并列现象，而被收束为同一条解析主线。

\endinput


\paragraph{零余类有限交截、nerve 的 flag 刚性、加权 clique--Euler 并集公式与共振整函数系数刚性}
在零点余类刚性、约数驱动的谱零点余类并结构、零点削减碰撞上界，以及共振梯整函数的 Laguerre--P\'olya 闭包、实零点结构与截断误差已经建立之后，还可继续沿同一条算术--解析主线抽取出下列七条结果。它们分别刻画零余类族任意有限交截的广义中国剩余判据、零点余类覆盖的 Helly--2 刚性与 nerve 的 flag 闭式、谱零点集合的精确加权 clique--Euler 公式、碰撞上界的约数兼容图闭式、共振整函数 \(\mathcal F_\varphi\) 的对数 Taylor 级数与倒数零矩、其偶系数的显式 Newton 递推与严格对数凹性，以及截断深度对相对误差预算的可审计控制。

\begin{theorem}[零余类族的任意有限交截由两两兼容性完全控制]\label{thm:xi-time-part71-zero-coset-finite-intersection-pairwise-compatibility}
\leanverified{paper\_xi\_time\_part71\_zero\_coset\_finite\_intersection\_pairwise\_compatibility}
设 \(M\) 为偶整数，且
\[
g_1,\dots,g_r\mid \frac{M}{2}.
\]
对每个 \(g\mid M/2\)，定义
\[
S_g(M):=
\left\{
k\in \ZZ/(M\ZZ):
\ k\equiv \frac{M}{2g}\pmod{M/g}
\right\}.
\]
则对任意有限族 \(\{g_1,\dots,g_r\}\)，都有
\[
\bigcap_{j=1}^{r}S_{g_j}(M)\neq\varnothing
\iff
2\gcd(g_i,g_j)\mid (g_j-g_i)
\qquad
(1\le i<j\le r).
\]
并且一旦交截非空，就有精确基数公式
\[
\left|
\bigcap_{j=1}^{r}S_{g_j}(M)
\right|
=
\gcd(g_1,\dots,g_r).
\]
\end{theorem}
\begin{proof}
由定理 \ref{thm:fold-zero-coset-rigidity}，
\[
S_{g_j}(M)
=
\left\{
k\in \ZZ/(M\ZZ):
\ k\equiv \frac{M}{2g_j}\pmod{M/g_j}
\right\}
\qquad (1\le j\le r).
\]
因此
\[
\bigcap_{j=1}^{r}S_{g_j}(M)\neq\varnothing
\]
当且仅当同余组
\[
k\equiv \frac{M}{2g_j}\pmod{M/g_j}
\qquad (j=1,\dots,r)
\]
可解。广义中国剩余定理断言：这一线性同余组可解，当且仅当任意两式相容。另一方面，引理 \ref{lem:fold-zero-coset-intersection} 已给出两式相容的充要条件恰为
\[
2\gcd(g_i,g_j)\mid (g_j-g_i).
\]
于是得到第一条等价关系。

现设该系统可解。由于解集模数为
\[
\operatorname{lcm}\!\left(\frac{M}{g_1},\dots,\frac{M}{g_r}\right)
=
\frac{M}{\gcd(g_1,\dots,g_r)},
\]
故在模 \(M\) 的意义下，解恰形成一个单一同余类的提升，其基数为
\[
\frac{M}{M/\gcd(g_1,\dots,g_r)}
=
\gcd(g_1,\dots,g_r).
\]
证毕。
\end{proof}

\input{subsubsec__xi-time-protocol-conclusions-part71a-zero-coset-nerve-flag}

\begin{theorem}[谱零点集合的精确加权 clique--Euler 公式]\label{thm:xi-time-part71-zero-spectrum-weighted-clique-euler}
\leanverified{paper\_xi\_time\_part71\_zero\_spectrum\_weighted\_clique\_euler}
设
\[
M:=F_{m+2},
\qquad
\mathcal Z_m:=\{k\in \ZZ/(M\ZZ):\ \widehat d_m(k)=0\},
\]
并假设 \(M\) 为偶数。定义
\[
\mathcal D_m:=
\left\{
d:\ d\mid (m+2),\ d<m+2,\ F_d\mid \frac{M}{2}
\right\}.
\]
在顶点集 \(\mathcal D_m\) 上定义兼容图 \(\Gamma_m\)：对不同顶点 \(d,e\in\mathcal D_m\)，规定
\[
d\sim e
\iff
2F_{\gcd(d,e)}\mid (F_e-F_d).
\]
若记 \(\mathrm{Cliq}(\Gamma_m)\) 为 \(\Gamma_m\) 的全部非空 clique，且对每个 clique \(C\) 记
\[
\gcd(C):=\gcd\{d:\ d\in C\},
\]
则谱零点集合满足严格公式
\[
|\mathcal Z_m|
=
\sum_{\varnothing\neq C\in \mathrm{Cliq}(\Gamma_m)}
(-1)^{|C|+1}F_{\gcd(C)}.
\]
\end{theorem}
\begin{proof}
由推论 \ref{cor:fold-zero-coset-union}，当 \(M\) 为偶数时，
\[
\mathcal Z_m
=
\bigcup_{\substack{1\le j\le m\\ g_j\mid M/2}}S_{g_j}(M),
\qquad
g_j=\gcd(F_{j+1},M)=F_{\gcd(j+1,m+2)}.
\]
去重之后，恰可写成
\[
\mathcal Z_m=\bigcup_{d\in\mathcal D_m}S_{F_d}(M).
\]

现取任意非空子集 \(I=\{d_1,\dots,d_r\}\subseteq \mathcal D_m\)。由定理 \ref{thm:xi-time-part71-zero-coset-finite-intersection-pairwise-compatibility}，
\[
\bigcap_{d\in I}S_{F_d}(M)\neq\varnothing
\]
当且仅当对所有 \(1\le i<j\le r\) 都有
\[
2\gcd(F_{d_i},F_{d_j})\mid (F_{d_j}-F_{d_i}).
\]
再由 Fibonacci 最大公因子恒等式
\[
\gcd(F_a,F_b)=F_{\gcd(a,b)},
\]
上式等价于
\[
2F_{\gcd(d_i,d_j)}\mid (F_{d_j}-F_{d_i}),
\]
也即 \(I\) 在 \(\Gamma_m\) 中形成一个 clique。

若交截非空，则再由定理 \ref{thm:xi-time-part71-zero-coset-finite-intersection-pairwise-compatibility}，
\[
\left|
\bigcap_{d\in I}S_{F_d}(M)
\right|
=
\gcd(F_{d_1},\dots,F_{d_r})
=
F_{\gcd(d_1,\dots,d_r)}
=
F_{\gcd(I)}.
\]
因此对并集
\[
\mathcal Z_m=\bigcup_{d\in\mathcal D_m}S_{F_d}(M)
\]
施行包含--排斥，便得到
\[
|\mathcal Z_m|
=
\sum_{\varnothing\neq C\in \mathrm{Cliq}(\Gamma_m)}
(-1)^{|C|+1}F_{\gcd(C)}.
\]
证毕。
\end{proof}

\begin{corollary}[碰撞上界可压缩为约数兼容图的加权 Euler 特征]\label{cor:xi-time-part71-collision-upper-bound-weighted-clique-euler}
沿用定理 \ref{thm:xi-time-part71-zero-spectrum-weighted-clique-euler} 的记号。则碰撞概率满足严格上界
\[
\mathsf{Col}_m
\le
1-\frac{1}{F_{m+2}}
\sum_{\varnothing\neq C\in \mathrm{Cliq}(\Gamma_m)}
(-1)^{|C|+1}F_{\gcd(C)}.
\]
特别地，若 \(\Gamma_m\) 不含三角形，亦即其非空 clique 仅有单点与边，则上界简化为
\[
\mathsf{Col}_m
\le
1-\frac{1}{F_{m+2}}
\left(
\sum_{d\in\mathcal D_m}F_d
-
\sum_{\{d,e\}\in E(\Gamma_m)}F_{\gcd(d,e)}
\right).
\]
\end{corollary}
\begin{proof}
由定理 \ref{thm:fold-collision-zero-reduction}，
\[
\mathsf{Col}_m\le 1-\frac{|\mathcal Z_m|}{F_{m+2}}.
\]
将定理 \ref{thm:xi-time-part71-zero-spectrum-weighted-clique-euler} 的精确公式代入，即得第一式。

若 \(\Gamma_m\) 不含三角形，则全部非空 clique 仅由单点与边组成，包含--排斥求和自动截断为两层：
\[
|\mathcal Z_m|
=
\sum_{d\in\mathcal D_m}F_d
-
\sum_{\{d,e\}\in E(\Gamma_m)}F_{\gcd(d,e)}.
\]
再代回碰撞上界即可。证毕。
\end{proof}

\begin{theorem}[共振整函数的对数 Taylor 级数与倒数零矩闭式]\label{thm:xi-time-part71-resonance-entire-log-taylor-zero-moments}
沿用定理 \ref{thm:fold-resonance-entire-lp} 的记号，记
\[
\mathcal F_\varphi(z):=\prod_{n=2}^{\infty}\cos(\pi z\varphi^{-n}).
\]
则在最近正零点半径
\[
|z|<\frac{\varphi^2}{2}
\]
内，有严格对数展开
\[
\log \mathcal F_\varphi(z)
=
-\sum_{r=1}^{\infty}\frac{\sigma_r}{r}\,z^{2r},
\qquad
\sigma_r:=
(2^{2r}-1)\zeta(2r)\frac{\varphi^{-2r}}{\varphi^{2r}-1}.
\]
等价地，正零点倒数偶矩满足
\[
\sum_{\substack{\rho\in\mathcal Z(\mathcal F_\varphi)\\ \rho>0}}\rho^{-2r}
=
(2^{2r}-1)\zeta(2r)\frac{\varphi^{-2r}}{\varphi^{2r}-1},
\]
而全零点求和式为
\[
\sum_{\rho\in\mathcal Z(\mathcal F_\varphi)}\rho^{-2r}
=
2(2^{2r}-1)\zeta(2r)\frac{\varphi^{-2r}}{\varphi^{2r}-1}.
\]
\end{theorem}
\begin{proof}
由定理 \ref{thm:fold-resonance-entire-lp}，\(\mathcal F_\varphi\) 为整函数，且
\[
\mathcal Z(\mathcal F_\varphi)
=
\left\{
\left(k+\tfrac12\right)\varphi^n:\ k\in\ZZ,\ n\ge 2
\right\}.
\]
当 \(|z|<\varphi^2/2\) 时，对所有 \(n\ge 2\) 都有
\[
|\pi z\varphi^{-n}|<\frac{\pi}{2},
\]
故可使用经典展开
\[
\log\cos x
=
-\sum_{r=1}^{\infty}\frac{(2^{2r}-1)\zeta(2r)}{r\pi^{2r}}\,x^{2r}
\qquad
(|x|<\pi/2).
\]
逐项代入 \(x=\pi z\varphi^{-n}\) 并交换求和，得到
\[
\log \mathcal F_\varphi(z)
=
\sum_{n=2}^{\infty}\log\cos(\pi z\varphi^{-n})
=
-\sum_{r=1}^{\infty}\frac{(2^{2r}-1)\zeta(2r)}{r}
\left(
\sum_{n=2}^{\infty}\varphi^{-2rn}
\right)z^{2r}.
\]
而
\[
\sum_{n=2}^{\infty}\varphi^{-2rn}
=
\frac{\varphi^{-4r}}{1-\varphi^{-2r}}
=
\frac{\varphi^{-2r}}{\varphi^{2r}-1},
\]
遂得
\[
\log \mathcal F_\varphi(z)
=
-\sum_{r=1}^{\infty}\frac{\sigma_r}{r}\,z^{2r}.
\]

另一方面，由定理 \ref{thm:fold-resonance-entire-lp} 的零点描述，
\[
\mathcal Z(\mathcal F_\varphi)\cap (0,\infty)
=
\left\{
\left(k+\tfrac12\right)\varphi^n:\ k\ge 0,\ n\ge 2
\right\}.
\]
故
\[
\sum_{\substack{\rho\in\mathcal Z(\mathcal F_\varphi)\\ \rho>0}}\rho^{-2r}
=
\sum_{n=2}^{\infty}\varphi^{-2rn}\sum_{k\ge 0}\left(k+\tfrac12\right)^{-2r}.
\]
又有
\[
\sum_{k\ge 0}\left(k+\tfrac12\right)^{-2r}
=
2^{2r}\sum_{k\ge 0}(2k+1)^{-2r}
=
(2^{2r}-1)\zeta(2r),
\]
故正零点公式成立。全零点关于原点成对对称，因此再乘以 \(2\) 即得最后一式。证毕。
\end{proof}

\begin{theorem}[共振整函数偶系数的显式 Newton 递推与严格对数凹性]\label{thm:xi-time-part71-resonance-entire-even-coefficient-newton-logconcavity}
\leanverified{paper\_xi\_time\_part71\_resonance\_entire\_even\_coefficient\_newton\_logconcavity}
将 \(\mathcal F_\varphi\) 写成偶幂级数
\[
\mathcal F_\varphi(z)=\sum_{n\ge 0}a_{2n}z^{2n},
\qquad
a_{2n+1}=0,
\]
并定义
\[
b_n:=(-1)^n a_{2n}
\qquad (n\ge 0).
\]
则对全部 \(n\ge 0\) 都有 \(b_n>0\)，并满足精确 Newton 递推
\[
n\,b_n
=
\sum_{r=1}^{n}(-1)^{r-1}\sigma_r\,b_{n-r}
\qquad (n\ge 1),
\]
其中 \(\sigma_r\) 如定理 \ref{thm:xi-time-part71-resonance-entire-log-taylor-zero-moments} 所示。进一步，序列 \((b_n)_{n\ge 0}\) 严格对数凹：
\[
b_n^2>b_{n-1}b_{n+1}
\qquad (n\ge 1).
\]
因此偶系数满足严格交替符号律
\[
a_0>0,\qquad a_2<0,\qquad a_4>0,\qquad \dots,
\]
且相邻比值 \((b_{n+1}/b_n)_{n\ge 0}\) 严格递减。
\end{theorem}
\begin{proof}
先由定理 \ref{thm:xi-time-part71-resonance-entire-log-taylor-zero-moments} 定义形式幂级数
\[
G(t):=\mathcal F_\varphi(\sqrt t)
=
\sum_{n\ge 0}a_{2n}t^n
=
\sum_{n\ge 0}(-1)^n b_n t^n.
\]
则
\[
\log G(t)
=
-\sum_{r=1}^{\infty}\frac{\sigma_r}{r}t^r.
\]
对 \(t\) 求导得
\[
\frac{G'(t)}{G(t)}
=
-\sum_{r=1}^{\infty}\sigma_r t^{r-1}.
\]
把
\[
G(t)=\sum_{n\ge 0}(-1)^n b_n t^n
\]
代入，并比较 \(t^{n-1}\) 的系数，即得
\[
n\,(-1)^n b_n
=
-\sum_{r=1}^{n}\sigma_r(-1)^{n-r}b_{n-r},
\]
等价于
\[
n\,b_n
=
\sum_{r=1}^{n}(-1)^{r-1}\sigma_r\,b_{n-r}.
\]

接着证明 \(b_n>0\) 与严格对数凹性。由定理 \ref{thm:fold-resonance-entire-lp}，
\(\mathcal F_\varphi\) 的正零点可记为
\[
0<\rho_1<\rho_2<\cdots,
\]
且由定理 \ref{thm:xi-time-part71-resonance-entire-log-taylor-zero-moments} 在 \(r=1\) 时的公式，
\[
\sum_{j\ge 1}\rho_j^{-2}<\infty.
\]
因此
\[
G(t)=\prod_{j=1}^{\infty}\left(1-\frac{t}{\rho_j^2}\right)
\]
在 \(t=0\) 邻域内成立，并且
\[
b_n=e_n(\rho_1^{-2},\rho_2^{-2},\dots)>0
\]
是正数列 \((\rho_j^{-2})_{j\ge 1}\) 的第 \(n\) 个初等对称函数。

对每个 \(N>n\)，定义截断多项式
\[
G_N(t):=\prod_{j=1}^{N}\left(1-\frac{t}{\rho_j^2}\right)
=
\sum_{s=0}^{N}(-1)^s b_s^{(N)}t^s.
\]
则 \(b_s^{(N)}\to b_s\)（\(s\) 固定）且 \(b_s^{(N)}>0\)。Newton 不等式给出
\[
\bigl(b_n^{(N)}\bigr)^2
\ge
\frac{n+1}{n}\frac{N-n+1}{N-n}\,
b_{n-1}^{(N)}b_{n+1}^{(N)}.
\]
令 \(N\to\infty\)，得到
\[
b_n^2
\ge
\frac{n+1}{n}\,b_{n-1}b_{n+1}
>
b_{n-1}b_{n+1}.
\]
这就证明了严格对数凹性。由于 \(b_n>0\)，于是
\[
a_{2n}=(-1)^n b_n
\]
严格交替。再将
\[
b_n^2>b_{n-1}b_{n+1}
\]
两边除以 \(b_{n-1}b_n>0\)，即可得到
\[
\frac{b_{n+1}}{b_n}<\frac{b_n}{b_{n-1}},
\]
故相邻比值严格递减。证毕。
\end{proof}

\begin{corollary}[共振梯截断深度的相对误差预算]\label{cor:xi-time-part71-resonance-truncation-relative-error-budget}
\leanverified{paper\_xi\_time\_part71\_resonance\_truncation\_relative\_error\_budget}
沿用命题 \ref{prop:fold-resonance-truncation-error} 的记号。对任意 \(\varepsilon>0\)，当 \(u=0\) 时有
\[
\frac{C_{\varphi,N}(0)}{C_\varphi(0)}=1
\qquad
(\forall N\ge 2).
\]
若 \(u\neq 0\)、\(N\ge 2\) 且
\[
N\ \ge\
-1+\max\!\left\{
\log_\varphi(2\pi|u|),\
\frac12\log_\varphi\!\left(
\frac{\pi^2u^2}{(1-\varphi^{-2})\log(1+\varepsilon)}
\right)
\right\},
\]
则有相对误差界
\[
1\le \frac{C_{\varphi,N}(u)}{C_\varphi(u)}\le 1+\varepsilon.
\]
\end{corollary}
\begin{proof}
当 \(u=0\) 时，
\[
C_{\varphi,N}(0)=C_\varphi(0)=1,
\]
结论平凡。以下设 \(u\neq 0\)。

由命题 \ref{prop:fold-resonance-truncation-error}，只要满足
\[
\pi|u|\varphi^{-(N+1)}\le \frac12,
\]
就有
\[
C_{\varphi,N}(u)\exp\!\left(
-\frac{\pi^2u^2\varphi^{-2(N+1)}}{1-\varphi^{-2}}
\right)
\le
C_\varphi(u)
\le
C_{\varphi,N}(u).
\]
因此
\[
1
\le
\frac{C_{\varphi,N}(u)}{C_\varphi(u)}
\le
\exp\!\left(
\frac{\pi^2u^2\varphi^{-2(N+1)}}{1-\varphi^{-2}}
\right).
\]
若进一步要求
\[
\frac{\pi^2u^2\varphi^{-2(N+1)}}{1-\varphi^{-2}}
\le
\log(1+\varepsilon),
\]
则立得
\[
\frac{C_{\varphi,N}(u)}{C_\varphi(u)}\le 1+\varepsilon.
\]
最后，把上述两条约束分别对 \(N\) 求解，便得到
\[
N\ \ge\
-1+\max\!\left\{
\log_\varphi(2\pi|u|),\
\frac12\log_\varphi\!\left(
\frac{\pi^2u^2}{(1-\varphi^{-2})\log(1+\varepsilon)}
\right)
\right\}.
\]
证毕。
\end{proof}

\noindent
上述七条结果表明，Fibonacci 模零谱与共振整函数 \(\mathcal F_\varphi\) 已经形成一条严格可逆的算术--解析双轨：一方面，谱零点集合 \(\mathcal Z_m\) 不再只是若干候选余类的松散并，而可被零余类覆盖的 flag nerve、约数兼容图的 clique 结构及其加权 Euler 特征精确决定，并进一步把碰撞上界压缩为有限图论计算；另一方面，\(\mathcal F_\varphi\) 的局部解析则不再停留于 Laguerre--P\'olya 性与 Jensen 多项式全实根性这样的抽象实根性，而会在对数 Taylor 系数、倒数零矩、偶系数 Newton 递推与截断深度预算中外显为完全可审计的 \(\zeta(2r)\)--\(\varphi\) 谱矩闭式。于是，零余类几何、有限交 nerve、碰撞上界与共振常数计算深度便被统一收束到同一条由 Fibonacci 约数格与黄金比例谱矩共同支配的闭合链条之中。

\input{subsubsec__xi-time-protocol-conclusions-part72-zero-coset-divisor-dichotomy-lowerbound}

\endinput


\endinput


\endinput


\endinput
